% ############################################################################
% ##  RETIRED 2026-08-03 (USER DECISION). NOT COMPILED BY ANYTHING. DO NOT EDIT.
% ##
% ##  THIS FILE WAS THE SOURCE OF RECORD AND IS NO LONGER.  The split mirror is:
% ##      assumptions.tex  ->  assumptions_main.tex + assumptions_supp.tex
% ##      results.tex      ->  results_s2_identification.tex
% ##                           + results_s3_estimation.tex
% ##                           + results_s3_supp.tex
% ##                           + results_s4_asymptotics.tex
% ##                           + results_supp.tex
% ##  scripts/build-theory/build.sh compiles the SPLIT as of 2026-08-03 and no
% ##  longer copies or requires this file.  paper/qmd/{main,supplement}.qmd have
% ##  always input the split.  proofs.tex is unsplit and remains live.
% ##
% ##  WHY IT WAS RETIRED RATHER THAN RECONCILED.  Two representations of the same
% ##  mathematics drifted, and the build gated this one while the manuscript
% ##  shipped the other -- so the file being validated was not the file being
% ##  published, and a defect repaired in one survived in the other.  Measured
% ##  before the decision, after normalising the \cref/\suppref and
% ##  \eqref/\suppeq translation that legitimately differs:
% ##      ~95 substantive lines lived ONLY here, ~157 ONLY in the split.
% ##  Three concrete instances, all found on 2026-08-02/03:
% ##    * def:onestep HERE asserted two bibliography keys "are absent from this
% ##      project's bibliography" when BOTH ARE PRESENT, and the split had been
% ##      citing them correctly throughout.  A claim about the artifact that the
% ##      artifact contradicted, in the file the build was validating.
% ##    * ass:bierens clause (vi), a 43-line retirement record, is live HERE and
% ##      absent from the split.
% ##    * ass:reprsource carries a 56-line PROOF INSIDE the assumption here,
% ##      which the split had already relocated to assumptions_supp.
% ##
% ##  KEPT, NOT DELETED, because nothing in this repository is committed and a
% ##  deletion would be unrecoverable.  It is also the only place the three items
% ##  above survive in their pre-split form, which is worth having if any of them
% ##  turns out to have been lost rather than moved.
% ##
% ##  IF YOU NEED SOMETHING FROM HERE: move it INTO the split, then check the
% ##  partition still holds -- every \label declared exactly once across the pair.
% ##  build.sh's "multiply defined" check is the live test of that.
% ############################################################################

% ============================================================================
%  CONSOLIDATED ASSUMPTION BLOCK -- project `panel-npiv-discrete`
%
%  Stage 3, dispatch 1 (proof_strategy_memo.md Sec.10, Steps 1-2).
%  Produced by the theorist, 2026-07-28.  ROUND 2: revised against
%  proofs_review.md (56/100).  Change log at the foot of this file.
%
%  CONTRACT.  Tags are STABLE: every later theorem cites by tag
%  (A0, A1, A1b, A1c, B1, A2, A2', H1, A3, A4, N1, A5, A8, A7, C1-C5),
%  never by restating.  The environment numbers may move; the tags may not.
%  FOUR THINGS ARE RETIRED, AND THEIR TAGS AND CLAUSE NUMBERS ARE KEPT RATHER
%  THAN DELETED, for the same reason the contract exists: several files cite
%  them, and a reader who follows one of those citations must land on the reason
%  and not on a gap.  A retired tag or clause is NOT AVAILABLE AS A HYPOTHESIS,
%  and no hypothesis list -- including a BLOCK citation of the parent tag -- may
%  carry one.  THE FOUR RETIREMENTS REST ON THREE DISTINCT GROUNDS AND THE
%  GROUND IS RECORDED AT EACH SITE:
%    A6 (ass:tcc, dispatch 6)  RETIRED BECAUSE FALSE -- prp:tccfails proves its
%       event is null.  Nothing replaces it as a hypothesis; def:a6b is a
%       property no theorem consumes.
%    A3 (ass:rank, dispatch 13)  RETIRED BECAUSE PROVED -- lem:momzero(iv)
%       derives it from A1 and A1b(iv) plus the finiteness tags already on every
%       list that carried it.  Striking it WEAKENS four hypothesis lists and
%       changes no conclusion.
%    A2'(ii) (ass:sourcecondp, dispatch 14)  RETIRED BECAUSE PROVED, on the A3
%       ground -- it follows from A2'(i) in full with B1(i), (ii), (iv.a) via
%       cor:elldiff and (P1).  The clause number is kept as a struck placeholder.
%    B1(vi) (ass:bierens, dispatch 14)  RETIRED BECAUSE IT IS ABOUT A DIFFERENT
%       ESTIMATOR -- the THIRD ground, and it is neither of the other two.  The
%       growing finite family T_n occurs nowhere outside its own sentence, while
%       def:hatT integrates against Pi_t over ALL of T.  What its heading claimed
%       -- identification of varphi_dagger -- is A1b(iii) here, assumed outright.
%  REGISTER RULE, ADDED IN DISPATCH 18 BECAUSE B1(vi)'s RETIREMENT WAS HALF DONE
%  FOR FOUR DISPATCHES: a retirement is complete only when the clause is struck,
%  when no hypothesis list carries it -- INCLUDING VIA A BLOCK CITATION OF THE
%  PARENT TAG -- and when nothing anywhere describes it in the present tense as
%  a live condition.  B1(vi) failed the second test at prp:phirate and the third
%  at three compiled sites; both are closed in that dispatch.
%  THIS HEADER READ "THE TWO RETIREMENTS" THROUGH DISPATCH 17, WHICH HAD BEEN
%  WRONG SINCE DISPATCH 14 ADDED THE OTHER TWO.
%  DISPATCH 7 also adds A5(vi') and moves def:intker and def:hatT into Part V,
%  because A7 is a statement ABOUT T-hat and a maintained assumption must not
%  name an undefined map.
%
%  HOW RESULTS CITE THESE, AND THE ONE THING YOU MUST NOT DO.
%  Each assumption's heading carries its tag, e.g. "Assumption (A2: Source
%  condition for theta)".  Theorems and lemmas cite the TAG in running text
%  ("Suppose B1(i)-(iv) hold").  The tag is the ONLY identifier.
%
%  *** NEVER WRITE \ref{ass:...} OR \cref{ass:...}. ***
%  The environment is declared \newtheorem*{assumption}{Assumption} -- with the
%  STAR, hence UNNUMBERED -- at preamble.tex:283 and in both rendered preambles.
%  A \ref to an unnumbered environment has no counter to read and silently
%  prints whatever counter last incremented: NOTHING AT ALL in the build.sh
%  corpus ("with s as in ."), and the SECTION NUMBER in the article
%  ("with s as in 2").  It is not an undefined reference -- the label exists --
%  so build.sh, render.sh and check_rendering_sync.py ALL report clean.  Only
%  reading the PDF finds it.
%
%  THROUGH 2026-08-01 THIS HEADER SAID THE OPPOSITE: that the environment was
%  \newtheorem{assumption}{Assumption}[section] and that "\cref to the
%  environment label returns the printed NUMBER".  Both halves were false, and
%  they are why 82 such calls accumulated across four files -- an author reading
%  this header would reasonably believe \ref worked.  build.sh now greps for the
%  pattern; see its "assumption refs" finding.
%
%  Requires the macros \IR, \un, \upmodels, \vdagger from preamble.tex.
%
%  Symbols follow quality_reports/theory/panel-npiv-discrete/notation_glossary.md.
%  In particular:  gamma  (not mu) is the Bierens signed measure;
%                  mu     is the source-condition index;
%                  p_K    is the kernel order;  p = dim(Z);
%                  d_Z    is the EFFECTIVE dimension of Delta Z, d_Z <= p-1;
%                  varsigma, s_0 are RESERVED LOCAL BOUND VARIABLES with no
%                           meaning across results -- and varsigma is NOT the
%                           Step-4 direction, which is vartheta (round-7 rename);
%                  vartheta in R^p is the coordinate-contraction direction,
%                           BINDING ON STEP 4;
%                  tau_DP is A7's discrepancy constant; tau = (tau_1,tau_2,tau_3)
%                           is Part F's perturbation parameter.
% ============================================================================

\section*{Assumptions}

% ---------------------------------------------------------------------------
% PART I.  SAMPLING, MODEL, IDENTIFICATION
% ---------------------------------------------------------------------------

\begin{assumption}[A0: Sampling]\label{ass:sampling}
$\{(Y_{i1},Y_{i2},X_{i1},X_{i2},Z_{i1},Z_{i2},W_i)\}_{i=1}^{n}$ are independent and identically
distributed draws from a distribution $P$ on
$\IR^{2}\times\mathcal X^{2}\times\IR^{2p}\times\{0,1\}$, with
$\mathcal X=[\underline x,\bar x]\subset\IR$ a compact \emph{interval}, $\underline x<\bar x$.
\end{assumption}

% A0 is not in the plan's list; it was implicit throughout and is stated here
% because every cross-fitting, empirical-process and rate argument from Step 6
% onwards uses it.
%
% NOTE FOR THE ASSUMPTION-AUDITOR, CORRECTED TWICE.  Through Step 2, A0 was used
% by NOTHING: every result was a statement about the population distribution P
% and held for a single observation.  THAT IS NO LONGER TRUE.  What is consumed,
% and by which clause:
%   lem:levels    COMPACTNESS AND CONNECTEDNESS OF X -- it integrates from the
%                 LEFT ENDPOINT of X, which presupposes an interval;
%   lem:balance   the i.i.d. clause, through the sample-splitting convention
%                 stated in its hypotheses (dispatch-3 round 7);
%   A8            "quasi-uniform mesh over X" is undefined on a disconnected X;
%   H1            "the Sobolev scale of order s" likewise;
%   lem:balance Step 4  Schumaker's direct theorem is a statement about splines
%                 on an interval.
%
% ROUND-7 CORRECTION 1: "COMPACT" WAS NOT ENOUGH, AND THE GAP WAS REAL.  A0 said
% only "X compact".  Compactness admits a Cantor set, a convergent sequence with
% its limit, or a countable union of intervals, and C1(ii) plus int f = 1 forces
% only finite positive Lebesgue measure -- WHICH A FAT CANTOR SET HAS.  On such
% an X all four consumers above are undefined or false.  A0 now says COMPACT
% INTERVAL.  The strengthening is free: the application has a scalar continuous
% regressor.
%
% ROUND-7 CORRECTION 2: THE lem:adjfubini APPEAL IS STRUCK, NOT LISTED.  Round 6
% recorded lem:adjfubini as a second A0 consumer, for sigma-finiteness of P.
% THAT IS OVER-ATTRIBUTION: P is a probability measure by def:model, hence
% sigma-finite with no appeal to A0 at all.  The sentence at the foot of that
% proof is corrected rather than the hypothesis list extended -- the cheaper of
% the two repairs, and the correct one.

\begin{assumption}[A1: Conditional exogeneity]\label{ass:condind}
$\Delta U \upmodels W \mid \Delta Z$ and $E[\Delta U \mid \Delta Z] = 0$.
\end{assumption}

\begin{assumption}[A1b: Identification]\label{ass:ident}~
\begin{enumerate}
\item[(i)] \emph{Normalisation.} $\vdagger\in\mathcal V$, that is, $\int_{\mathcal X}\vdagger\,d\pi=0$.
\emph{(The corresponding statement for a generic $\varphi\in\mathcal V$ is \cref{def:spaces} and not an
assumption; only membership of the truth is content.)}
\emph{Mode of consumption, stated once and binding on every citation of this clause: A1b(i) is consumed
for \textbf{well-posedness} and never for a bound. What it buys is that the truth --- and hence the path
$\vdagger+\tau_1h$ for $h\in\mathcal V$ --- lies in the parameter space. A result whose conclusion
concerns a generic $h\in\mathcal V$, and whose proof needs only $\int h\,d\pi=0$, consumes
\cref{def:spaces} and must cite \cref{def:spaces}; citing A1b(i) there charges the result with a
hypothesis it does not use, because the two statements are typographically identical and semantically
disjoint. The genuine consumers are \cref{lem:secondord} \textup{(Step 1)}, \cref{prp:orthog}\textup{(i)},
\cref{cor:udep} \textup{(clause (a), at $\vdagger$)}, \cref{lem:biasexp} and, transitively,
\cref{lem:solset} and \cref{thm:rootn}.
\emph{\Cref{lem:biasexp}'s use is the archetype}: A1b(i) is what puts
$h_k=\hat\varphi^{(-k)}-\vdagger$ in $\mathcal V$, hence what makes the balancing pairing
$\langle\nu_{g_k}/\pi,h_k\rangle_{L^{2}(\pi)}$ of \eqref{eq:biasdecomp} the object
\cref{lem:nusigma}\textup{(iii)} computes and the object A2 asserts a range condition about --- a
\emph{well-posedness} use and not a bound. \Cref{thm:rootn} carries
A1b(i) in group \textup{(T-1)} for the same use, inherited.
\emph{\Cref{lem:prelimtheta} Step~4, \cref{lem:epterm} Step~3 and \cref{lem:sandwich} Steps~3 and~6$'$
are \textbf{not} consumers}, although each applies \cref{lem:nusigma}\textup{(i)} at $h_k$:
that clause is stated for every $h\in L^{2}(\mathcal X,\pi)$, which is all its proof establishes.}
% ARCHAEOLOGY.  lem:biasexp was missing from this binding list through dispatch 13,
% which is the dangerous direction: the list was wrong and the citation was right,
% so a reader following the list would have struck a correct citation.  Both the
% assumption-auditor and the regularity-checker reached that conclusion
% independently, the latter classifying it as "a citation gap and not a mathematical
% one".  thm:rootn -- the paper's only theorem -- was missing from the same list
% through dispatch 17, the same defect one level up; it survived the dispatch-14
% repair because that repair ran site-by-site rather than list-by-list.  No other
% clause in this block carries a list declared binding; A1d's "sole consumer
% cor:a1cprim" and A2c's "exactly one consumer lem:reprsmooth" are the only
% comparable claims and both were re-checked and are correct.  The three
% non-consumers would each have needed A1b(i) under lem:nusigma(i)'s round-13
% statement "for every h in V".
\item[(ii)] \emph{Measure-separability.} Every function of $X_1$ that is almost surely equal to a
function of $X_2$ is almost surely equal to a constant.
\item[(iii)] \emph{Global identification of $\vdagger$.} $\mathcal V_{CI}=\{\vdagger\}$, where
$\mathcal V_{CI}:=\{\varphi\in\mathcal V:\ \Delta Y-\Delta\varphi\upmodels W\mid\Delta Z\}$.
\item[(iv)] \emph{Identification of $\theta_\dagger$.} $E[\Delta Z\,\Delta Z']$ is nonsingular, and
$E\|\Delta Z\|^{2}<\infty$, $E[\Delta U^{2}]<\infty$, $E\|\Delta Z\,\Delta U\|^{2}<\infty$.
\end{enumerate}
\end{assumption}

% (iii) FILLS THE PLACEHOLDER at draft `:111', which reads `edddee'.  The
% intended content is unrecoverable; (iii) is the theorist's proposal, flagged
% for user confirmation in theory_memo.md Sec.7, item N-3.  Rationale: the
% sentence introducing the block promises "restrictions for identification";
% (i) and (ii) are normalisations removing the additive constant and the kernel
% of Delta; what is missing is injectivity, and injectivity of the map
% varphi |-> law of (Delta Y - Delta varphi, W, Delta Z) is exactly (iii).
%
% AN OPEN GAP, RECORDED RATHER THAN CLOSED.  (iii) IS the conclusion
% "varphi_dagger is identified", assumed.  NO PRIMITIVE ROUTE TO IT IS OFFERED
% HERE, and A1c below is not one: A1c is the LINEARISED condition, it neither
% implies nor is implied by (iii), and pointing at it as the "NPIV analogue"
% would gesture at a bridge this block has just said does not exist.  A methods
% referee will ask what observable feature of P delivers (iii); supplying a
% primitive sufficient condition is open work, recorded in theory_memo.md
% Sec.5.  It is NOT to be closed by weakening the statement.
%
% DISPATCH 2: THE FEVE--FLORENS ROUTE WAS TRIED AND IT DOES NOT REACH (iii).
% Feve & Florens (2014, Prop. 2.2) identify the structural function in a
% first-differenced panel with endogenous regressors from two primitives: strong
% identification of (Z_1,Z_2) by the instrument, and measurable separability of
% Z_1 and Z_2.  Transplanted here those become A1d and A1b(ii), and the
% transplant WORKS -- but it lands on A1c (see cor:a1cprim), not on (iii).  The
% obstruction is exact and is worth stating in the paper rather than hiding.
% In their model the restriction is a CONDITIONAL MEAN and the operator is
% LINEAR in varphi, so ker K = {0} IS identification.  Here the restriction is
% CONDITIONAL INDEPENDENCE: membership of V_CI is the single family
%     E[Wtilde varsigma(Delta Y - Delta varphi) | Delta Z] = 0,   varsigma bounded,
% indexed by test functions of the OBSERVED residual at the trial value, whereas
% running the linearised argument globally would need the larger family
%     E[Wtilde varsigma(Delta Y - Delta varphi) varkappa(Delta U) | Delta Z] = 0,
% indexed also by test functions of the TRUE residual -- which is not observed
% at a trial varphi and about which the restriction says nothing.  The two
% families coincide when varphi = varphi_dagger and not otherwise, which is
% precisely the gap between first-order and global identification.
%
% THIS MAKES RECORDING THE GAP DEFENSIBLE RATHER THAN LAZY.  Feve & Florens
% write of their own Proposition 2.2 that "the assumption ... is too strong".
% So the honest primitive is already hard in the LINEAR case, where the operator
% is fixed and the restriction is a conditional mean.  Citing that is the right
% way to present (iii): not "we could not find a primitive" but "the primitive
% is known to be demanding in the linear special case, and the nonlinearity
% introduced by conditional independence separates it from the linearised
% condition that the primitive does deliver".
%
% (iv) IS NEW IN ROUND 2 and is what actually identifies theta_dagger:
%   theta_dagger = (E[Delta Z Delta Z'])^{-1} E[Delta Z (Delta Y - Delta varphi_dagger)].
% Round 1 routed this through A3, which is CONDITIONAL ON A2.  That made
% identification of theta_dagger contingent on a source condition, which is
% false and contradicts the separation recorded in theory_memo.md Sec.3.
% Nothing in this dispatch, and nothing in the identification of theta_dagger,
% depends on A2.
%
% DISPATCH-13 CORRECTION TO THIS COMMENT, AND IT IS A CORRECTION AND NOT A
% RETRACTION.  Through dispatch 12 the sentence above read "...through A3, which
% is nonsingularity of a DIFFERENT matrix (J, involving ell_star) and which is
% itself conditional on A2."  THE PARENTHETICAL IS FALSE.  By lem:momzero(iv),
% under A1 the ell_star summand of J is the ZERO MATRIX, so J = E[Delta Z Delta Z']
% wherever J is defined: A3's matrix and A1b(iv)'s matrix ARE THE SAME MATRIX.
% What made round 2's repair correct is the CONDITIONALITY and nothing else --
% ell_star exists only under A2, so eq:JSig is a well-formed display only under
% A2, and identification must not be routed through a display that may not exist.
% THE REPAIR STANDS; ITS STATED GROUND IS HALVED.  A3 is now retired (ass:rank);
% A1b(iv) is what identifies theta_dagger AND, through lem:momzero(iv), what
% supplies the rank condition the one-step needs.  So (iv) does MORE work than
% round 2 assigned it, not less.

\begin{assumption}[A1c: First-order local identification]\label{ass:completeness}
The linearisation is injective: $\ker\mathcal A_{\mathrm{cond}}=\{0\}$, where
$\mathcal A_{\mathrm{cond}}:\mathcal V\to L^{2}(\Pi)$ is the Fr\'echet derivative of
\cref{prp:frechet} evaluated at $(\vdagger,\theta_\dagger)$.
\end{assumption}

% THE STATEMENT IS THE OPERATOR ONE.  THE PASSAGE TO THE DRAFT'S
% CONDITIONAL-MEAN FORM IS NOW A THEOREM (prp:kerchar), AND THE PARAGRAPH BELOW
% IS THE RECORD OF WHAT THAT THEOREM COST, NOT A LIVE DEBT.
%
% ROUND-7 DESTALING.  Through round 6 this comment opened by asserting IN
% CAPITALS that the two forms "are not equivalent" and that "the passage between
% them is not free", listed five ingredients as unavailable, and only forty lines
% later recorded that prp:kerchar -- in this same block's own results file --
% PROVES the passage.  A reader stopping at the head carried away a debt the file
% closes.  The head is therefore rewritten: what follows is the accounting of the
% passage, and the passage HOLDS under A1, A4(i), B1(i)-(iv), C1(ii) and C2'(i).
% What is NOT free, and remains the reason A1c is stated in operator form, is
% that the passage consumes those five tags -- so the operator form is the one
% every consumer can cite unconditionally.  Write, for a direction h,
%     m_1(u,z) := E[Delta h | Delta U = u, Delta Z = z, W = 1],
%     m(u,z)   := E[Delta h | Delta U = u, Delta Z = z],
% and let f_1(u|z), f(u|z) be the corresponding conditional densities of
% Delta U.  The conditional-mean form says m_1 = m; the operator form says
% (after the binary identity and the Bierens passage) f_1 m_1 = f m weighted
% and integrated.  m_1 = m DOES NOT GIVE f_1 m_1 = f m unless f_1 = f, which
% is A1.  Passing from one to the other consumes, at least: A1; A4(i) so that
% the W = 1 cell is non-null; a lem:bierens Step-5 passage from an
% L^2(Pi)-a.e. statement to a pointwise one; continuity in u; and a DECLARED
% reference measure on (u,z), since "for a.e. (u,z)" is meaningless until one
% is fixed.  NONE OF THOSE IS AVAILABLE AT THE POINT WHERE A1c IS STATED -- which
% is why A1c is stated in operator form.  All of them ARE available by the time
% prp:kerchar is reached, and the corrections to this list are recorded below.
%
% The operator form is also the form every consumer reads: def:spaces, A2, H1
% and lem:adjoint all name ker A_cond or the singular system of A_cond, none
% of them the conditional-mean display.  So the operator form is primary and
% the draft's display is recorded below as an interpretation, not an
% equivalence.
%
% RENAMED, because "completeness" promises more than the condition delivers.
% A1c IS FIRST-ORDER ONLY, AND NO STRENGTHENING OF IT WILL BUY LOCAL
% UNIQUENESS.  By prp:compact, A_cond is COMPACT; an injective compact
% operator on an infinite-dimensional space is not bounded below, so
% inf over ||h|| = 1 of ||A_cond h|| = 0.  Take h = e_j, the j-th right
% singular function.  By prp:frechet(i), a competing solution at
% varphi_dagger + eps e_j would require
%     eps sigma_j = ||A_cond (eps e_j)|| <= C_A eps^2, i.e. eps >= sigma_j / C_A,
% and sigma_j decreases to 0.  The first-order argument therefore excludes
% competitors only at distances below sigma_j / C_A, a threshold that shrinks
% to zero along the singular directions: A1c does not deliver a neighbourhood
% of uniqueness at all.  GLOBAL IDENTIFICATION IS A1b(iii) AND MUST BE ASSUMED
% SEPARATELY; A1c neither implies nor is implied by it.
%
% The draft's `eq-condcomppanel' (`:127--130') is the conditional-mean display
% above, and it needs correcting in its own right: it conditions on (U,W) only,
% which is the unconditional reading abandoned in plan Sec.17, and it should
% condition on Delta Z as well.
% ROUND 2: retagged from "C6" to "A1c"; CLEANUP PASS: restated and renamed.
%
% DISPATCH 2 ADDS THE MISSING BRIDGE, AND IT IS NOW A THEOREM.  The paragraph
% above says that the draft's conditional-mean display and the operator form
% ker A_cond = {0} "are not equivalent, and the passage between them is not
% free", then lists what the passage would consume: A1; A4(i); a lem:bierens
% Step-5 passage from an L^2(Pi)-a.e. statement to a pointwise one; continuity
% in u; and a declared reference measure on (u,z).  prp:kerchar PROVES the
% equivalence, and the accounting is BETTER than that list, not merely equal to
% it.  Two corrections to the list itself:
%
%   (1) THE REFERENCE MEASURE WAS NEVER MISSING.  A first draft of this note
%       claimed the fifth ingredient was supplied by the strengthened B1(iv.a).
%       THAT IS FALSE.  prp:kerchar never uses varpi_u: clauses (b) and (c) are
%       stated P-almost surely and the reference measure on (u,z) is
%       P_{Delta U, Delta Z}, which is available unconditionally.  B1(iv.a)
%       enters only through lem:bierens Step 5, and only its ROUND-2
%       full-support form.  CONSEQUENCE, and it is the stronger statement:
%       prp:kerchar AND cor:a1cprim SURVIVE A COMPLETE ROLLBACK OF THE B1(iv.a)
%       STRENGTHENING.  The A1c bridge and the N-10 closure are independent
%       results that happen to have been proved in the same dispatch.
%
%   (2) A1 AND A4(i) ARE CONSUMED BY ONLY ONE HALF.  prp:kerchar is split:
%       (a) <=> (b) needs neither, and (b) <=> (c) needs both.  cor:a1cprim
%       reads clause (b), so IT NEEDS NEITHER A1 NOR A4(i) -- only B1(i)-(iv),
%       C1, C2'(i), A1b(i)-(ii) and A1d.  Listing A1 and A4(i) on the corollary,
%       as a first draft did, would have charged it for hypotheses it does not
%       use.
%
% What remains open is A1b(iii), which is a different statement.  See
% prp:kerchar, cor:a1cprim and rem:ffcompare in results.tex.

\begin{assumption}[A1d: Instrument completeness given the residual]\label{ass:instrcomp}
For every $\Lambda\in L^{2}(P_{X_1,X_2})$,
\[
E\big[\widetilde W\,\Lambda(X_1,X_2)\,\big|\,\Delta U,\Delta Z\big]=0\ \text{a.s.}
\qquad\Longrightarrow\qquad
\Lambda(X_1,X_2)=0\ \text{a.s.}
\]
\end{assumption}

% A1d IS NOT MAINTAINED.  It is offered as a PRIMITIVE SUFFICIENT CONDITION for
% A1c and is cited by exactly one result, cor:a1cprim.  No theorem in the
% programme assumes it; A1c remains the maintained condition.  It is stated in
% the block rather than inline so that cor:a1cprim can cite it by tag, per house
% style, and so that the assumption-auditor sees a declared single consumer.
%
% WHAT IT IS.  A1d is the exact analogue, for the conditional-independence
% design, of the "strong identification of (Z_1,Z_2) by W" hypothesis of
% Feve & Florens (2014, Prop. 2.2): injectivity of the map
% Lambda |-> E[Wtilde Lambda | Delta U, Delta Z] on functions of the two periods'
% endogenous regressors.  In their linear conditional-mean panel model that
% hypothesis, together with measurable separability of Z_1 and Z_2 -- our
% A1b(ii) -- DELIVERS IDENTIFICATION OUTRIGHT, because their operator is linear
% in varphi.  Here it delivers A1c and no more: our operator equation is
% NONLINEAR in varphi, so injectivity of the linearisation is a first-order
% statement.  The separation is not a defect of the adaptation; it is what
% distinguishes a conditional-independence restriction from a conditional-mean
% one, and it is discussed at cor:a1cprim and rem:ffcompare.
%
% AND IT IS STRONG.  Feve & Florens themselves write that the assumption of
% their Proposition 2.2 "is too strong".  That judgement transfers: A1d is a
% completeness condition on a two-dimensional argument (X_1,X_2) with a
% one-dimensional binary contrast Wtilde doing the separating, so it will be
% harder here than there, not easier.  It is recorded as a route, priced
% honestly, and NOT adopted.

% ---------------------------------------------------------------------------
% PART II.  THE INDEX SET AND THE WEIGHT FAMILY
% ---------------------------------------------------------------------------

\begin{assumption}[B1: Bierens family and range space]\label{ass:bierens}~
\begin{enumerate}
\item[(i)] \emph{Transform.} $\Phi:\IR^{p}\to\IR^{p}$ is Borel, one-to-one, and bounded:
$c_\Phi:=\sup_{z}\|\Phi(z)\|_\infty<\infty$. The maintained choice is $\Phi$ componentwise $\arctan$,
rescaled so that $c_\Phi=1$.
\item[(ii)] \emph{Weight.} $w=\exp$. Write $\bar w(t):=e^{c_\Phi\|t\|_1}$, so that
$|w(t'\Phi(z))|\le\bar w(t)$ for every $z$.
\item[(iii)] \emph{Index set.} $\mathcal T\subset\IR^{p}$ has strictly positive Lebesgue measure.
\item[(iv)] \emph{Range reference measure.} $\Pi=\Pi_u\otimes\Pi_t$, where
\begin{enumerate}
\item[(iv.a)] $\Pi_u$ is a finite measure on $\IR$ that is absolutely continuous with respect to
Lebesgue measure, with a density $\varpi_u$ that is bounded, strictly positive at every point of $\IR$,
and continuously differentiable with $\|\varpi_u'\|_\infty<\infty$; in particular $\Pi_u$ has full
topological support, $\Pi_u(O)>0$ for every nonempty open $O\subseteq\IR$;
\item[(iv.b)] $\Pi_t$ is a finite measure on $\mathcal T$ for which there exists a Borel set
$\mathcal T_0\subseteq\mathcal T$ such that $\mathrm{Leb}|_{\mathcal T_0}$ and $\Pi_t|_{\mathcal T_0}$
are \emph{not mutually singular};
\item[(iv.c)] $c_\Pi^{2}:=\Pi_u(\IR)\int_{\mathcal T}\bar w(t)^{2}\,\Pi_t(dt)<\infty$.
\end{enumerate}
\item[(v)] \emph{Correction measure.} $\gamma$ is a finite signed measure on $\mathcal T$ with
$\int_{\mathcal T}\bar w(t)\,|\gamma|(dt)<\infty$; and, where $\partial_u L$ is required,
$|\partial_u\lambda(u,t)|\le\Lambda_1(t)$ for all $u$, with
$\int_{\mathcal T}\Lambda_1(t)\,\bar w(t)\,|\gamma|(dt)<\infty$.
\item[(vi)] \textbf{\sout{RETIRED.}} \emph{It describes an estimator this block
does not implement.} The clause read: \emph{the family used on the $\varphi$-side is a finite set
$\mathcal T_n\subset\mathcal T$ with $\#\mathcal T_n\to\infty$ and
$\overline{\bigcup_n\mathcal T_n}\supseteq\mathcal T$} \sout{[retired]}. It \emph{is not maintained, is not
available for use, and is a hypothesis of no result in this programme.} \emph{The clause number is kept
rather than reused, on the A3 and \textup{(P-b)} precedent, so that a stale citation lands
somewhere.}
\emph{That second claim was false for four dispatches, and what made it false is recorded here rather than
repaired silently.} A retirement is complete only when the clause is struck, when no hypothesis list carries
it \emph{including through a block citation of the parent tag}, and when nothing describes it in the present
tense as a live condition. B1(vi) failed the second test at \cref{prp:phirate}, whose list read ``B1'' in full
and now reads B1(i)--(v); and it failed the third at three compiled sites --- the Discussion of this block,
\cref{rem:Tcond} in \texttt{results.tex}, and that remark's rendering in
\texttt{paper/sections/theorems/lem-bierens.tex} --- each of which called it an identification restriction in
the present tense. All four are closed. \emph{A clause that is struck in one file and live in three is worse
than a clause nobody struck}, because the strike is what stops the next reader from checking.
\emph{Why it is retired --- and this one is retired for the third of the three reasons, neither
``it is a theorem'' \textup{(A3)} nor ``it is false'' \textup{(A6)} but ``it is about a different
estimator''.} \emph{The \textbf{growing finite family} $\mathcal T_n$ --- a finite subset of $\mathcal T$
with $\#\mathcal T_n\to\infty$ --- is instantiated \textbf{nowhere in this project}, in the compiled text
or in the simulation code.} The
estimator does not use a growing finite grid: \cref{def:hatT} integrates against $\Pi_t$ over \emph{all} of
% ARCHAEOLOGY.  The sentence read "cal T_n occurs nowhere in this project outside the
% sentence above" through dispatch 20 and was FALSE on the register in which it was
% written: the varphi-step simulation carries an identifier Tn for a FIXED quadrature
% node set -- a constant number of Halton nodes for a Pi_t uniform on cal T, whose
% own header says in terms that it is B1(iii)--(iv.b)--(iv.c) and NOT B1(vi) -- so
% the STRING occurred and the OBJECT did not.  The claim is restated at the object
% rather than at the string, as of dispatch 21.
$\mathcal T$, and every operator norm in Parts~F--H is an $L^{2}(\Pi)$ norm over the whole index set. The
clause is a leftover from a Dom\'inguez--Lobato-style design with a growing finite family, which this block
considered and did not adopt.
\emph{What it was for, and where that work is actually done.} Its heading claimed identification of
$\vdagger$. Identification of $\vdagger$ in this block is \textbf{A1b(iii)}, assumed outright, and A1b(iii)'s
own comment records that no primitive route to it is offered. A growing family would be one ingredient of
such a route --- \emph{and the counterexample-hunter's sub-probe shows it would be a genuine ingredient and
not decoration}: at a \emph{fixed finite} grid $\Pi_t=\sum_{k\le K}\delta_{t_k}$ with $\Delta Z$ supported on
$K+1$ atoms, the $K\times(K+1)$ matrix $[e^{t_k\Phi(z_j)}]$ has a null vector, and perturbing along it makes
the Bierens family hold while the conditional restriction fails. \emph{So the clause is correctly labelled an
identification condition; what it is not is a condition of this estimator.} It is recorded here rather than
deleted so that a later dispatch reinstating the growing-grid design finds the statement and the witness
together. \emph{Note also that (vi) was never sufficient for A1b(iii)}: it constrains the index family and
says nothing about $\mathcal V_{CI}$.
\end{enumerate}
\end{assumption}

% ---------------------------------------------------------------------------
% B1(ii) IS RESTRICTED TO w = exp, AND THAT IS A CORRECTION, NOT A CONVENIENCE.
%
% Round 1 stated B1(ii) as "real-analytic and non-polynomial", which is the
% textbook Bierens phrasing.  UNDER THAT CONDITION lem:bierenscore IS FALSE.
% Counterexample, inside the paper's own maintained transform: let V = +-1 with
% probability 1/2 each, let Phi be the rescaled arctan (an ODD function), so
% Phi(-1) = -Phi(1), and set epsilon := V.  Then E|epsilon| < infinity and
% E[epsilon | V] = V != 0, yet for w = cos and EVERY t,
%      E[epsilon cos(t Phi(V))] = (1/2)cos(t Phi(1)) - (1/2)cos(-t Phi(1)) = 0.
% (ROUND-8: the value Phi(1) was written "a" through round 7, which collided with
% the DEGREE OF ILL-POSEDNESS a of plan Sec.2.6 -- and glossary Sec.1.6 declares
% the sans-serif constant slot expressly to keep generic constants away from that
% very letter.  The value is written out; no symbol is introduced.  This also
% removes the juxtaposition cos(ta) of the RESERVED Bierens index t with a.)
% B1(i)-(iii) all hold and the converse fails.  Structurally: Phi ODD against
% w EVEN means the family sees only the EVEN part of the image measure, and
% Step 4 of the proof of lem:bierenscore recovers only the even moments.  The
% failure is not abstract -- it breaks lem:bierens itself: with Delta Z
% two-valued at +-c and w = cos, condition (c) holds whenever the UNCONDITIONAL
% restriction does, while the Delta-Z-conditional restriction can fail through
% opposite-signed violations that cancel.
%
% w = exp is the maintained choice in plan Sec.2.2, it is what Bierens (1990)
% Thm 1 actually characterises, and it is what any implementation uses.  The
% general-w statement is recorded as rem:swintercept in results.tex, with the
% intercept-indexed Stinchcombe--White family that repairs the failure; it is
% not claimed here because it would change the range space and is consumed
% nowhere in the programme.
%
% CONSEQUENCE FOR PLAN Sec.1: the notation table listed "cos s / sin s" as
% admissible Bierens weights.  Both are inadmissible against an odd Phi.
% ---------------------------------------------------------------------------
%
% Reading of the rest of B1.
%  (i) BOUNDEDNESS does THREE jobs.  It makes w = exp integrable against
%      epsilon with NO moment condition whatsoever on Delta Z, since
%      |w(t'Phi(z))| <= bar-w(t) uniformly in z.  It makes the image measure
%      B |-> E[eps 1{Phi(V) in B}] COMPACTLY SUPPORTED, which is what makes its
%      moment problem determinate in the proof of lem:bierenscore.  And it is
%      why the arctan device may not be skipped.
%      ONE-TO-ONE is what makes sigma(Phi(Delta Z)) = sigma(Delta Z), so that a
%      statement about Phi(Delta Z) is a statement about Delta Z.  Borel image
%      is NOT assumed: an injective Borel map between Polish spaces is
%      automatically bi-measurable (Lusin--Souslin), and for the maintained
%      choice Phi is a homeomorphism onto an open cube, so this is elementary.
%
%  (iv) COMPACTNESS OF T IS SUFFICIENT BUT NOT NECESSARY, and (iv) is now
%      stated at the strength the proofs consume and no more.
%      (iv.a) Step 5 of lem:bierens needs the zero set of M(.,t) to be DENSE
%          in R, which follows from full topological support alone; equivalence
%          to Lebesgue measure, assumed in round 1, is far stronger and idle.
%
%          DISPATCH-2 STRENGTHENING, AND IT IS WHAT CLOSES N-10.  (iv.a) now
%          also requires Pi_u to be ABSOLUTELY CONTINUOUS with a bounded,
%          everywhere strictly positive density varpi_u.  This IMPLIES the old
%          full-support clause (a positive density gives every nonempty open set
%          positive measure), so every proof written against the round-2 form of
%          (iv.a) -- lem:bierens Step 5 is the only one -- goes through verbatim.
%
%          WHY IT IS NEEDED.  The adjoint A*_cond produces, in raw form, an
%          object measurable with respect to F_XZW = sigma(X_1,X_2,Delta Z,W):
%          the u-integration against Pi_u sits against the conditional density
%          f_{Delta U | F_XZW}(u), and until Pi_u(du) can be written varpi_u(u)du
%          that pairing is NOT a conditional expectation and has no
%          (Delta U, Delta Z)-representative.  With the density in hand,
%              int f_{Delta U | F}(u) ell(u,t) Pi_u(du)
%                    = E[ varpi_u(Delta U) ell(Delta U, t) | F_XZW ],
%          and the Bierens induction map of lem:adjoint(ii) is a CONSTRUCTION
%          rather than the assumption A2'(0) used to be.  See lem:adjoint and
%          theory_memo.md Sec.7 item N-10.
%
%          WHY THE DERIVATIVE CLAUSE IS THERE, AND IT IS THE N-10 CLOSURE'S OWN
%          DEBT.  Once lem:adjoint(ii) makes ell_star = W ell^Pi_star, the
%          factor varpi_u(u) sits inside ell_star, so
%              partial_u ell_star(u,z)
%                 = varpi_u'(u) int ell^Pi(u,t) w(t'Phi(z)) Pi_t(dt)
%                 + varpi_u(u) int partial_u ell^Pi(u,t) w(t'Phi(z)) Pi_t(dt).
%          A2'(i) bounds ||partial_u ell_star||_inf.  WITHOUT varpi_u in C^1 THAT
%          DERIVATIVE NEED NOT EXIST, and A2'(i) would be a bound on an object
%          the block does not deliver.  Before N-10 was closed the defect was
%          invisible: A2'(0) made ell_star primitive, so its derivative was
%          whatever the assumption said it was.  Constructing ell_star exposed
%          the debt.  The clause also turns the Pi_u SELECTION LEVER from an
%          aspiration into a theorem: with varpi_u in C^1 and ||varpi_u'||_inf
%          finite, the display above HOLDS whenever partial_u ell^Pi exists, so
%          the design's influence on q and on sigma_q -- hence on
%          prp:sharpquad's sharp constant -- is computable rather than merely
%          conjectured.  BINDING ON STEP 4.
%
%          WHAT IT COSTS: NOTHING ABOUT P.  Pi_u is a reference measure on the
%          RANGE side, chosen by the analyst, exactly like Phi and w.  Any
%          logistic, Gaussian or Cauchy Pi_u satisfies all four clauses --
%          bounded, positive, C^1, bounded derivative -- and every
%          implementation already uses such a Pi_u.  It is a restriction on the
%          normalisation, not on the data-generating process, and it is the only
%          place in the block where that distinction does real work.
%
%          WHAT DOES *NOT* DEPEND ON ANY OF THIS.  prp:kerchar and cor:a1cprim
%          use NO clause of (iv.a) beyond the round-2 full-support one, and they
%          use that only through lem:bierens Step 5.  Their statements are P-a.s.
%          and their reference measure on (u,z) is P_{Delta U, Delta Z}, which is
%          available unconditionally.  THEY SURVIVE A ROLLBACK OF THIS ENTIRE
%          STRENGTHENING.  An earlier note claimed the opposite; see the
%          correction under A1c.
%      (iv.b) Step 5 needs: no Pi_t-null set swallows T_0.  That is EXACTLY
%          non-mutual-singularity of Leb|_{T_0} and Pi_t|_{T_0}.  (If N is
%          Pi_t-null and Leb(T_0 \ N) = 0, then Leb|_{T_0} sits on N and
%          Pi_t|_{T_0} on its complement, i.e. they are mutually singular;
%          conversely non-singularity forbids such an N.)  Equivalently, by the
%          Lebesgue decomposition, (iv.b) SAYS THAT Pi_t HAS A NONZERO
%          ABSOLUTELY CONTINUOUS PART.  Non-mutual-singularity already forces
%          Leb(T_0) > 0 and Pi_t(T_0) > 0, since the zero measure is mutually
%          singular with everything.  Round 1 assumed EQUIVALENCE on T_0, which
%          is two-sided; only one direction was ever invoked, and even that is
%          stronger than needed.
%
%          WHAT (iv.b) EXCLUDES -- corrected in the cleanup pass.  It excludes
%          EVERY LEBESGUE-SINGULAR Pi_t, not merely the atomic ones: a Cantor
%          measure is non-atomic, has uncountable support, and is excluded.
%          Round 2's gloss ("it excludes a purely atomic Pi_t, and that
%          exclusion is the point") named the smallest of the excluded classes.
%
%          AND (iv.b) IS SUFFICIENT, NOT NECESSARY.  What the proof really
%          needs is an ANALYTIC-UNIQUENESS condition on supp Pi_t: by
%          lem:bierenscore Step 2 the map t |-> M(u,t) is real-analytic on all
%          of R^p, so it is enough that its zero set cannot contain supp Pi_t
%          up to a Pi_t-null set unless it vanishes identically.  For p = 1 any
%          Pi_t whose support has an accumulation point suffices -- a Cantor
%          measure included, since a Cantor set is perfect -- by the
%          one-variable identity theorem, and lem:bierens survives such a Pi_t.
%          For p >= 2 the condition is that supp Pi_t not lie inside a proper
%          real-analytic variety; positive Lebesgue measure guarantees that,
%          and so do many Lebesgue-null sets, but perfectness alone does not
%          (a hyperplane is perfect and is the zero set of a coordinate).
%          (iv.b) is retained as the clean sufficient version.
%
%          What (iv.b) does exclude in every dimension is a FINITE grid.  THAT
%          WAS THE GROUND on which (vi) was classed an identification assumption
%          and not an approximation device.  (vi) IS RETIRED (dispatch 14); the
%          exclusion (iv.b) performs is unaffected, and the sentence is kept in
%          the corrected register so that a later dispatch reinstating the
%          growing-grid design finds it.
%      (iv.c) is the integrability that makes L^2(Pi) the right range space;
%          for w = exp it reads Pi_u(R) int e^{2 c_Phi ||t||_1} Pi_t(dt)
%          < infinity, and any compact T with a finite Pi_t satisfies it.
%
%  (v)  By lem:density the estimator SOLVES FOR L rather than integrating a
%      weighting measure, so the first half of (v) is nearly moot on the
%      theta-side: no t-integral survives there.  It binds only if one insists
%      on the integral representation of L.  The SECOND half is the explicit
%      dominating function for differentiating L under the t-integral in u --
%      round 1 asserted that the exponential-moment condition "licenses that
%      directly" without naming either the integral or the dominating function.
%
%  (vi) IS RETIRED (dispatch 14) AND IS NOT AVAILABLE AS A HYPOTHESIS.  The
%      reading below is the reading it HAD while it was maintained, kept because
%      the clause statement is kept: it was the Dominguez--Lobato requirement,
%      an IDENTIFICATION assumption and not a regularity condition, since with a
%      FIXED finite grid the operator equation need not identify
%      varphi_dagger.  Precedent: Poirier (2017); lead case Donald, Imbens &
%      Newey (2009).  WHAT REPLACED IT: nothing -- identification of
%      varphi_dagger in this block is A1b(iii), assumed outright, and the
%      counterexample-hunter's finite-grid witness is recorded at the struck
%      clause so that the statement and the witness are found together.

% ---------------------------------------------------------------------------
% PART III.  SOURCE CONDITIONS AND THE SMOOTHNESS LINK
% ---------------------------------------------------------------------------

\begin{assumption}[A2: Source condition for $\theta$]\label{ass:sourcecond}
Each coordinate of $\nu/\pi$ lies in $\mathcal R(\mathcal A^\ast_{\mathrm{cond}})$, the range of the
adjoint of $\mathcal A_{\mathrm{cond}}:\mathcal V\to L^{2}(\Pi)$, where $\nu:=\nu_{\Delta Z}$ and
$\nu_g(x):=E[g\mid X_2=x]f_{X_2}(x)-E[g\mid X_1=x]f_{X_1}(x)$ as in \eqref{eq:nudef}. Equivalently
(Picard):
$\nu/\pi\perp\ker\mathcal A_{\mathrm{cond}}$ in $L^{2}(\pi)$ and
$\sum_{j\ge1}\sigma_j(\mathcal A_{\mathrm{cond}})^{-2}\,\langle\nu/\pi,\,e_j\rangle_{L^{2}(\pi)}^{2}<\infty$.
\end{assumption}

% ---------------------------------------------------------------------------
% THE OBJECT IS nu / pi, NOT nu.  ROUND-2 CORRECTION, AND IT PROPAGATES.
%
% prp:compact computes the kernel of A_cond as k(u,t;x) = nu_{g_{u,t}}(x)/pi(x).
% The adjoint therefore maps into functions of the form nu_{tilde g}/pi, so
% R(A*_cond) consists of nu-objects DIVIDED BY pi.  A membership test settles it
% independently of any range computation.  By prp:compact Step 2 every kernel
% slice lies in V, so R(A*_cond) is contained in V BY CONSTRUCTION OF THE
% ADJOINT, and V is exactly {zeta : int zeta pi dx = 0}.  Now:
%     zeta = nu     requires int nu pi dx = E[Delta Z (pi(X_2) - pi(X_1))] = 0,
%                   which is FALSE for a generic reference density;
%     zeta = nu/pi  requires int nu dx = E[Delta Z (1 - 1)] = 0, TRUE always.
% So nu/pi is in the right space and nu is not.  Equivalently: nu_g/pi lies in V
% automatically, for every g, since int (nu_g/pi) d pi = int nu_g dx = 0.
%
% NOTE the reason.  This does NOT come from R(A*) subset (ker A)^perp together
% with Delta 1 = 0: once def:spaces fixes the domain at V, the constants are not
% in the domain and that orthogonality argument is unavailable.  It comes from
% the construction of the adjoint, which is simpler and is what prp:compact
% Step 2 actually proves.
%
% This is a bookkeeping consequence of taking the adjoint in L^2(pi), which is
% forced by plan Sec.19.1's dual-pairing principle: the rate for h is in
% L^2(pi), so the residual and the range condition belong there too.  The
% memo's Sec.0.1 identity "A*_cond ell = nu_{Wtilde ell}" is an L^2(dx)
% statement; under the L^2(pi) convention it carries a pi^{-1}.  ONE OF THE TWO
% HAS TO GIVE, AND THE L^2(pi) CONVENTION WINS.  lem:adjoint (Step 3) inherits
% this and must state A*_cond ell = nu_{Wtilde ell}/pi.  Recorded because tags
% are declared stable and Step 3 must not re-open it.
%
% DISPATCH 2: DISCHARGED, AND WITH ONE FURTHER CORRECTION.  lem:adjoint(iii)
% states A*_cond ell = nu_{Wtilde (W ell)} / pi, with the induction map W of
% lem:adjoint(ii) between ell's range-space form and its observable-space form.
% The pi^{-1} is as recorded above; what the round-2 note could not yet say is
% that the "Wtilde ell" inside nu_g is NOT the solution of the adjoint equation
% but its image under W.  Writing nu_{Wtilde ell} without the W is a type error
% of the same kind the round-2 note corrected for pi, one level in.
%
% A2 ALSO ACQUIRES A PRIMITIVE READING.  By lem:adjoint(iii), A2 holds if and
% only if nu_{Delta Z + Wtilde (W ell)} == 0 for some ell in L^2(Pi) -- i.e. if
% and only if the ORTHOGONALITY EQUATION the estimator targets is solvable
% WITHIN THE REACHABLE CLASS R(W).  Since R(W) is dense but not known to be all
% of L^2(P_{Delta U, Delta Z}), A2 IMPLIES solvability of that equation and the
% converse is open.  A2 may therefore be strictly stronger than what the sieve
% estimator needs.  RECORDED, NOT CLAIMED (theory_memo.md finding F-20);
% claiming it would require characterising R(W).
% ---------------------------------------------------------------------------

\begin{assumption}[A2$'$: Source condition with usable representer]\label{ass:sourcecondp}
A2 holds, and the selected solution
$\ell^\Pi_\star\in\ell^\Pi_\dagger+\ker\mathcal A^\ast_{\mathrm{cond}}$ of
$\mathcal A^\ast_{\mathrm{cond}}\ell^\Pi=-\nu/\pi$ \emph{admits a version} satisfying
\begin{enumerate}
\item[(i)] \emph{Range-space regularity.}
\begin{enumerate}
\item[(i.a)] $\Lambda^{(0)}_\star:=\displaystyle\int_{\mathcal T}\Big(\sup_{u\in\IR}\big|\ell^\Pi_\star(u,t)\big|\Big)\,\bar w(t)\,\Pi_t(dt)<\infty$;
\item[(i.b)] for $\Pi_t$-almost every $t$ the map $u\mapsto\ell^\Pi_\star(u,t)$ is continuously
differentiable, and
$\Lambda^{(1)}_\star:=\displaystyle\int_{\mathcal T}\Big(\sup_{u\in\IR}\big|\partial_u\ell^\Pi_\star(u,t)\big|\Big)\,\bar w(t)\,\Pi_t(dt)<\infty$;
\end{enumerate}
and $\ell_\star:=\mathcal W\ell^\Pi_\star$ is defined \emph{pointwise} from that version, through
\eqref{eq:Wdef};
\item[(ii)] \textbf{\sout{RETIRED.}} \emph{It is a consequence of \textup{(i)}.}
The clause read
$E[\widetilde W^{2}\|\ell_\star\|^{2}]<\infty$ \sout{[retired]} and \emph{is not maintained, is not
available for use, and is a hypothesis of no result in this programme.} \emph{The clause number is kept
rather than reused, on the A3 precedent and for the same reason \textup{(P-b)} is kept: a stale
citation must land somewhere.}
\emph{Why it is retired --- and it is retired because it is a \textbf{theorem}, exactly as A3 was, not
because it is false and not because it is too strong.} Under (i) in full, together with B1(i), (ii),
(iv.a), \cref{cor:elldiff} gives $\|\ell_\star\|_\infty\le\|\varpi_u\|_\infty\Lambda^{(0)}_\star<\infty$,
and $|\widetilde W|\le1$ is (P1), free from $W\in\{0,1\}$ and $\varrho\in[0,1]$; hence
$E[\widetilde W^{2}\|\ell_\star\|^{2}]\le\|\ell_\star\|_\infty^{2}<\infty$ with no further hypothesis.
\emph{The attribution is (i) \textbf{in full} and not (i.a) alone}: the only route to
$\|\ell_\star\|_\infty<\infty$ in this block is \cref{cor:elldiff}, that corollary is stated at A2b, and
A2b requires (i.a) \emph{and} (i.b). \emph{Every tag in that derivation is already carried by every list
that could have carried (ii)}, so the retirement changes no conclusion anywhere.
\emph{The clause has \textbf{no consumer}}, as \cref{lem:balance}, \cref{lem:solset} and
\cref{lem:momzero}'s closing paragraph each state independently.
\emph{What survives, and it is the half the clause was really carrying.} At Step~8
% ARCHAEOLOGY.  The round-8 change log of this file recorded "A2'(ii) is consumed by
% none of them and is a Step-8 clause".  A4's comment (d) justified retaining it by
% asserting that "the second-order and bias terms consume" it, which was false, and
% a maintained clause with a false retention ground is what an audit is for.
(\texttt{prp:effbound}), where the variance is computed and where tier~T3 relaxes (i), $\|\ell_\star\|_\infty$
is no longer available and $E[\widetilde W^{2}\|\ell_\star\|^{2}]<\infty$ becomes a genuine restriction ---
it is exactly the Khan--Tamer irregularity boundary, since holding the conditional laws given $W$ fixed and
letting $\varrho\equiv\varepsilon$ gives $E[\widetilde W^{2}\|\ell_\star\|^{2}]\asymp\varepsilon^{-1}$.
\emph{That is a \textbf{forward} obligation and not a present consumer}, and it is recorded here rather than
carried as a maintained clause: Step~8 must state the condition it needs, at the tier where it needs it, and
under T3 it will not be this clause because (i) will not be available to make it redundant. \emph{Flagged
for Step~8}, on the model of \cref{rem:obligations};
\item[(iii)] \emph{Anchor.} The debiasing weight attached to $\ell_\star$ is
$L_\star:=L^{\,b_\star}_{\ell_\star}$ of \cref{def:weight} for a measurable
$b_\star:\IR^{p}\to\IR^{p}$ with $E\big\|b_\star(\Delta Z)\big\|^{2}<\infty$.
\end{enumerate}
\emph{Status of clause \textup{(i)}, stated here so that the tag names exactly one object.} A$2'$(i) is a
\emph{maintained assertion about $\ell^\Pi_\star$, and about $\ell^\Pi_\star$ alone}. The same two
conditions read at a generic $\ell^\Pi\in L^{2}(\Pi)$ are not an assumption but a \emph{property} that a
candidate representer may or may not have; that property is tagged \textbf{A2b} below and is what
\cref{def:weight}(a) names \emph{admissible representer}. \emph{A result that needs the property at a
generic representer cites A2b and never A$2'$(i)}, exactly as a result needing the centring
property of a generic weight cites N1 and never a maintained clause. This is the
treatment N1/N1$'$ received in round 7 and it is applied here for the same reason: one tag, one object.
\end{assumption}

\begin{assumption}[A2b: Admissible representer (a \emph{property}, not maintained)]\label{ass:admrep}
A function $\ell^\Pi\in L^{2}(\Pi)$ is an \emph{admissible representer} if it admits a version satisfying
A2$'$(i.a)--(i.b) with $\ell^\Pi$ in place of $\ell^\Pi_\star$ --- that is, with
\[
\Lambda^{(0)}:=\int_{\mathcal T}\Big(\sup_{u\in\IR}\big|\ell^\Pi(u,t)\big|\Big)\bar w(t)\,\Pi_t(dt)<\infty,
\qquad
\Lambda^{(1)}:=\int_{\mathcal T}\Big(\sup_{u\in\IR}\big|\partial_u\ell^\Pi(u,t)\big|\Big)\bar w(t)\,\Pi_t(dt)<\infty
\]
and $u\mapsto\ell^\Pi(u,t)$ continuously differentiable for $\Pi_t$-almost every $t$ --- and if
$\ell:=\mathcal W\ell^\Pi$ is defined \emph{pointwise} from that version through \eqref{eq:Wdef}.
\emph{A2b is NOT MAINTAINED.} It is a hypothesis a result may impose on a representer it
quantifies over. Under A2$'$(i) the selected solution $\ell^\Pi_\star$ satisfies
A2b, with $\Lambda^{(0)}=\Lambda^{(0)}_\star$ and $\Lambda^{(1)}=\Lambda^{(1)}_\star$; that
implication is the \emph{only} relation between the two tags.
\end{assumption}

\begin{assumption}[A2c: Source condition of index $\mu''$ on the representer (a \emph{property}, not maintained)]\label{ass:reprsource}
Fix $\mu''>\tfrac32$. For $g\in L^{2}(\IR,du)$ write $(\mathcal Fg)(\xi):=\int_{\IR}g(u)e^{-\mathrm iu\xi}du$
for the Fourier transform in $u$, and let
\begin{equation}\label{eq:usobolev}
H^{\mu''}(\IR)\;:=\;\Big\{g\in L^{2}(\IR,du):\
\big\|g\big\|^{2}_{H^{\mu''}(\IR)}:=\tfrac{1}{2\pi}\!\int_{\IR}\!\big(1+\xi^{2}\big)^{\mu''}\big|(\mathcal Fg)(\xi)\big|^{2}d\xi<\infty\Big\}
\end{equation}
be the $u$-Sobolev scale of order $\mu''$ on $\IR$, \emph{taken with respect to Lebesgue measure $du$ and
not with respect to $\Pi_u$}. A function $\ell^\Pi\in L^{2}(\Pi)$ \emph{satisfies the source condition of
index $\mu''$} if it admits a jointly measurable representative --- \emph{fix one; everything below is
read at it} --- each of whose $p$ coordinates satisfies: there is a Borel set
$\mathcal T_\ast\subseteq\mathcal T$ with $\Pi_t(\mathcal T\setminus\mathcal T_\ast)=0$ such that the
slice $\ell^\Pi(\cdot,t)$ lies in $H^{\mu''}(\IR)$ for \emph{every} $t\in\mathcal T_\ast$, and
\begin{equation}\label{eq:lambdamu}
\Lambda^{(\mu'')}\;:=\;\int_{\mathcal T_\ast}\big\|\ell^\Pi(\cdot,t)\big\|_{H^{\mu''}(\IR)}\;\bar w(t)\,\Pi_t(dt)\;<\;\infty ,
\end{equation}
the budget being formed coordinatewise as in A2$'$(i).

\emph{Two things in \eqref{eq:lambdamu} are established here rather than assumed, because neither is
automatic and the second cannot be had from the transform.}

\textup{(m.1)} \emph{The exceptional set is a Borel set by construction, and stating it that way is a
choice.} The weaker phrasing ``for $\Pi_t$-almost every $t$ the slice lies in $H^{\mu''}(\IR)$'' delivers
the displayed form at no cost: it means that the set of $t$ failing the slice condition is contained in
some $\mathcal N\in\mathfrak B(\mathcal T)$ with $\Pi_t(\mathcal N)=0$, and
$\mathcal T_\ast:=\mathcal T\setminus\mathcal N$ is then Borel, carries full $\Pi_t$-measure, and
contains \emph{only} points at which the slice condition holds. \emph{Nothing anywhere requires
$\mathcal T_\ast$ to be the largest such set}, which need not be Borel. What the choice buys is that
$\IR\times\mathcal T_\ast$ belongs to $\mathfrak B(\IR)\otimes\mathfrak B(\mathcal T)$ --- the property
\cref{lem:reprsmooth}'s Step 2 needs when it asserts joint measurability there, and the property an
``almost every'' quantifier does not by itself supply.

\textup{(m.2)} \emph{The integrand is measurable, by pairing against one fixed countable family of
functions of $u$.} There are real $\varkappa_1,\varkappa_2,\ldots\in L^{2}(\IR,du)$ with
\begin{equation}\label{eq:hdual}
\big\|g\big\|_{H^{\mu''}(\IR)}\;=\;\sup_{k\ge1}\Big|\int_{\IR}g(u)\,\varkappa_k(u)\,du\Big|
\qquad\text{for every real }g\in H^{\mu''}(\IR).
\end{equation}
\emph{Construction, with no negative-index space and no completion.} Call a function of $\xi$
\emph{Hermitian} when its value at $-\xi$ is the complex conjugate of its value at $\xi$, which is what
transforms of real functions are. The Hermitian elements of $L^{2}(\IR,(2\pi)^{-1}d\xi)$ form a closed
real-linear subspace of a separable space, hence a separable real Hilbert space; take a countable subset
of its closed unit ball, dense there, and replace each member by its truncations to
$\{|\xi|\le R,\ |\cdot|\le R\}$ over rational $R$ --- truncation preserves Hermitian symmetry, because
the truncation set is symmetric, does not increase the norm, and converges in $L^{2}$ --- so the subset
may be taken to consist of \emph{bounded, compactly supported} Hermitian functions and to remain dense in
the ball. Let $\varkappa_k$ be the inverse Fourier transform of $(1+\xi^{2})^{\mu''/2}$ times the $k$-th
member: that product is bounded with compact support, hence in $L^{1}\cap L^{2}$, so $\varkappa_k$ is a
bounded continuous element of $L^{2}(\IR,du)$, and it is real because the product is Hermitian. For real
$g\in H^{\mu''}(\IR)$, Plancherel in the normalisation of \eqref{eq:usobolev} gives
\[
\int_{\IR}g\,\varkappa_k\,du
\;=\;\frac1{2\pi}\int_{\IR}\big[(1+\xi^{2})^{\mu''/2}(\mathcal Fg)(\xi)\big]\,
\overline{\big[(1+\xi^{2})^{-\mu''/2}(\mathcal F\varkappa_k)(\xi)\big]}\,d\xi ,
\]
whose second factor is the $k$-th member of the dense family and whose first is a Hermitian element of
$L^{2}(\IR,(2\pi)^{-1}d\xi)$ of norm exactly $\|g\|_{H^{\mu''}(\IR)}$, by \eqref{eq:usobolev}. The
supremum over $k$ of the modulus of a continuous linear functional over a dense subset of the closed unit
ball is the functional's norm, which is \eqref{eq:hdual}.

At each fixed $k$ the map $t\mapsto\int_{\IR}\ell^\Pi(u,t)\varkappa_k(u)\,du$ is $\Pi_t$-measurable
on $\mathcal T_\ast$: the product $(u,t)\mapsto\ell^\Pi(u,t)\varkappa_k(u)$ is
$\mathfrak B(\IR)\otimes\mathfrak B(\mathcal T)$-measurable, Tonelli on the $\sigma$-finite product
$(\IR,du)\otimes(\mathcal T,\Pi_t)$ makes the $u$-integrals of its positive and negative parts measurable
in $t$, and both are finite at each $t\in\mathcal T_\ast$ because $\ell^\Pi(\cdot,t)$ and $\varkappa_k$
are both in $L^{2}(\IR,du)$. A countable supremum of measurable functions is measurable, so
$t\mapsto\|\ell^\Pi(\cdot,t)\|_{H^{\mu''}(\IR)}$ is $\Pi_t$-measurable on $\mathcal T_\ast$ and
\eqref{eq:lambdamu} is a well-defined element of $[0,\infty]$, of which A2c asserts finiteness.
\emph{This is the same countable-supremum device \cref{lem:reprsmooth} uses at its Step 3 and
\cref{def:dirclass} at \eqref{eq:shiftnormone}, applied one level up} --- and it is used \emph{instead}
of a Fubini statement about $(\xi,t)\mapsto(\mathcal F\ell^\Pi(\cdot,t))(\xi)$, which is not available,
for the reason \cref{lem:reprsmooth}'s Step 2 prints: on $L^{2}(\IR,du)$ the transform is an
$L^{2}$-limit and not an integral. \eqref{eq:hdual} pairs the slice against a family fixed \emph{once},
in the variable $u$ in which the slice is jointly measurable, so no such statement is needed.

% EXTRACTED FROM THE STATEMENT OF A2c ON 2026-08-03.  Seven lines stood here
% saying what A2c IS -- not maintained, a hypothesis lem:reprsmooth imposes on a
% representer it quantifies over, its only role, and the geometry it is stated in.
% Status and motivation are not conditions, so they are now rem:reprsourceread.
% WHAT DELIBERATELY STAYED, because it disambiguates rather than comments: "taken
% with respect to Lebesgue measure du and not with respect to Pi_u", and "fix one;
% everything below is read at it" for the jointly measurable representative.  Both
% fix WHICH object the clause is about, so a reader who drops them reads a
% different condition.
% NO \label MOVED: eq:usobolev and eq:lambdamu both stayed in the statement.
\end{assumption}

\noindent\emph{Notes on A2c.} That it is a property and not a maintained condition, whose sole role is to
imply A2$'$(i) at the starred representer, and the $\bar w$-weighted geometry it is stated in, are \cref{rem:reprsourceread}.

% ---------------------------------------------------------------------------
% WHY A2c EXISTS, WHY IT IS NOT MAINTAINED, AND WHAT IT DOES NOT DO.
%
% THE OBLIGATION IT DISCHARGES.  (O1) of rem:obligations named L-21 as the
% discharge site for A2'(i.a)-(i.b) IN SUP-INSIDE FORM ON ell^Pi_star.  L-21 is
% now lem:reprsmooth, and A2c is the condition it runs on.  A2'(i) REMAINS A
% MAINTAINED ASSUMPTION; what changes is that it is no longer an unverifiable
% primitive -- it is implied by one smoothness condition on one object, stated in
% the language the block already uses for varphi_dagger (H1, a Hilbert scale).
%
% WHY NOT THE ROUTE THE MEMO REGISTERS.  proof_strategy_memo.md's L-21 row
% proposes "a source condition of index mu'' ... giving ell_star in
% R((A A*)^{mu''}) -> C^1".  THAT ROUTE IS THREE STEPS SHORT, and each gap is
% recorded here rather than papered over:
%   (a) WRONG REFERENT.  It targets ell_star, the OBSERVABLE-space object.  The
%       round-5 repair moved A2'(i) onto the RANGE-space object ell^Pi_star
%       (see clause (c) of the long note below), so a condition on ell_star
%       cannot discharge it -- and cor:elldiff, which bridges the two, CONSUMES
%       A2'(i) and so cannot be used to get there.
%   (b) NO SCALE LINK ON THE RANGE SIDE.  R((A_cond A*_cond)^{mu''}) is a
%       subspace of L^2(Pi) defined by the singular system of A_cond.  Membership
%       in it says NOTHING about sup_u or about u-differentiability unless the
%       A_cond A*_cond scale is MATCHED to a u-smoothness scale.  H1 supplies
%       exactly such a matching on the DOMAIN side (Sobolev order s generated by
%       D, matched to sigma_j asymp j^{-a}); there is no analogue on the range
%       side, and inventing one would be a SECOND Hilbert-scale link, assumed and
%       not derived, about how the left singular functions tilde-e_j oscillate in
%       u.  That is a strictly larger assumption than A2c, on an object the
%       analyst does not choose.  ROUTE REJECTED, and the reason is recorded
%       because it is the obvious thing a reader will propose.
%   (c) NO t-INTEGRATION.  (i.a) and (i.b) are integrals over T of a sup over u,
%       against bar-w d Pi_t.  The memo's row states only the embedding and omits
%       the t-integration entirely, which is half of what the clauses assert.
%       eq:lambdamu supplies it, in the SAME geometry.
%
% WHAT A2c COSTS BEYOND A2, STATED PLAINLY.  ell^Pi_star in L^2(Pi) gives, for
% a.e. t, ell^Pi_star(., t) in L^2(Pi_u) = L^2(varpi_u du).  B1(iv.a) makes
% varpi_u bounded and strictly positive POINTWISE, NOT bounded below, so
% L^2(du) is contained in L^2(Pi_u) and NOT conversely: a slice can lie in
% L^2(Pi_u) with fat tails and fail to lie in L^2(du).  A2c therefore adds a
% TAIL-DECAY requirement in |u| on top of smoothness.  That is the same species
% of condition as C4''s residual truncation, and it is stated rather than hidden.
% The asymmetry is the mirror image of the one recorded at A2 for pi: a bound on
% a density in one direction does not give the reverse containment.
%
% WHY A2c IS CHECKABLE, WHICH IS THE WHOLE POINT.  By eq:usobolev, A2c is a
% weighted summability condition on the u-Fourier transform of the slices -- the
% same KIND of object as A2's Picard sum, and reportable in the application in
% the same way (plan Sec.9.2 item 3).  Picard summability in the A_cond system
% cannot supply it (see item 2 of the FOUR THINGS note below); a summability
% condition in the u-frequency variable can, and does.
%
% mu'' IS INDEPENDENT OF mu.  mu = s/(2a) is a smoothness index for
% varphi_dagger on X; mu'' is a smoothness index for ell^Pi_star in u.  Nothing
% links them, and the smoothness-budget table lists them separately.
%
% THE INTEGER CASE NEEDS NO TRANSFORM, and rem:l21route records it: at mu'' = 2,
% H^2(R) is {g, g', g'' in L^2(du)} and eq:embed holds ~~with c_2 = 2~~ WITH THE
% CONSTANT 2 by an elementary argument.  TWO CORRECTIONS TO THE ROUND-9 FORM OF
% THIS NOTE, both from the dispatch-6 review.
%   (1) THE SYMBOL WAS MISUSED.  c_2 is eq:embedconst at mu'' = 2, and its value
%       is SQRT(2), since int (1+xi^2)^{-1} dxi = pi.  The elementary route gives
%       the WEAKER constant 2.  Writing "c_2 = 2" put two different numbers on
%       one symbol.  Both constants are correct in their own route; the sharper
%       one is c_2 = sqrt(2) and it is the one eq:l21budget is stated with.
%   (2) THE H^2 NORM CARRIES A FACTOR 2 ON ||g'||.  Under eq:usobolev the weight
%       at mu'' = 2 is (1 + xi^2)^2 = 1 + 2 xi^2 + xi^4, so
%       ||g||^2_{H^2} = ||g||^2 + 2||g'||^2 + ||g''||^2 -- eq:h2norm.  The SET is
%       {g, g', g'' in L^2(du)} either way, the two quadratic forms being
%       equivalent with ratio in [1, 2]; the NORM is not.  Step 1' of
%       lem:reprsmooth's proof is correct against the corrected norm, so the
%       comparison sqrt(2) < 2 is like for like.
% An applied reader who does not want the Fourier
% definition can read A2c at mu'' = 2 and lose only mu'' in (3/2, 2).
% ---------------------------------------------------------------------------

% ---------------------------------------------------------------------------
% WHY A2b IS A SEPARATE TAG (round-8 fix, auditor Finding 4 -- IT GATED STEP 4).
%
% THE TAG LETTER WAS SWEPT BEFORE IT WAS CLAIMED, which is the round's own
% organising discipline.  MY FIRST CHOICE WAS "R1", following the auditor's
% "N1' treatment" wording.  R1 IS ALREADY OCCUPIED: results.tex:270 carries
% \begin{remark}[R1]\label{rem:R1}, R1 being a RESULT name from the plan (as are
% L-16, R2, O-11).  Claiming it would have reproduced, inside the fix for
% Finding 4, exactly the defect the warden's structural observation describes --
% a resolution that reasons about the symbol it renames and not about the slot it
% claims.  A2b is used instead: letter-suffixed A-family tags are the project's
% established device for a related-but-distinct condition (A1b, A1c, A1d), the
% string "A2b" occurs nowhere else in any file, and it is visually separable from
% A2' in a way that "N2" would not have been from the memo's escalation items
% N-2, N-7, N-10, N-18.
%
% THE DEFECT.  def:weight(a) defined "admissible representer" by re-using the
% TEXT of A2'(i.a)-(i.b) with ell^Pi in place of ell^Pi_star.  The tag A2'(i)
% therefore named TWO objects at once:
%   - a MAINTAINED ASSERTION about the selected solution (consumed by
%     cor:elldiff at ell^Pi_star, lem:solset, lem:secondord, prp:orthog(ii)-(iii),
%     lem:balance);
%   - a PROPERTY PREDICATE over generic representers (consumed by def:weight(a),
%     hence lem:solset(ii), prp:orthog(i), eq:PsiL).
% N1 got two things this did not: an explicit status sentence, and a SEPARATE TAG
% (N1') when applied to a different object.  A2'(i) got neither.
%
% WHY IT GATED STEP 4.  L-21's target is A2'(i.a)-(i.b) IN SUP-INSIDE FORM ON
% ell^Pi_star.  With the same text doing duty as a predicate over generic
% representers, and cor:elldiff's hypothesis line naming ell^Pi_star while FOUR
% sites applied it generically (results.tex def:weight(a) and eq:PsiL;
% proofs.tex lem:solset(ii) and prp:orthog(i)), a Step-4 author could not read
% off WHICH OBJECT L-21 discharges.  Splitting the tag makes that mechanical:
% L-21 discharges A2'(i) at ell^Pi_star and nothing else; A2b is what a result
% ASSUMES about a representer it quantifies over, and is never discharged by
% L-21 because it is never asserted.
%
% CONSEQUENCE FOR cor:elldiff.  The corollary is now STATED at a generic
% admissible representer under A2b, with ell^Pi_star as the SPECIALISATION.  That
% is a strictly stronger result on strictly weaker hypotheses -- the same
% move prp:orthog(i) received in round 7 -- and it removes the four generic
% applications' reliance on a hypothesis line that named ell^Pi_star.
%
% ORDERING NOTE.  A2b is stated HERE, in the assumption block, and NOT by
% forward reference to def:weight(a), because cor:elldiff (results.tex) precedes
% def:weight in the file and may not cite it.  def:weight(a) now DEFINES the
% phrase "admissible representer" TO MEAN A2b, which is a naming convention and
% not a second statement of the condition.
%
% ROUND-9 CORRECTION TO THIS COMMENT BLOCK, FOUND BY SWEEPING AND NOT BY A
% FINDING.  Three lines above read "R1" where they meant A2b -- the tag this very
% block records as REJECTED, on the ground that R1 already names rem:R1.  The
% round-8 rename was applied to every statement and every citation and NOT to the
% prose that explains the rename, so the block declared a slot unclaimable in one
% paragraph and claimed it in the next.  That is glossary Sec.1.5b's defect
% species one level in -- the sweep for the letter being CLAIMED was run over the
% .tex statements and not over the comment that governs them -- and it is
% corrected in place rather than left as a comment "nobody compiles", because a
% comment block in the file of record is what the next author reads.
% ---------------------------------------------------------------------------

% ---------------------------------------------------------------------------
% CLAUSE (iii) IS NEW IN ROUND 7, AND IT CLOSES A CRITICAL: L_star WAS USED
% EVERYWHERE AND DEFINED NOWHERE.
%
% THE DEFECT.  Through round 6 the symbol L_star appeared in eq:Psitau,
% lem:secondord, prp:orthog(ii)-(iii), A5(iv) and the glossary's eps_n row, and
% was DEFINED IN NO FILE.  The only thing resembling a definition was the aside
% "W ell^Pi_star = partial_u L_star" inside a proof step, which makes L_star an
% ANTIDERIVATIVE IN u OF ell_star WITH A FREE CONSTANT OF INTEGRATION b(Delta Z).
% Nothing constrained b.  Under N1 the gap closes in two lines -- N1 pins b --
% BUT prp:orthog DELIBERATELY WITHHOLDS N1, which is the whole content of its
% part (iii).  With N1 withheld and b arbitrary: eq:Psitau's expectation at tau = 0
% need not exist; E[L_star | Delta Z] on the right of prp:orthog(iii)'s
% equivalence need not exist; and the "=>" direction CONSTRUCTS ITS PERTURBATION
% OUT OF THAT NON-EXISTENT OBJECT, r = eps_varrho sign(E[L_{star,j} | Delta Z]).
%
% THE REPAIR, IN TWO PARTS.  def:weight in results.tex DEFINES the weight
% attached to a representer as L^b_ell(u,z) = int_0^u ell(v,z) dv + b(z), an
% anchor b and a representer ell being exactly the data that determine it.
% Clause (iii) then prices the anchor, and NOTHING ELSE.
%
% WHY THIS IS THE WEAKEST FORM THAT CARRIES THE PROOF, AND WHAT MORE WOULD BUY.
% By cor:elldiff, ||ell_star||_inf <= ||varpi_u||_inf Lambda^(0)_star < infinity
% under (i).  Hence, for EVERY u and z,
%       ||L_star(u,z)||  <=  ||b_star(z)||  +  ||ell_star||_inf |u| ,
% so that, with A1b(iv)'s E[Delta U^2] < infinity and lem:nusigma(i)'s
% E[(Delta h)^2] <= 4 bar-c_pi ||h||^2,
%   (a) E ||L_star(Delta U - tau_1 Delta h, Delta Z)|| < infinity for EVERY
%       tau_1 -- no shift restriction is needed on the L_star term, which is why
%       lem:secondord's shift proviso is imposed on the DIRECTION D and not here;
%   (b) E[ ||L_star|| | Delta Z ] < infinity a.s., which is what makes
%       prp:orthog(iii)'s conditional expectation and its sign perturbation
%       well defined;
%   (c) E[ Wtilde^2 ||L_star||^2 ] <= 2 E||b_star||^2
%                                   + 2 ||ell_star||_inf^2 E[Delta U^2] < infinity,
%       which is the VARIANCE of the moment -- see the correction under A4(d).
% Anything stronger -- L_star in L^2(P) as a standing convention, or a sup-norm
% bound on L_star -- would buy NOTHING beyond (a)-(c) and would be false for the
% affine weights of cor:udep, whose L(u,z) = c(z) u is unbounded in u.
%
% AND IT COSTS NOTHING AT EITHER ANCHOR OF INTEREST, WHICH IS THE POINT.
%   canonical anchor  b_star == 0:  E||b_star||^2 = 0.
%   N1 anchor  b_star(z) = -E[ int_0^{Delta U} ell_star(v,z) dv | Delta Z = z ]:
%       ||b_star(z)|| <= ||ell_star||_inf E[ |Delta U| | Delta Z = z ], so by
%       conditional Jensen E||b_star||^2 <= ||ell_star||_inf^2 E[Delta U^2]
%       < infinity by A1b(iv).
% So (iii) BINDS ONLY on an anchor the analyst has chosen badly.  It is stated
% because the anchor is a free design object, not because either canonical choice
% is in doubt.
%
% VERSIONS (m-3).  ell^Pi_star is an equivalence class in L^2(Pi), and neither
% "sup_u |ell^Pi_star(u,t)|" nor "continuously differentiable in u" is invariant
% under Pi-null modification -- while cor:elldiff's conclusions are evaluated at
% the RANDOM point (Delta U, Delta Z).  A2' therefore says "admits a version"
% and fixes ell_star pointwise from that version.  Without this the sup-norm
% clauses are not well posed as stated.
%
% VECTOR NORM IN (ii).  (ii) was written E[Wtilde^2 ell_star^2] for an
% R^p-valued ell_star.  It is now E[Wtilde^2 ||ell_star||^2], which is the
% coordinatewise statement summed and is well typed; nothing downstream changes.
% ---------------------------------------------------------------------------
%
% ---------------------------------------------------------------------------
% ROUND-2's CLAUSE (0) IS DELETED, BECAUSE IT IS NOW A THEOREM.  Round 2 added
% an "evaluability" clause asserting ell_star in closure(G_gamma), so that
% ell_star(Delta U, Delta Z) would be a well-defined random variable.  The
% clause repaired a real type collision -- a solution of A*_cond ell = -nu/pi
% lives on (u,t) with t the BIERENS INDEX, while (i), (ii) and eq:JSig (through
% dispatch 12: "A3") evaluate
% ell_star at (u,z) with z a VALUE of Delta Z -- but it repaired it by
% ASSERTING the referent instead of constructing it.  Escalated as N-10.
%
% lem:adjoint(ii) NOW CONSTRUCTS THE MAP.  Under the strengthened B1(iv.a) the
% Bierens induction operator
%     (W ell^Pi)(u,z) := varpi_u(u) int_T ell^Pi(u,t) w(t'Phi(z)) Pi_t(dt)
% is a BOUNDED linear map L^2(Pi) -> L^2(P_{Delta U, Delta Z}) with dense range,
% and lem:adjoint(iii) shows that it COMMUTES WITH THE ADJOINT EQUATION:
%     A*_cond ell^Pi = nu_{Wtilde (W ell^Pi)} / pi.
% So for EVERY solution ell^Pi of the adjoint equation the observable-space
% object ell_star := W ell^Pi is automatically square integrable and
% ell_star(Delta U, Delta Z) is automatically a well-defined random variable.
% An assumption that has become a theorem must leave the block; leaving it in
% would show up in the usage matrix as an unused-or-redundant clause and would
% overstate what is assumed.  A2'(i) and A2'(ii) KEEP THEIR NUMBERS, so every
% downstream citation of A2'(i) or A2'(ii) is unaffected.
%
% THREE THINGS THE CONSTRUCTION DOES NOT DO, recorded so that Step 4 does not
% inherit them silently.
%   (a) W IS NOT SURJECTIVE.  Its range is dense (lem:adjoint(ii)) but need not
%       be all of L^2(P_{Delta U, Delta Z}), so a variance-minimising selection
%       posed over the full space may be unreachable.  That is prp:effbound's
%       problem (Step 8), not lem:adjoint's.
%   (b) W IS NOT INJECTIVE in general; what IS proved is
%       ker W subset ker A*_cond, which is all the selection argument of L-16
%       needs, since a direction killed by W is killed by the adjoint too.
%   (c) CLAUSE (i) IS NOW STATED ON THE RANGE-SPACE OBJECT, AND THAT IS A
%       ROUND-3 REPAIR, NOT A RESTATEMENT.  Rounds 1-2 wrote (i) as
%       ||ell_star||_inf < infinity AND ||partial_u ell_star||_inf < infinity.
%       AFTER N-10 CLOSED, THAT WAS ILL-POSED: ell_star = W ell^Pi_star is
%       CONSTRUCTED, and nothing in the block made partial_u ell_star EXIST.
%       B1(iv.a)'s C^1 clause bounds only the varpi_u' term of eq:Wderiv; the
%       second term needs partial_u ell^Pi to exist with a dominating function,
%       WHICH B1(v) DOES NOT SUPPLY -- B1(v) dominates partial_u lambda in the
%       L-representation, and since W ell^Pi_star = partial_u L_star, that is one
%       derivative short of what eq:Wderiv consumes.
%
%       THE REPAIR PUTS THE CONDITION WHERE IT BELONGS.  Differentiability is a
%       property of the SELECTED SOLUTION, not of an object the analyst chooses,
%       so it cannot go in B1 at all.  (i.a) and (i.b) are conditions on
%       ell^Pi_star; cor:elldiff then DERIVES
%           ||ell_star||_inf        <= ||varpi_u||_inf Lambda^(0)_star,
%           ||partial_u ell_star||_inf
%              <= ||varpi_u'||_inf Lambda^(0)_star + ||varpi_u||_inf Lambda^(1)_star,
%       together with eq:Wderiv itself.  What used to be assumed is now a
%       conclusion, and (i.b) is exactly the dominating function that licenses
%       differentiating under the Pi_t-integral.
%
%       B1(v) IS NOT AMENDED.  Its second half remains the dominating function
%       for the GENERIC L-representation quoted in rem:Tcond; at the SELECTED
%       ell_star the domination now comes from (i.b).  Adding a second-derivative
%       clause to B1(v) was the alternative and was rejected: it would impose a
%       smoothness condition on lambda, an object the analyst does not choose
%       once lem:adjoint(ii) fixes gamma = Pi_t and partial_u lambda =
%       varpi_u ell^Pi.  THE TWO SYMBOLS ARE RELATED BY ONE u-INTEGRATION AND
%       ARE NOT INTERCHANGEABLE; C13's reservation of lambda stands.
%
%       THE SUPREMUM IS INSIDE THE INTEGRAL, AND THAT IS A ROUND-5 CORRECTION.
%       Round 4 wrote (i.a) and (i.b) as  sup_u int (.) dPi_t < infinity.  THAT
%       DOES NOT DOMINATE.  Leibniz's rule needs int sup_u (.) dPi_t < infinity;
%       the sup-outside form gives L^1(Pi_t)-boundedness in u WITHOUT uniform
%       integrability, and cor:elldiff's proof silently used the stronger one.
%       Two repairs were available:
%         (alpha) move the sup inside -- STRICTLY STRONGER, since
%                 sup_u int <= int sup_u, so the round-4 clauses are implied and
%                 every printed conclusion survives with the same constants;
%         (beta)  keep the sup outside and weaken the conclusion to: ell_star is
%                 LIPSCHITZ in u with constant Lambda^(1)_star, and eq:Wderiv
%                 holds for a.e. u rather than every u.
%       (alpha) IS ADOPTED, AND THE DECISIVE REASON IS NOT THE ONE FIRST GIVEN.
%       Round 5 justified (alpha) by pointing at lem:secondord(i) and cor:udep,
%       which evaluate partial_u ell_star at Delta U, a RANDOM point.  cor:udep
%       is a weak instance: its use of partial_u ell_star is at the FUNCTION
%       level (whether it is identically zero), not at a random argument.
%       THE DECISIVE CONSUMER IS THE C2'(i) REMOVAL.  Under (beta), eq:Wderiv
%       holds only for a.e. u, so using it at Delta U requires Delta U to charge
%       no Lebesgue-null set -- i.e. requires an absolute-continuity condition,
%       which is exactly what C2'(i) supplies and what this round REMOVED from
%       lem:secondord and prp:orthog as idle.  (beta) would put it straight
%       back.  SO (alpha) AND THE C2'(i) REMOVAL STAND OR FALL TOGETHER, and
%       choosing (beta) would have undone a separate correction in the same
%       round.  A2'(i) is a sup-norm condition in any case, so tightening it in
%       the sup direction is in keeping.
%
%       WHAT REMAINS OPEN, AND THE TARGET HAS BEEN RAISED.  (i.a) and (i.b) are
%       still primitives.  L-21 -- a source condition of index mu'' on
%       ell^Pi_star plus a Sobolev embedding -- is the route to them and is
%       ~~unproved~~ **PROVED IN ROUND 9 as lem:reprsmooth, on the condition A2c.**
%       **L-21 MUST NOW DELIVER THE SUP-INSIDE FORM.**  theory_memo.md
%       Sec.D5 item 2 previously described L-21's target as the sup-outside
%       form; that is the form that does NOT dominate, so aiming L-21 at it would
%       leave cor:elldiff unproved after Step 4 had apparently discharged it.
%       Recorded before Step 4 runs, which is when it is cheap.
%       **AND IT IS THE SUP-INSIDE FORM THAT lem:reprsmooth DELIVERS**: A2c is
%       stated in the bar-w L^1(Pi_t) geometry of (i.a)-(i.b) themselves, and
%       lem:reprsmooth(ii) is a SLICEWISE inequality, so the sup never has to be
%       moved across the t-integral at all.  The raised target is met, not
%       approximated.
%       The gain over rounds 1-2 is unchanged: the clauses are now WELL FORMED
%       and are stated on an object whose existence the block guarantees.
%
% ONE CONSEQUENCE FOR STEP 4.  By eq:Wderiv the DESIGN CHOICE Pi_u enters
% partial_u ell_star, hence q, hence BOTH sigma_q AND the operator K with kernel
% q f_{X_1,X_2}, hence the second-order form of lem:secondord(i).  Pi_u is
% therefore a THIRD selection lever alongside the choice of ell^Pi within its
% affine set, and SG-15 is where it is swept.
%
% DO NOT WRITE "(1/2)||sigma_q/pi||_inf IS prp:sharpquad's SHARP CONSTANT."  An
% earlier version of this comment did.  FOUR things are wrong with it and all
% four are recorded at lem:solset(iii):
%   (a) sigma_q is R^p-valued (F-15), so a NORM presupposes a direction
%       vartheta in R^p with ||vartheta|| = 1: the object is
%       ||vartheta' sigma_q / pi||_inf, never ||sigma_q/pi||_inf.  (The direction
%       letter is vartheta and NOT varsigma; varsigma is a reserved LOCAL bound
%       variable and round 6 wrongly made it binding.  Round-7 rename.)
%   (b) the form on L^2(pi) is B = (vartheta' sigma_q/pi). - 2 Ktilde, and B is
%       NOT self-adjoint, since q f_{X_1,X_2} is not a symmetric kernel;
%   (c) a quadratic form sees only the symmetric part, so the sharp constant is
%       (1/2)||B_sym|| with B_sym = (B + B*)/2, and (1/2)||B|| is only an UPPER
%       bound;
%   (d) ||vartheta' sigma_q/pi||_inf is a LOWER bound on ||B_sym||, via Weyl's
%       theorem and the essential spectrum -- useful, and it is what makes the
%       form non-compact, but it is neither the constant nor an upper bound.
%
% AND THE SPACE IS L^2(pi), NOT V -- ROUND-7 CORRECTION, AND IT GATES STEP 4.
% def:spaces says in italics that every operator in this paper has domain V.  B
% IS THE ONE DECLARED EXCEPTION, and it has to be, because V = {h : int h pi = 0}
% is NOT INVARIANT under multiplication by vartheta' sigma_q / pi.  Three
% consequences that Step 4 must not inherit silently:
%   (i)  ||vartheta' sigma_q / pi||_inf < infinity is not automatic and is now
%        proved, not assumed: |sigma_q(x)| <= ||partial_u ell_star||_inf
%        (f_{X_1}(x) + f_{X_2}(x)) by (P1), so
%        ||vartheta' sigma_q/pi||_inf <= 2 bar-c_pi ||partial_u ell_star||_inf,
%        finite by C1(ii) and cor:elldiff.  Without it the essential-spectrum
%        statement would be about an UNBOUNDED operator.
%  (ii)  Weyl and the norm ||B_sym|| are statements on ALL of L^2(pi), whereas
%        lem:secondord(i)'s quadratic form is evaluated only at h in V, so the
%        form's sharp constant is (1/2)|| P_V B_sym P_V |_V ||.  THE TWO ARE NOT
%        THE SAME OPERATOR.
% (iii)  The lower bound transfers anyway, and lem:solset(iii)'s proof now says
%        why: V has codimension one, so B_sym - P_V B_sym P_V has rank at most
%        two, a finite-rank perturbation changes no essential spectrum, and
%        sigma_ess(P_V B_sym P_V |_V) = sigma_ess(B_sym).
% BINDING ON STEP 4.
%
% NOTE the status of eq:Wderiv: it is a CONDITIONAL identity, and (i.b) is its
% hypothesis, not a theorem.  Calling the lever "delivered" overstates it.

% FOUR THINGS ABOUT A2'.
%
% 1. (ii) IS IMPLIED BY (i).  Since |Wtilde| <= 1 -- which is free from
%    W in {0,1} and varrho in [0,1], and needs NO overlap condition -- (i)
%    gives E[Wtilde^2 ||ell_star||^2] <= || ||ell_star|| ||_inf^2 < infinity.
%    The memo lists both without noting the implication.  Both are retained
%    deliberately: tier T3 (the trimmed estimator, memo Sec.6.6) RELAXES (i)
%    and keeps (ii), so (ii) must exist as a separate, citable clause.
%    Recorded in theory_memo.md Sec.6, finding F-1.
%
% 2. PICARD CANNOT SUPPLY (i).  The Picard criterion in A2 delivers only
%    ell^Pi_star in L^2(Pi) -- an L^2 membership statement.  It says nothing
%    about sup-norms, and (i) is precisely a sup-norm statement.  The route to
%    (i) (memo L-21, ~~to be proved at Step 4~~ **PROVED at Step 4, round 9:
%    lem:reprsmooth on the condition A2c**) is a SOURCE CONDITION OF INDEX mu''
%    ON ell^Pi_star, followed by a Sobolev embedding valid once mu'' exceeds the
%    embedding threshold for the relevant scale -- **which is mu'' > 3/2, the
%    threshold for H^{mu''}(R) -> C^1_b(R) in one dimension**.  mu'' is genuinely
%    independent of mu: mu is a smoothness index for varphi_dagger, mu'' one for
%    ell^Pi_star.  ROUND-7 REFERENT CORRECTION: this note, the smoothness-budget
%    table and the glossary all said "mu'' for ell_star" after the round-5 repair
%    had moved (i) onto the RANGE-SPACE object ell^Pi_star.  L-21 must deliver
%    (i.a)-(i.b), which are conditions on ell^Pi_star, so mu'' is an index for
%    ell^Pi_star.  **ROUND-9: the glossary row for mu'' was the last uncorrected
%    instance of that defect and is now fixed; the round-7 correction reached this
%    note and the budget table and not the row.  Found by sweeping the referent,
%    not by a finding.**  Without L-21, (i) is an unverifiable primitive and the
%    entire second-order analysis rests on it -- which is why it is stated as an
%    assumption here and discharged, not asserted.
%    WHAT THE DISCHARGE DOES AND DOES NOT CHANGE.  A2'(i) is STILL MAINTAINED and
%    still carried as a hypothesis by every consumer.  What lem:reprsmooth removes
%    is its status as an UNVERIFIABLE primitive: it is now implied by one
%    smoothness-plus-tail condition on one object, in a language the block already
%    uses (H1's Hilbert scale) and with a reportable diagnostic.  A2c does NOT
%    become maintained, on the A1d model.
%
% 3. (iii) IS PRICED SEPARATELY FROM (ii) AND THE TWO ARE NOT THE SAME OBJECT.
%    (ii) is about ell_star = partial_u L_star; (iii) is about L_star itself,
%    through its anchor.  Neither implies the other: an L^2 bound on a function
%    does not bound its derivative (lem:levels(iv)), and a sup bound on a
%    derivative does not bound the function (the anchor is free).  The moment's
%    VARIANCE is governed by E[Wtilde^2 ||L_star||^2] and therefore by (i)+(iii),
%    not by (ii).  See the round-7 correction under A4(d).
%
% 4. A2'' IS NOT WRITTEN HERE, deliberately.  The round-2 range-inclusion
%    condition R(B) subset R(A*A) is empty unless sigma_q == 0 (memo Sec.3.5,
%    verified by the strategist-critic).  Do not reintroduce it.
%
\begin{assumption}[H1: Hilbert-scale link]\label{ass:hilbscale}
$\mathcal A_{\mathrm{cond}}:\mathcal V\to L^{2}(\Pi)$ has infinite rank and its singular values satisfy
$\sigma_j(\mathcal A_{\mathrm{cond}})\asymp j^{-a}$ for some $a>0$; the source set for $\vdagger$ is a
ball of the \emph{classical} Sobolev space $H^{s}(\mathcal X)$ --- for non-integer $s$, the interpolation
space between the neighbouring integer orders, \textbf{with no boundary condition imposed at either
endpoint of $\mathcal X$}; and the operator is \emph{matched} to that scale in the sense of
\citet[\S8.5]{engl2000}, which is here the following stated comparison: there is a finite $\mathsf c$,
not depending on $\zeta$, such that every $\zeta\in\mathcal V\cap H^{s}(\mathcal X)$ admits
\begin{equation}\label{eq:h1match}
\zeta\;=\;\big(\mathcal A^{\ast}_{\mathrm{cond}}\mathcal A_{\mathrm{cond}}\big)^{s/(2a)}h
\qquad\text{with }\ \|h\|_{L^{2}(\pi)}\;\le\;\mathsf c\,\|\zeta\|_{H^{s}} .
\end{equation}
\emph{Two notational points.} \emph{First}, bare sans-serif $\mathsf c$ is this project's \emph{generic
finite positive constant}: bound by the quantifier of the clause in which it appears, free to take a
different value in every such clause, and never carried from one result to another. This clause binds it,
in the words ``there is a finite $\mathsf c$, not depending on $\zeta$'', and no site downstream reads its
value.
\emph{Second, the pre-image is written $h$ and \textbf{not} $\varsigma$}, $h$ being the
\emph{domain}-side bound variable \cref{rem:spacesread}\textup{(1)} declares --- the counterpart of
$\Xi,\Xi'$ on the range side --- and the pre-image being exactly that: a generic element of $L^{2}(\pi)$, existentially
bound per $\zeta$ inside this clause. \emph{The object is the spectral source pre-image}, the same species
as in (P-a.3) and \cref{lem:scaleinv}\textup{(iii)}, read here at the \emph{population} operator on all of
$\mathcal V$.
\emph{$H^{s}$ is the \textbf{classical} Sobolev space and \textbf{not} a scale generated by
$\mathcal D$}: for a Hilbert scale generated by a
differential operator on a bounded interval the second scale step imports the generator's boundary
conditions, so the space the generator produces at order $s'$ coincides with $H^{s'}(\mathcal X)$
only for
$s'<\tfrac32$ (Neumann-type) or $s'<\tfrac12$ (Dirichlet-type) --- \emph{$s'$ throughout this assumption
and A8's derivation denoting a \textbf{generic} order of the scale, bound where it appears
and used nowhere else in this project} --- while \eqref{eq:lxwindow} forces
% ARCHAEOLOGY.  Through dispatch 21 the constant line read "the constant is the
% Sec.1.6 sans-serif slot and carries no name of its own"; that pointer was false --
% Sec.1.6 scopes bare sans-c to one named PROOF site (prp:nophifree Step 3) -- and
% the glossary was amended rather than the pointer patched (C88).  varsigma is a
% Sec.1.6 LOCAL and a local may not be made binding (C25); it had been made binding
% here, in a maintained assumption's statement.  eq:h1match is the matching written
% as a condition on one named pair of objects rather than as machinery, which is
% dispatch 21's change; through dispatch 20 the clause read "the Sobolev scale H^s
% of order s GENERATED BY D", and that phrase is struck.  The generator's space at
% order s' was written X_{s'} through dispatch 23 and is written out: italic X is
% the BINDING regressor X_1, X_2, subscripted by a period index, and a scale space
% subscripted by an order on the same letter is separated from it by nothing but
% what the subscript happens to be; swept d24, X_{s'} -> A 1 / R 0 / P 0.  s' was
% written sigma and tau through dispatch 21, against the four-way sigma rule and
% against the reserved tau_DP.  Glossary C87, C106.
$s>\tfrac{27}{16}$ at $\alpha=\tfrac38$ and \eqref{eq:hxwindow} forces more still. \emph{A B-spline on a
quasi-uniform mesh satisfies no boundary condition}, so under the literal generated reading
\eqref{eq:sieveinhs} would be \textbf{false} at every admissible $s$ and the four conditions it underwrites
--- A5(vi), A5(vi$'$), \cref{prp:phirate}(P-a.1) and \eqref{eq:foldrates}'s
second display --- would again be events on the \emph{empty} set, which is the failure the amendment of
A8 closed and which the generated reading reopened verbatim.
\emph{What held it together was a directive and not a condition}: one sentence in a scope note in
\texttt{proofs.tex} instructing that this scale be read on the same scale as \eqref{eq:jackson}, i.e.\
classically --- in a different file, and cited at neither this assumption, \cref{def:srcball}, nor
\eqref{eq:sieveinhs}. \emph{It is settled here instead, in the clause, once}: the space is the classical
one, and the comparison the generator was there to produce is \eqref{eq:h1match}, assumed outright.
\emph{Nothing downstream loses anything}: only the single index $s$ is ever used, no site reads a second
member of the scale, and \eqref{eq:h1match} is exactly what the generated construction was invoked to
deliver. \emph{Consequently the source index in A2 is}
\[
\mu_{\text{\textup{H1}}} \;:=\; \frac{s}{2a}.
\]
\emph{The subscripted name is used \textbf{only} where the two indices must be told apart, which is the
disposition \cref{rem:pasrc}\textup{(5)} makes binding.} Bare $\mu$ in this project is (P-a.3)'s
\emph{primitive} index at the finite-rank sieve compression of the \emph{realised} operator at the anchor
$\theta_\circ$; $\mu_{\text{\textup{H1}}}$ is the index of A2 at the \emph{population}
operator $\mathcal A_{\mathrm{cond}}$ on all of $\mathcal V$ at $\theta_\dagger$, and $s$ and $a$ in the
display are objects at $\theta_\dagger$ as well, against the $\theta_\circ$-specific $a$ of
\cref{def:crossfit}(b). \emph{They are equal only under the interpolation proviso this file states and
declines}, and the three sites that computed one and read it into a condition on the other --- the two
daggered lines of the smoothness-budget table, the sentence deducing $\mu>\tfrac56$ from $s>a+\tfrac12$,
and \cref{rem:rootnscope}\textup{(1)} --- are qualified in place.
\emph{\textbf{The same identity is \textbf{not} asserted for \texttt{prp:phirate}.}}
\texttt{prp:phirate}'s hypothesis (P-a.3) states a
% ARCHAEOLOGY.  The change of scope is the round-12 amendment to this assumption.
\textbf{spectral} source condition, at the sieve compression of the \emph{realised} operator, with $\mu$
and its radius $\mathsf R_\ast$ \emph{primitive}. What H1 supplies to that route is the \emph{scale}
$H^{s}$ --- which \cref{def:srcball}, A5(vi$'$), (P-a.1), (P-c) and (P-g) are written in ---
and nothing more; the passage from a Hilbert-scale source condition plus a link condition to a spectral
one is operator interpolation at exponent $2\mu$, licensed only at $2\mu\le1$, and
\texttt{prp:phirate} lives at $2\mu>\tfrac53$. \cref{rem:pasrc} carries the full disposition, including
the commutativity proviso under which the identity above \emph{would} extend to the sample operator, and
what that proviso costs.
\end{assumption}

% WHY H1 MATTERS, and it does -- BUT ROUND 12 NARROWS THE CLAIM, AND THE
% NARROWING IS THE CRITICAL FINDING OF THAT ROUND.
%
% WHAT H1 STILL DOES, and it is most of what the block uses it for:
%   (a) it FIXES THE SCALE H^s in which def:srcball, A5(vi'), (P-a.1), (P-c) and
%       (P-g) are written -- without it those five conditions name no norm;
%   (b) it supplies A2's PICARD FORM, via the singular system sigma_j asymp
%       j^{-a} of A_cond on V;
%   (c) it delivers the mild-ill-posedness and mu > 1/2 lines of the budget
%       table AS AN INTERPRETATION: s > a + 1/2 gives
%           mu = s/(2a) > (a + 1/2)/(2a) = 1/2 + 1/(4a) > 1/2,
%       with room, and mu > 1/2 is necessary for A5(iii)'s frontier to be
%       satisfiable at any r whatsoever.
%
% WHAT IT DOES NOT DO, AND ROUNDS 1-11 ASSUMED IT DID.  The clause
% "consequently the source index in prp:phirate is mu = s/(2a)" IS WITHDRAWN FROM
% THE RATE ROUTE.  Three facts force this and each is proved inside the block.
%   (1) (F-f) -- the classical Landweber/discrepancy theorem the rate rests on --
%       consumes a SPECTRAL source condition, v = (L* L)^mu varsigma with
%       ||varsigma|| <= R_*, at the operator being regularised.  H1 supplies a
%       HILBERT-SCALE condition plus a link condition.  Those are different
%       hypotheses.
%   (2) THE BRIDGE BETWEEN THEM IS OPERATOR INTERPOLATION AT EXPONENT 2 mu, AND
%       IT IS AVAILABLE ONLY AT 2 mu <= 1 -- Heinz's range, which is also the
%       interpolation range of engl2000 Sec. 8.5, THE VERY PIN H1 NAMES.
%       eq:muhalf forces 2 mu > 5/3; and lem:srctransfer(ii.a) CONSTRUCTS a 2x2
%       witness refuting operator monotonicity of z -> z^p at EVERY p > 1.  So no
%       configuration this block admits escapes, and the file already contains
%       the proof that the device is unavailable.
%   (3) FINITE-DIMENSIONALITY OF S_n DOES NOT RESCUE IT.  On S_n range membership
%       is vacuous -- every positive-definite operator there is surjective -- so
%       all the step's content is the uniform-in-n PRE-IMAGE NORM bound, and that
%       bound IS the interpolation inequality.  Finite rank makes the failure
%       quantitative rather than qualitative; it does not remove it.
%
% THE REPAIR, AND WHY IT IS A WEAKENING RATHER THAN A LOSS.  (P-a.3) now states
% the SPECTRAL condition directly, at the compression eq:compress, with mu and
% R_* PRIMITIVE.  The alternative -- keep the Hilbert-scale pair and add a
% COMMUTATIVITY proviso between A-hat* A-hat and the compressed scale operator --
% would preserve mu = s/(2a) exactly, and lem:srctransfer(ii.b) already proves the
% transfer under it.  BUT OPTION A = OPTION B + A PROVISO, so the hypothesis
% actually assumed here is STRICTLY WEAKER and needs strictly less proved.  The
% proviso is not diagnosable -- it asserts spectral alignment of a random,
% rank <= n_m operator built from 2n design points with a spline-Sobolev scale --
% and it would add a THIRD constant def:lfstep must print before the data are
% seen.  rem:pasrc(4) records it as the alternative, with its price stated.
%
% AND THE COLLAPSE WAS ALREADY PARTLY ILLUSORY, WHICH IS THE HONEST FORM OF THE
% VERDICT.  H1's matching is a statement about A_cond ON ALL OF V.  (P-a) is about
% a FINITE-RANK, SIEVE-COMPRESSED operator at a RANDOM base point, and
% lem:srctransfer(iii.a) records that a link condition CANNOT hold on V at all,
% because A-hat has finite rank (prp:hatAprop(ii)).  So H1 never delivered
% mu = s/(2a) for the operator (P-a) is about: (P-a) HAS BEEN ASSUMING ITS OWN
% MATCHING ALL ALONG.  Round 12 makes an existing assumption visible rather than
% adding one.
%
% CONSEQUENCE FOR ESCALATION ITEM 15, stated rather than absorbed.  It does NOT
% close favourably on the rate route: the parameter count there is (mu, s, a, r)
% with mu primitive at (P-a.3) and s, a entering through the scale, the budget
% table and A2 -- FOUR parameters, one of them soft.  That is exactly the cost
% the regularity-checker's third remedy names, and it is paid in the open at
% rem:pasrc(4) and rem:phirscope.  It DOES still close favourably for A2's Picard
% form and for the balancing route, neither of which invokes (F-f).
% Reference: Engl, Hanke & Neubauer (2000), Sec. 8.5.
%
% NOTE THE DOMAIN.  The singular system is that of A_cond restricted to V.  On
% all of L^2(pi) the constants lie in the kernel and the system would carry a
% zero singular value, which is not the system H1 or A2 are about.  The domain
% is fixed once, in def:spaces of results.tex.
%
% NOTE ALSO THE RANK, added in the cleanup pass.  "sigma_j asymp j^{-a} for all
% j" SILENTLY FORCES INFINITE RANK, and finite rank is not a pathology here:
% both counterexample-hunter DGPs produce a RANK-ONE A_cond satisfying every
% other hypothesis of the block.  Nothing else breaks -- A2's two-clause Picard
% form is well posed at finite rank, since the sum is then finite -- but the
% singular system spans (ker A_cond)^perp, NOT V, so with finite rank H1's
% hypothesis is simply false and H1 must not be invoked.  The infinite-rank
% requirement is therefore stated rather than smuggled in through the decay.

% ---------------------------------------------------------------------------
% PART IV.  DESIGN, OVERLAP, NORMALISATION, RATES
% ---------------------------------------------------------------------------

\begin{assumption}[A3: RETIRED --- local rank, now \cref{lem:momzero}\textup{(iv)}]\label{ass:rank}
\textbf{\sout{RETIRED.}} \emph{This assumption is not maintained, it is not available for use, and no
result in this programme takes it as a hypothesis. It is retired because it is a \textbf{theorem} --- not
because it is false, as A6 was, and not because it is too strong.} It read: \emph{suppose
A2 holds, so that $\ell_\star$ exists; then}
% ARCHAEOLOGY.  It was maintained through dispatch 12.
\[
J \;:=\; E[\Delta Z\,\Delta Z'] \;+\; E\big[\widetilde W\,\ell_\star(\Delta U,\Delta Z)\,\Delta Z'\big]
\quad\text{\emph{is nonsingular.}}
\qquad\text{\sout{[retired]}}
\]
\emph{Why it is retired.} \Cref{lem:momzero}\textup{(iv)} proves that under A1's \emph{first} clause the
second summand is the zero matrix --- $\ell_\star(\Delta U,\Delta Z)\Delta Z'$ is
$\sigma(\Delta U,\Delta Z)$-measurable and $E[\widetilde W\mid\Delta U,\Delta Z]=0$ --- so
$J=E[\Delta Z\Delta Z']$ whenever $J$ is defined at all, and nonsingularity of that matrix is A1b(iv). The
proof consumes A1, A1b(iv), A$2'$(i), A$2'$(iii) and B1(i), (ii), (iv.a) --- \emph{A1 for the vanishing and
A1b(iv) for the nonsingularity; A$2'$(i), A$2'$(iii) and the three B1 clauses only to make the second
summand a finite matrix in the first place, through \cref{cor:elldiff}}. \emph{Every one of those tags was
already carried by every list
that carried A3}, so striking the tag is a strict weakening of four hypothesis lists with no change to any
conclusion: \cref{lem:sandwich}\textup{(i)} and \textup{(i$'$)}, \cref{thm:rootn}'s group (T-2), and
\cref{cor:ratefloor} through (T-2).

\emph{What survives, and it is the half the conditional statement was invented to carry.} A3 was stated
conditional on A2 because $\ell_\star$ exists only under A2, so the display above is a well-formed object
only under A2 and ``A3 holds while A2 fails'' had no truth value. \emph{That remains exactly true of the
display} and \cref{rem:rootnscope}\textup{(4)} still says so, in the sharper form that what fails under
$\neg$A2 is the \emph{existence} of $\ell_\star$, $L_\star$ and $\psi_\star$ and not a rank condition. It
also remains true that A3 was never what identifies $\theta_\dagger$: that is A1b(iv). \emph{What is
withdrawn is the further claim, made at A1b(iv)'s comment and at three sites in
\texttt{results.tex}, that the two are \textbf{different} matrices.} They are the same matrix wherever
both are defined.
% ARCHAEOLOGY.  The round-2 correction (M-3) was right about the conditionality and
% wrong about the identity, and the identity was never computed.

\emph{What is \textbf{not} conceded.} No conclusion weakens, no rate moves, and no other tag changes
status. In particular A$2'$(i.b) is \textbf{not} freed --- see \cref{lem:momzero}\textup{(iv.1)} --- and
A2 is not freed, being consumed at \cref{lem:biasexp}, \eqref{eq:biasb1}, for the first-order bias, which
is where the source condition does its work. \emph{And the counterexample probe is redirected as it always
was}: the object to probe is A2.
\end{assumption}

% DISPATCH 13.  A3 IS RETIRED, AND THE RETIREMENT IS A STRENGTHENING.  The
% history matters more than the disposition and is recorded in the statement.
%
% WHAT THE THREE EARLIER ROUNDS SAID, IN ORDER.
%   Round 1  routed identification of theta_dagger through A3.  WRONG: that made
%            identification contingent on A2.
%   Round 2  (M-3) repaired it by routing identification through A1b(iv), giving
%            as the ground that A3 is "nonsingularity of a DIFFERENT matrix (J,
%            involving ell_star) and is itself conditional on A2".  THE
%            CONDITIONALITY HALF IS RIGHT.  THE "DIFFERENT MATRIX" HALF WAS
%            ASSERTED AND NEVER COMPUTED.
%   Dispatch 10  displayed J at eq:JSig for the first time and hardened the
%            uncomputed half into an ITALICISED claim inside def:moment: "J is
%            the matrix of A3 and it is NOT E[Delta Z Delta Z']".
% So a warning about CONDITIONALITY hardened, across three rounds, into a false
% claim about IDENTITY.  THE DIAGNOSIS IS A SIGMA-FIELD SCOPE SLIP of exactly the
% species (P4b)'s warning and (P4c) exist to police: E[Wtilde . | .] kills a
% sigma(Delta U, Delta Z)-measurable factor and does NOT kill a
% sigma(X_1, X_2)-measurable one.  The nearby object that genuinely does NOT
% vanish is q = E[Wtilde partial_u ell_star | X_1, X_2] of the glossary, which
% lem:secondord and prp:sharpquad consume; a reader who has q in mind will read
% E[Wtilde ell_star Delta Z'] as non-vanishing too, and that is the error.
%
% THE FILE ALREADY CONTAINED THE REFUTATION, AT TWO SITES, WHICH IS WHY THIS IS A
% BOOKKEEPING FAILURE AND NOT A MATHEMATICAL ONE:
%   lem:nusigma(v)     E[Wtilde Xi(Delta U, Delta Z)] = 0 for every integrable
%                      Xi, under A1.  PROVED since dispatch 1.
%   rem:rootnhyp(1)    names that identity as exactly what A1's first clause
%                      buys, in Part I itself.
% Neither was ever evaluated at Xi(u,z) = ell_star(u,z) z'.
%
% FOUR FAILURE MODES CHECKED BEFORE ACTING, ALL NEGATIVE.
%  (a) SELECTION.  Does ell_star depend on W through the SELECTION of ell^Pi_star
%      from ell^Pi_dagger + ker A*_cond?  It may depend on the JOINT LAW,
%      including the law of W.  IRRELEVANT: by eq:Wdef, ell_star = W ell^Pi_star
%      is a fixed Borel map (u,z) -> R^p, and a law-dependent choice of a
%      deterministic function is still a deterministic function.  Measurability
%      with respect to sigma(Delta U, Delta Z) is a property of the map's
%      ARGUMENTS, of which W is not one.
%  (b) N1.  Does the normalisation reintroduce W?  NO, twice over: N1 constrains
%      the ANCHOR b of def:weight(b) and leaves ell untouched, and eq:JSig
%      carries ell_star and not L_star.
%  (c) TRUNCATION.  A4(i) truncates varrho-hat to [eps, 1 - eps] at def:hatT, so
%      Wtilde-hat is NOT conditionally centred.  CORRECT, AND THE FINDING DOES
%      NOT EXTEND TO IT: J is defined at the TRUE varrho (def:moment says so in
%      its own text), A4(i) makes the truncation free THERE, and the second
%      summand of eq:Jhat does not vanish at any n.  lem:sandwich Step 3 pays its
%      cost and is unchanged.  Recorded at lem:momzero(iv.2).
%  (d) INTEGRABILITY for the tower step.  |Wtilde| <= 1 by (P1),
%      ||ell_star||_inf < oo by cor:elldiff, E||Delta Z|| < oo by A1b(iv) --
%      which is the SAME bound lem:momzero Step 3 already used to conclude
%      ||J|| < oo.  Nothing new is needed.
%
% WHAT DOES NOT FOLLOW, AND THE BLOCK WOULD COLLAPSE IF IT DID.  The vanishing
% does NOT extend to the varphi-direction: Delta h is sigma(X_1, X_2)-measurable
% and NOT sigma(Delta U, Delta Z)-measurable, so E[Wtilde ell_star Delta h] is
% not zero -- indeed A2 is precisely the assertion that it equals
% -E[Delta Z Delta h] for every h in V (prp:orthog(i), eq:orthphi).  Were it zero
% for every h one would get nu_{Delta Z} = 0, i.e. the instrument does not move
% the regressor.  NEYMAN ORTHOGONALITY IS UNTOUCHED: prp:orthog is about the
% three NUISANCE directions, not about partial_theta.  Sigma is untouched and in
% fact SPLITS, eq:sigsplit, with neither summand vanishing.

\begin{assumption}[A4: Overlap and propensity regularity]\label{ass:overlap}~
\begin{enumerate}
\item[(i)] \emph{Overlap.} There is $\varepsilon_\varrho\in(0,\tfrac12)$ with
$\varrho(\Delta Z)\in[\varepsilon_\varrho,\,1-\varepsilon_\varrho]$ almost surely.
\item[(ii)] \emph{Smoothness.} $\varrho$ belongs to a H\"older ball of order $r>0$ on the support of the
non-degenerate coordinates of $\Delta Z$, whose number is $d_Z$ with $1\le d_Z\le p-1$.
\emph{The lower bound $d_Z\ge1$ is load-bearing.} At $d_Z=0$ the exponent $\alpha=r/(2r+d_Z)$ is
$\tfrac12$ \emph{exactly}, and:
% ARCHAEOLOGY.  The lower bound is stated in dispatch 14; it was assumed silently
% before it, at four sites and in one place substantively.
\begin{itemize}
\item \emph{One site genuinely breaks, and it is the $\varphi$-step.} (P-d)'s sieve-entropy condition is
$L_x\ln n=o(n\hat\delta_n^{2})$. At $\alpha=\tfrac12$, A5(i)'s $\hat\delta_n=\max\{n^{-1/2}\sqrt{\ln n},r_n\}$
is $n^{-1/2}\sqrt{\ln n}$, so $n\hat\delta_n^{2}=\ln n$ and the condition reads $L_x=o(1)$ ---
\textbf{no integer $L_x\ge1$ satisfies it}. Equivalently, the window \eqref{eq:lxwindow} is empty at every
$s$, its right-hand side $9\alpha/\big(8(1-2\alpha)\big)$ diverging as $\alpha\uparrow\tfrac12$.
\cref{prp:phirate} therefore has no admissible sieve dimension, group (T-5) of \cref{thm:rootn} is
unsatisfiable, and the theorem is \emph{vacuous} at $d_Z=0$ rather than false.
\item \emph{Three further sites print ``$\alpha<\tfrac12$'' as a \textbf{reason} and are false as printed,
although each conclusion happens to survive at the boundary}: \cref{lem:prelimtheta}'s statement, which
writes $\alpha=r/(2r+d_Z)<\tfrac12$; that lemma's Step~6, whose collapse argument gives
$\kappa\alpha<\tfrac12$ from ``$\kappa<1$ and $\alpha<\tfrac12$'' --- true at $\alpha=\tfrac12$ from
$\kappa<1$ alone; and \cref{rem:h6bandwidth}'s ``impossible because $\alpha<\tfrac12$ always''. \emph{A
true conclusion resting on a false printed reason is the harder defect to find}, which is why they are
listed.
\end{itemize}
\emph{$d_Z=0$ is not a pathology and it is not being excluded as one}: it
says $\varrho$ has no non-degenerate argument, hence is a constant, hence is estimable at the parametric
rate --- a genuinely easier propensity problem. \emph{What fails there is not the propensity but the
$\varphi$-step's entropy budget}, because $\hat\delta_n$ then sits at $n^{-1/2}$ and leaves no room between
the noise level and the sieve dimension. \emph{A block that assumes a strict inequality should say which
primitive delivers it}, and this is that primitive.
\item[(iii)] \emph{Estimation --- and the measurability that makes the display a probability statement.}
$\hat\varrho$ is computed on an auxiliary fold; its \emph{own pointwise function}
$(\omega,z)\mapsto\hat\varrho(\omega)(z)$ --- the one \cref{def:hatT} evaluates at each $\Delta Z_i$ --- is
\textbf{jointly measurable}, with respect to the product of the $\sigma$-algebra generated by the
auxiliary fold with $\mathfrak B(\IR^{p})$; and
\[
\big\|\hat\varrho-\varrho\big\|_{L^2}=O_P\big(n^{-r/(2r+d_Z)}\big)=:O_P(r_n).
\]
\emph{The $\sigma$-algebra of the fold is written out and given no symbol here}, because
\cref{def:hatT} --- which names it $\sigma(\mathcal I_a)$ by adopting \cref{lem:balance}'s convention ---
comes several hundred lines later in this file, and a maintained clause must not name an undefined object;
this is the write-it-out precedent C12/C24/C78, applied for the same reason A7's own statement forced
\cref{def:hatT} into Part~V.
\emph{The joint-measurability half is load-bearing}: two things this clause is asked to support are not
well posed without it.
\emph{First, its own display.} ``$O_P$'' is a statement about the distribution of
% ARCHAEOLOGY.  Through dispatch 15 this clause carried NO measurability statement
% of any kind, and the gap was recorded as OPEN; the half is new in dispatch 16.
$\omega\mapsto\|\hat\varrho-\varrho\|_{L^2}$, which is an \emph{integral of a random function}; joint
measurability makes it a measurable function of the sample, by Tonelli's theorem applied to the
non-negative jointly measurable integrand $(\omega,z)\mapsto|\hat\varrho(\omega)(z)-\varrho(z)|^{2}$
against the law of $\Delta Z$. Per-realisation Borel measurability of $z\mapsto\hat\varrho(\omega)(z)$
does not supply it.
\emph{Second, \cref{lem:opnoise} Step~4}, which freezes $\hat\varrho$ on the auxiliary fold and applies
\textup{(P8)}'s freezing lemma. That lemma's hypothesis is joint measurability of the \emph{composite},
and the composite there is $|\hat\varrho(\omega)(\cdot)-\varrho(\cdot)|$ read at the main fold's
$\Delta Z$-draws, which the clause above supplies by composition with the Borel coordinate maps.
\emph{The shape is \cref{def:crossfit}\textup{(c$'$.1)}'s, one Part later}: a requirement on the
% ARCHAEOLOGY.  Through dispatch 15 lem:opnoise Step 4 was the one consumer of the
% freezing lemma whose hypothesis no clause of this block supplied; the repair was
% located here rather than papered over there.
estimator's own pointwise function and \textbf{not} a container declaration. That matters and is not a
stylistic preference --- no equivalence class is formed and no representative is chosen, so
\cref{def:hatT}'s truncation, being a pointwise operation, continues to make
$\varepsilon_\varrho\le|\widehat{\widetilde W}_i|\le1$ hold \emph{surely}, and
\cref{def:crossfit}\textup{(a)}'s reading of that bound at the folds is exact rather than almost sure.
\emph{And it costs nothing about $P$}: it is a construction requirement on an estimator, satisfied by any
$\hat\varrho$ given by a formula jointly measurable in the auxiliary sample and the evaluation point ---
which every kernel and every series construction is, by inspection of the formula. \emph{The
Carath\'eodory route \cref{def:crossfit}\textup{(c$'$.1)} records for \eqref{eq:prodmeas} is available
here too but is \textbf{not} the same statement}: it delivers joint measurability from continuity in one
argument and measurability in the others, so here it would run on continuity of
$z\mapsto\hat\varrho(\omega)(z)$ at each realisation --- true of a kernel or series construction and not
of a measurable selection, and in any case a sufficient condition and not the requirement, which is the
display above.
\end{enumerate}
\end{assumption}

% WHAT A4(i) DOES, AND -- ROUND-2 CORRECTION -- WHAT IT DOES NOT.
%
%  (a) IT MAKES THE TWO CONDITIONAL LAWS WELL DEFINED.  Both W-cells are
%      non-null given Delta Z = z for a.e. z, so F_1(.|z) and F_0(.|z) are
%      pinned down a.e. rather than arbitrary on a null set.  Used in
%      lem:quantu.
%
%  (b) IT CANCELS THE PREFACTOR.  At sigma(Delta Z), and only there,
%          E[Wtilde 1{Delta U <= u} | Delta Z = z]
%            = varrho(z)(1 - varrho(z)) { F_1(u|z) - F_0(u|z) },
%      and A4(i) bounds varrho(1-varrho) below by eps(1-eps) > 0, so the
%      conditional restriction can be divided through.  This is what makes
%      lem:bierens an EQUIVALENCE rather than an implication: without A4(i) the
%      restriction is vacuous wherever varrho in {0,1}.
%
%  (c) IT DOES *NOT* MAKE A_cond A REWEIGHTING OF THE RAW W-CONTRAST OPERATOR.
%      Round 1 asserted -- following plan Sec.18.3 -- that A_cond and the raw
%      W-contrast operator "differ by the bounded, boundedly-invertible weight
%      varrho(1-varrho), hence have the same range and the same kernel".
%      THAT IS FALSE AND IS WITHDRAWN.  The identity in (b) holds only at
%      sigma(Delta Z).  The operator's kernel requires conditioning on
%      (X_1, X_2, Delta Z), where, writing p_X := P(W = 1 | X_1, X_2, Delta Z),
%          E[Wtilde f_{Delta Y | X_1,X_2,Delta Z,W}(u) | X_1,X_2,Delta Z]
%             = p_X (1 - varrho) f_1(u) - (1 - p_X) varrho f_0(u),
%      which equals varrho(1-varrho){f_1 - f_0} ONLY IF p_X = varrho, i.e. only
%      if W is independent of (X_1,X_2) given Delta Z.  That says THE INSTRUMENT
%      DOES NOT SHIFT THE ENDOGENOUS REGRESSOR GIVEN Delta Z, and it MUST NOT BE
%      ASSUMED.  See rem:draftrec(ii) in results.tex
%      for the correct comparison, which is between the two operators' KERNELS
%      and claims neither equality of range nor equality of kernel.
%      CONSEQUENCE: plan Sec.18.3's bullet "A4's mechanism" is refuted as
%      stated, and the reconciliation it asked for is delivered in the weaker,
%      correct form.  Escalated in theory_memo.md Sec.7, item N-7.
%
%  (d) IRREGULARITY -- AND THE REFERENT IS L_star, NOT ell_star.  ROUND-7
%      CORRECTION.  As overlap degrades, A2 can continue to hold AS A RANGE
%      STATEMENT while the moment's variance diverges: that is irregular
%      identification in the sense of Khan & Tamer (2010).  THE VARIANCE OBJECT
%      IS E[Wtilde^2 ||L_star||^2], because psi = Delta Z Delta U + Wtilde L_star
%      carries L_star and not its derivative.  Rounds 1-6 named
%      E[Wtilde^2 ell_star^2] here -- THE DERIVATIVE -- which is the object
%      lem:balance(i) and lem:secondord(iv) consume through eps_n, and is a
%      different quantity.  This is the same referent slip that A5(iv)'s round-5
%      correction fixed in the opposite direction (there the block priced the
%      LEVEL where the consumers needed the DERIVATIVE; here it priced the
%      DERIVATIVE where the variance needs the LEVEL).
%
%      WHAT SURVIVES -- REWRITTEN IN DISPATCH 14, BECAUSE THE ROUND-7 VERSION OF
%      THIS PARAGRAPH CONTAINED TWO FALSE STATEMENTS AND ONE LEGITIMATE GROUND.
%      The two false ones are struck; the legitimate one is kept and is now the
%      only ground stated.
%
%        FALSE 1 (struck).  "A2'(ii) stays as written, because it is the clause
%        the second-order and bias terms consume."  A2'(ii) IS CONSUMED BY NO
%        RESULT IN THIS BLOCK.  results.tex says so three times -- at
%        lem:balance's hypothesis line, at lem:solset and at lem:momzero's closing
%        paragraph -- and the ROUND-8 CHANGE LOG OF THIS VERY FILE says so:
%        "A2'(ii) is consumed by none of them and is a Step-8 clause."  A comment
%        block asserting the opposite of the change log twelve hundred lines below
%        it is the defect, not the citation.
%
%        FALSE 2 (struck).  The bare "(i) + (iii)" in the sentence beginning "What
%        the irregularity discussion now points at" was ambiguous between A2'(i),
%        A2'(iii) and -- inside a comment attached to A4 -- A4(i), A4(iii); read
%        the second way it attributes the finiteness of E[Wtilde^2 ||L_star||^2]
%        to OVERLAP, and THAT IS WRONG.  A4(i) PLAYS NO ROLE IN THE VARIANCE
%        BOUND.  The bound is
%            E[Wtilde^2 ||L_star||^2] <= 2 E||b_star||^2
%                                        + 2 ||ell_star||_inf^2 E[Delta U^2],
%        whose FOUR ingredients are: |Wtilde| <= 1, which is (P1) and is FREE
%        from W in {0,1} and varrho in [0,1] -- the NOTE below this comment has
%        said so since round 2; ||ell_star||_inf < infinity, which is cor:elldiff
%        under A2'(i) IN FULL -- (i.a) AND (i.b), since cor:elldiff is stated at
%        A2b and A2b requires both, and since A2'(i)'s closing clause is what
%        makes ell_star pointwise well defined at all -- with B1(i), (ii),
%        (iv.a); E||b_star||^2 < infinity, which is A2'(iii); and
%        E[Delta U^2] < infinity, which is A1b(iv).  The tags are therefore
%        A2'(i), A2'(iii), A1b(iv), B1(i), (ii), (iv.a) -- and NOT A4(i).  The
%        counterexample-hunter reached this from the other side and put it in one
%        line: "A4(i) plays no role in the variance bound."
%
%        WHAT SURVIVES, AND IT IS ENOUGH TO RETAIN THE POINT.  Overlap does bear
%        on the SIZE of the variance and not on its finiteness, and the mechanism
%        is exact rather than qualitative: with the two conditional laws given W
%        held fixed and varrho == eps, A_cond^{(eps)} = eps(1-eps) B with B free of
%        eps, so RANGES ARE INVARIANT and A2 holds at every eps > 0 or at none --
%        the range statement overlap does not touch -- while
%        ell_star^{(eps)} = {eps(1-eps)}^{-1} ell^{(eps)} and, under A1,
%        E[Wtilde^2 ell_star^2] = E[(ell^{(eps)})^2]/{eps(1-eps)} asymp
%        eps_varrho^{-1} -> infinity.  That is Khan-Tamer irregularity, it lives
%        where A4(i) FAILS, and it is why A4(i) is retained.  NOTE THE OBJECT: the
%        divergence above is displayed for E[Wtilde^2 ||ell_star||^2], the
%        DERIVATIVE, which is A2'(ii)'s object; the same computation runs for
%        E[Wtilde^2 ||L_star||^2], the LEVEL, through ||ell_star||_inf.
%
%      A2'(ii) IS THEREFORE RETIRED IN PLACE (see the clause itself), and the only
%      ground the block ever had for keeping it -- a FORWARD Step-8 need at
%      prp:effbound, where the variance is computed and where the T3 tier relaxes
%      A2'(i) -- is a forward obligation and not a present consumer.  It is
%      recorded at the retirement note and flagged for Step 8.
%
%  NOTE.  A4(i) is NOT needed for |Wtilde| <= 1, which is free from W in {0,1}
%  and varrho in [0,1].  Round 1 listed A4(i) as a hypothesis of prp:wellposed,
%  prp:frechet and prp:compact, where it is idle; it has been removed from all
%  three.
%
%  (e) DISPATCH 7: A4(i) NOW ALSO LICENSES A CONSTRUCTION, AND THE CONSTRUCTION
%      IS WHAT KEEPS THE SAMPLE OPERATOR NON-DEGENERATE.  def:hatT truncates
%      varrho-hat to [eps_varrho, 1 - eps_varrho].  A4(i) is what makes the
%      truncation FREE -- varrho already lies in that interval, so the metric
%      projection is a pointwise contraction and A4(iii) is preserved -- and
%      what makes it USEFUL: |Wtilde-hat_i| >= eps_varrho SURELY, for every i.
%      WITHOUT IT the block has a real gap, and it is not a formality.  A4(i)
%      constrains varrho and says NOTHING about varrho-hat, so with an untruncated
%      (or [0,1]-truncated) first stage the event Wtilde-hat_i = 0 has positive
%      probability at each i.  Two consequences, of different sizes.
%        - TOTAL degeneracy, Wtilde-hat_i = 0 for EVERY i in I_m, makes T-hat
%          identically zero, on which event every statement about it -- including
%          prp:tccfails -- is vacuous.  This event is negligible under overlap and
%          cross-fitting, but "negligible" is a theorem, not a convention.
%        - PARTIAL degeneracy is NOT negligible.  For any first stage whose raw
%          fitted values leave [0,1] on a set of probability of order r_n^2 -- a
%          series linear-probability fit, say -- the expected number of main-fold
%          units with Wtilde-hat_i = 0 is of order n r_n^2 = n^{1 - 2 alpha},
%          which DIVERGES for every alpha < 1/2.  Each such unit annihilates
%          g-hat_{.,.,i}, so a row and a column of the Gram matrix eq:gram vanish
%          and lambda_min(G_n) = 0 EXACTLY.  prp:tccrepair(ii)'s nonsingularity
%          hypothesis then fails and eq:tccrad returns a radius of zero, so the
%          file's only demonstration that A6b is non-vacuous would have no content
%          on an event of non-vanishing probability.
%      The truncation removes the exact-zero mechanism DETERMINISTICALLY.  It does
%      NOT make lambda_min(G_n) bounded below -- near-collinear g-hat_i remain
%      possible, which is Q-3 and is a simulation question (probe CP-8).
%
%  (f) AND A4(iii)'s AUXILIARY FOLD IS LOAD-BEARING FOR THE SAME PROPERTY, not
%      only for the rate.  Drop cross-fitting -- estimate varrho on the full
%      sample with any rule that interpolates -- and varrho-hat(Delta Z_i) = W_i
%      for every i almost surely, so T-hat is identically zero and prp:tccfails
%      is FALSE with probability one rather than on a null set.  The fold split
%      is therefore listed in prp:tccfails's hypotheses and not only in
%      def:hatT's convention.
%
%  d_Z is the EFFECTIVE dimension: Delta Z contains a constant, which
%  contributes nothing to a nonparametric regression, so d_Z <= p - 1.

\begin{assumption}[N1: Conditional centring of a weight]\label{ass:centering}
A debiasing weight $L$ \emph{satisfies N1} if
$E\big[L(\Delta U,\Delta Z)\mid\Delta Z\big]=0$ almost surely. \emph{N1 is a property that a weight may
or may not have; it is not maintained, and a result that needs it says so.}
\end{assumption}

\begin{assumption}[N1$'$: Conditional centring of a direction]\label{ass:centeringdir}
A perturbation direction $D$ in the weight \emph{satisfies N1$'$} if
$E\big[D(\Delta U,\Delta Z)\mid\Delta Z\big]=0$ almost surely.
\end{assumption}

% N1'S STATUS IS FIXED HERE, ONCE, AND EVERY OTHER SITE CITES THIS ONE.
% ROUND-7 DECISION.  Through round 6 the tag N1 carried THREE INCOMPATIBLE
% STATUSES across five sites -- a CONSTRUCTION here and in the glossary ("L IS
% NORMALISED so that ..."), an ASSERTION in the Discussion ("N1 costs nothing and
% BUYS orthogonality"), and an unused HYPOTHESIS on lem:secondord, lem:solset and
% lem:balance -- plus, at lem:secondord(iii), a fourth use in which the tag was
% applied to the DIRECTION D rather than to the weight L.
%
% THE COST OF THE CONSTRUCTION READING IS THAT IT VOIDS THE RESULT N1 EXISTS FOR.
% If L_star satisfies N1 BY DEFINITION then prp:orthog(iii)'s "only if" has no
% content, and the sentence "do not assume N1" in its statement is incoherent --
% one cannot suspend a definitional normalisation.  That is defect M-8, which
% round 6 repaired at prp:orthog and at no other site.  ONLY THE PROPERTY
% READING KEEPS (iii) NON-VACUOUS, and (iii) is the load-bearing claim of the
% whole rate architecture: it is what demotes varrho from a root-n nuisance to a
% product-rate one.  So the property reading is adopted, here, and:
%   - N1 is STRUCK from lem:secondord, lem:solset and lem:balance, where it was
%     listed and never consumed;
%   - N1 is NOT assumed anywhere in prp:orthog, which characterises it;
%   - the direction condition gets its own tag, N1', so that the tag N1 names
%     exactly one object;
%   - lem:density(iii) reads "the subspace of weights satisfying N1", which is
%     the property reading and needed no change;
%   - cor:udep continues to ASSUME N1, which is now a genuine hypothesis.
%
% WHAT REMAINS TRUE, AND IS WHY N1 IS CHEAP RATHER THAN FREE.  By def:weight the
% weight attached to a representer ell is determined by its anchor b, and
% re-anchoring at b = -E[ int_0^{Delta U} ell(v, Delta Z) dv | Delta Z ] produces
% a weight satisfying N1 with the SAME ell -- so it changes neither partial_u L
% nor, by (P3) and A1, the value of E[psi].  A2'(iii) shows that anchor is
% admissible.  So N1 CAN ALWAYS BE ARRANGED; it does not HOLD automatically, and
% the difference is exactly what prp:orthog(iii) measures.  Perturbing
% varrho -> varrho + tau_3 r gives d/d tau_3 E[psi] = -E[r(Delta Z) L], which
% equals -E[r(Delta Z) E[L | Delta Z]] and vanishes for every bounded r under N1
% and only under N1.  Without it varrho-hat needs a root-n rate.

\begin{assumption}[A5: Nuisance rates]\label{ass:rates}~
Let $h:=\hat\varphi-\vdagger$, and let $\hat\delta_n$, $\delta_n$, $\epsilon_n$, $r_n$ be as in the
notation glossary.
\begin{enumerate}
\item[(i)] \emph{Operator noise.}
$\hat\delta_n=\max\{\hat\delta_n^{(T)},\hat\delta_n^{(\varrho)}\}$ with
$\hat\delta_n^{(T)}=O_P\big(n^{-1/2}\sqrt{\ln n}\big)$ and $\hat\delta_n^{(\varrho)}=O_P(r_n)$;
in every configuration of interest $\hat\delta_n=\hat\delta_n^{(\varrho)}=O_P(r_n)$.
\item[(ii)] \emph{Structural rate, with the logarithm of \eqref{eq:phirate} carried \textbf{in the
symbol}.} $\|h\|_{L^2(\pi)}=O_P(\delta_n)$ with
\[
\delta_n\;:=\;\big(\hat\delta_n\ln n\big)^{2\mu/(2\mu+1)} ,
\]
and $\delta_n=o_P(n^{-1/4})$.
\emph{The factor $\ln n$ is not optional}: \cref{prp:phirate} delivers \eqref{eq:phirate} at
$(\hat\delta_n\ln n)^{2\mu/(2\mu+1)}$, and (P-f),
\cref{rem:h6bandwidth}\textup{(0)} and \cref{prp:stopfeas}\textup{(iv)} all read that log-inflated
$\delta_n$. \emph{A log-free reading of this clause is off by $(\ln n)^{2\mu/(2\mu+1)}$.}
\Cref{rem:rootnlog} carries the full disposition, including which of the four rate products
% ARCHAEOLOGY.  The factor is new in dispatch 10 and is a CORRECTION, not a
% strengthening: rounds 1--9 wrote this clause at delta-hat_n^{2mu/(2mu+1)}, so the
% clause the block STATED was not the object the proposition DISCHARGES, and the
% block was already inconsistent with itself, not merely with the proposition.
\cref{thm:rootn} needs are affected. In one line: the second half of this clause and clause (iii) are
\textbf{unchanged}, because \eqref{eq:frontier} is a strict inequality between polynomial exponents and its
strictness absorbs the logarithm; clause (iv) becomes \textbf{stronger} by the factor $(\ln n)^{2\mu/(2\mu+1)}$,
and that is the only place the logarithm bites; clause (v) becomes \textbf{weaker}, its right-hand side being
inflated; and \eqref{eq:muhalf} does not move, having been computed at the inflated $\delta_n$ already.
\item[(iii)] \emph{Propensity frontier.} Writing $\hat\delta_n=n^{-r/(2r+d_Z)}$, condition (ii) holds
if and only if
\begin{equation}\label{eq:frontier}
\frac{2\mu}{2\mu+1}\cdot\frac{r}{2r+d_Z}\;>\;\frac14
\qquad\Longleftrightarrow\qquad
r\;>\;\frac{d_Z\,(2\mu+1)}{2\,(2\mu-1)},
\end{equation}
which in particular requires $\mu>\tfrac12$, whatever the value of $r$.
\item[(iv)] \emph{Debiasing weight, on the derivative.} The estimator $\hat\ell$ of the representer, a
random element of $L^{2}(P_{\Delta U,\Delta Z})$ measurable with respect to the $\sigma$-algebra generated
by the auxiliary fold and read at one fixed \emph{jointly} measurable representative, satisfies
\[
\big\|\hat\ell-\ell_\star\big\|_{L^{2}(P)}=O_P(\epsilon_n),
\qquad \delta_n\,\epsilon_n=o_P\big(n^{-1/2}\big).
\]
\emph{That fold and that measurability are fixed by \cref{lem:balance}'s sample-splitting convention,
which also \emph{names what the measurability is of}; and the naming is not decoration.} The
first display is an $O_P$ statement about $\omega\mapsto\|\hat\ell-\ell_\star\|_{L^{2}(P)}$, an integral of
a random function; joint measurability of the fixed representative makes it a measurable function of the
sample by Tonelli's theorem, and \emph{membership plus an unspecified ``measurability'' does not}.
\emph{Nothing here is strengthened}: the joint measurability
% ARCHAEOLOGY.  Through dispatch 15 the convention this clause points at stated
% exactly membership and a measurability whose target was not named, so the clause
% was well posed only by reference to something that did not state it -- the same
% defect as A4(iii)'s and lem:balance's own, written a third time and at one remove.
% The repair is at lem:balance and the naming half is dispatch 16's.
is scoped to what \cref{lem:balance} Steps~0 and~2 use, which is strictly less than
\cref{def:crossfit}\textup{(c$'$.1)}'s \eqref{eq:prodmeas}, because \cref{lem:balance} evaluates $\hat\ell$
at an undistorted $(\Delta U,\Delta Z)$ draw. \emph{In Part~I it is \eqref{eq:prodmeas} that governs and
not this convention}: what \cref{def:crossfit}\textup{(d)} reads at the folds is this clause's
\emph{display}, and the measurability of $\hat\ell^{(-k)}$ is \cref{def:crossfit}\textup{(c$'$.1)}'s own.
\emph{What this clause does and does not say about $\hat L$.} It prices
$\hat\ell$, and nothing else. The identity
$\|\partial_u(\hat L-L_\star)\|_{L^{2}(P)}=\|\hat\ell-\ell_\star\|_{L^{2}(P)}$ is \textbf{an assertion,
not notation}, and is \emph{not} made here: it holds if and only if $\hat L$ is differentiable in $u$ at
every point with $\partial_u\hat L=\hat\ell$, that is, if and only if $\hat L=L^{\,\hat b}_{\hat\ell}$ in
the sense of \cref{def:weight}(b) --- which is exactly the construction obligation \cref{rem:dirmember}
records as owed by Step 6 or 7 and \emph{undischarged in this block}. \cref{lem:balance}(i)--(ii) consume
the clause in the form printed above, at $\hat\ell$ itself, and need nothing further;
\cref{lem:secondord}(iv) consumes it at $\|\partial_uD\|_{L^{2}(P)}$ with $D=\hat L-L_\star$, and may
therefore be instantiated at the estimated direction only \emph{after} that obligation is discharged.
% ARCHAEOLOGY.  Rounds 1--7 printed the identity as the first of two equalities
% inside the clause, which asserted the construction silently -- a forward
% obligation migrated into the block unstated.  Round-8 correction.
\item[(v)] \emph{Plug-in.} $\sqrt{\hat m}\,n^{-1/2}=o_P(\delta_n)$.
\item[(vi)] \emph{Sobolev stability of the first-stage error, at the final iterate.}
$\|h\|_{H^{s}}=O_P(1)$, with $s$ as in H1. \emph{The clause is a statement about a
finite quantity, and what makes it one is \eqref{eq:sieveinhs} together with H1}:
$h=\hat\varphi-\vdagger$ with $\hat\varphi\in\mathcal S_n$ at the configuration
\cref{lem:pathcont} states, so $\hat\varphi\in H^{s}$ by \eqref{eq:sieveinhs} and $\vdagger\in H^{s}$
by H1's source set; without the first of those $\|h\|_{H^{s}}=+\infty$ surely and the
$O_P(1)$ statement would hold on the \emph{empty} event rather than with probability tending to one.
\item[(vi$'$)] \emph{Sobolev stability along the linearisation path, as a w.p.$\to1$ envelope at a
\textbf{fixed} radius.} With $\hat\varphi^{(0)},\dots,\hat\varphi^{(\bar k)}$ the stages of
\cref{def:lfstep} and $\bar k$ finite, there is one finite $\mathsf R$, not depending on $n$, such that
% ARCHAEOLOGY.  This clause named only the first membership through dispatch 20.
\[
P\Big(\max_{0\le k\le\bar k}\big\|\hat\varphi^{(k)}-\vdagger\big\|_{H^{s}}\;\le\;\mathsf R\Big)
\;\longrightarrow\;1 .
\]
\emph{Both memberships are needed for that display to be a probability statement about a finite
quantity}: every stage lies in $\mathcal S_n$ once $\hat\varphi^{(0)}$ does, by
\cref{def:lfstep}\textup{(b)},\textup{(d)}, and hence in $H^{s}$ by \eqref{eq:sieveinhs}, while
$\vdagger\in H^{s}$ by H1's source set --- the display being a norm of a \emph{difference}, either
membership failing makes it $+\infty$ at every realisation and the event \textbf{empty} rather than
merely improbable.
\emph{That same $\mathsf R$ serves at every site \cref{def:srcball} enumerates, and the enumeration is
\textbf{there and nowhere else}}: \cref{def:srcball} carries the list, this clause and
\cref{lem:pathcont} and \cref{rem:paenvelope} point at it, and no second copy is maintained.
\emph{The quantifier is an existential over
% ARCHAEOLOGY.  Four copies of one list were maintained through dispatch 20, of
% lengths seven, five, five and four, so that a reader assembling the consumer set
% from any one of them got a different answer on a constant def:srcball declares
% "the radius that is not free".
one radius meeting all of those demands at once}, which is what \cref{lem:pathcont} consumes and what
\cref{rem:paenvelope} explains; an $O_P(1)$ reading does not serve, because \cref{def:srcball}'s ball must
be deterministic and an $O_P(1)$ bound is a family indexed by tolerance rather than a number.
\end{enumerate}
\end{assumption}

% THE FOLD READING OF A5(ii), A5(iv), A5(vi) -- DECLARED IN dispatch 11, IN
% results.tex def:crossfit(d), AND RECORDED HERE BECAUSE THIS IS THE CLAUSES' HOME.
% All three clauses above are stated at a SINGLE estimated nuisance: h is
% varphi-hat - varphi_dagger, ell-hat is one function.  Part I's estimator
% (def:crossfit) computes F of each, one per fold complement, so every Part-I
% result that cites them consumes them at max_{k <= F}.  THAT IS eq:foldrates,
% and it is NOT a silent widening:
%    A5(ii) first half and A5(vi) at the folds are DERIVED, in thm:rootn Step 0,
%      from (T-5) with F fixed and n_{-k} asymp n -- so they cost nothing;
%    A5(iv) at the folds is a STRENGTHENING IN FORM of the clause as written
%      (same order for fixed F, but a different statement), and it is therefore
%      carried by hypothesis (R-e.1) of thm:rootn EXPLICITLY, not read into this
%      clause.
% A4(iii) received the same treatment one clause earlier, in def:crossfit(a), and
% (R-a) and (R-e.2) already wrote max_k.  THE POINT OF RECORDING IT HERE: a reader
% who meets A5(iv) in this file and then meets it cited inside Part I must be able
% to see that the two are not the same statement and where the difference is paid.

% A5(vi) IS NEW IN ROUND 7 AND IS A TAG FOR AN INPUT THAT WAS ALREADY BEING
% CONSUMED UNTAGGED.  lem:balance(iii) applies the spline direct theorem to
% zeta = h = varphi-hat - varphi_dagger and needs ||h||_{H^s} = O_P(1); through
% round 6 that condition sat in the lemma's statement as the last untagged
% substantive hypothesis in Part F.  Two things must be said about it.
%   (a) IT IS NOT IMPLIED BY THE SOURCE CONDITION.  A2 and H1 constrain
%       varphi_dagger; h is an LF ITERATE, and its Sobolev norm is a property of
%       the regularisation path, not of the truth.  Recorded as N-18(a).
%   (b) ROUND-9 TEXT, NOW STRUCK: "IT IS CONSUMED ON ONE ROUTE ONLY.
%       lem:balance(ii), the rate route, does not use it..."  THAT IS FALSE FROM
%       STEP 6 ONWARD.  lem:pathcont needs the H^s bound to place the iterate in
%       the SOURCE-CONDITION ball V_n of def:srcball, and it needs it AT EVERY
%       STAGE, not at the final one.  So the Route column changes to "both" AND
%       the content changes -- which is why (vi') is added rather than (vi)
%       merely re-routed.
% A5(vi') IS NEW IN DISPATCH 7 AND IT IS NOT A RESTATEMENT OF (vi).  (vi) prices
% the FINAL iterate; lem:pathcont(c) and prp:phirate Step 5 consume the maximum
% over k = 0, ..., bar-k.  (vi) is (vi') at k = bar-k, so (vi') implies (vi) and
% not conversely.  bar-k is FINITE and does not grow with n (prp:phirate), so a
% finite union of O_P(1) statements is O_P(1) and (vi') is not a uniformity
% requirement over a growing family -- which is the only reason it is affordable
% as a maintained clause.
% WHAT (vi') IS NOT.  It is not discharged by prp:phirate.  prp:phirate delivers
% the L^2(pi) half of containment; the H^s half is assumed, here.  The
% alternative -- Bernstein's inequality for splines on S_n, giving
% ||zeta||_{H^s} <= c L_x^s ||zeta||_{L^2(pi)} -- would derive it at the cost of
% L_x^s delta_n = O(1), a genuine extra restriction that is NOT imposed.
% It is stated as an assumption rather than left inline so that the
% assumption-auditor sees a declared condition with a named consumer.

% A5(i) IS CONFIRMED BY lem:opnoise, AND ONE HALF OF IT IS SHARPER THAN WRITTEN
% -- BUT THE SHARPENING IS TUNING-DEPENDENT AND MUST BE READ WITH C5'(i).
%  (a) WHAT IS SHARPER.  The T-component's L^2(Pi) NORM is O_P(n^{-1/2}) with NO
%      logarithm (lem:opnoise(i)): L^2(Pi) is a Hilbert space, the summands are
%      bounded in norm, and a second-moment computation plus Chebyshev suffices.
%      No uniformity in (u,t) is taken, so no logarithm is bought.  A5(i)'s
%      sqrt(ln n) is therefore satisfied with room.
%  (b) IT IS NOT IDLE.  prp:stopfeas(i) consumes the sqrt(ln n) to turn a
%      second-moment bound into a FIXED-CONSTANT exceedance statement: Chebyshev
%      at s = tau n^{-1/2} sqrt(ln n) leaves O(1 / ln n) -> 0.  Delete the
%      logarithm and the statement holds only at a tau depending on epsilon.
%  (c) THE GUARD, AND IT IS WHY (a) IS NOT WRITTEN INTO THE CLAUSE.  The
%      SMOOTHING BIAS is O(h_u^{p_K}) and under the ROUND-9 tuning
%      h_u ~ n^{-1/16} that is exactly n^{-1/2} = o(n^{-1/2} sqrt(ln n)).  Under
%      the AMENDED tuning of C5'(i) below, h_u ~ delta-hat_n^{1/p_K}, the bias is
%      asymp delta-hat_n and is NO LONGER o(n^{-1/2}) -- it is the largest term
%      in eq:opnoise, by design.  So the two clauses interact: the T-component
%      closes at a fixed constant AGAINST delta-hat_n under EITHER tuning, and
%      against the sharper delta-hat_n^{(T)} = n^{-1/2} sqrt(ln n) only under the
%      round-9 one.  prp:stopfeas(i) is stated at delta-hat_n for that reason.
%      DO NOT HARD-CODE A BANDWIDTH INTO A5(i).

% A5(iv) IS PRICED ON THE DERIVATIVE, AND THAT IS A ROUND-5 CORRECTION.  Rounds
% 1-4 wrote it as ||L-hat - L_star||_{L^2} = O_P(eps_n).  BOTH CONSUMERS NEED
% THE DERIVATIVE INSTEAD:
%    lem:balance(i)-(ii) bound ||nu_g / pi|| by 2 bar-c_pi^{1/2} ||ell-hat -
%        ell_star||, because the first-order bias is generated by ell = partial_u L
%        and not by L;
%    lem:secondord(iv) bounds the (varphi, L) cross term by
%        2 bar-c_pi^{1/2} ||partial_u D|| ||h||, again the derivative.
% AN L^2 RATE FOR A FUNCTION DOES NOT TRANSFER TO ITS DERIVATIVE -- the same
% one-directional transfer lem:levels(iv) records for D^{-1}_0, here in the
% nuisance rather than in the structural direction.  So the round-4 A5(iv) was
% priced on an object no consumer uses, and stating the derivative version is
% not a strengthening for its own sake but a correction of the referent.
% NOTE what is NOT claimed: the LEVEL rate ||L-hat - L_star|| is not implied by
% the derivative rate either, and nothing at Step 5 needs it.  If a later step
% needs the level rate it must be added as a separate clause, not inferred.
%
% ROUND-7 POSTFIX: NO A5 CLAUSE IS ADDED FOR def:dirclass MEMBERSHIP, AND THE
% DECISION IS RECORDED HERE BECAUSE THIS IS WHERE A READER WILL LOOK FOR IT.
% lem:secondord(iii)-(iv) are stated at a FIXED admissible direction pair, and
% Step 6 or 7 will read them at D = L-hat - L_star, r = varrho-hat - varrho,
% h = varphi-hat - varphi_dagger.  That instantiation needs MEMBERSHIP in
% def:dirclass -- ||D||_{0,h} < infinity, and for (iv) differentiability with
% ||D||_{1,h} < infinity -- which is a DIFFERENT statement from any rate, and
% which A5 does not contain.  Three reasons it does not become a clause:
%   (a) [CORRECTED IN DISPATCH 11 -- THE ROUND-7 GROUND WAS FALSE AND THIS FILE
%       REFUTES IT NINETY LINES ABOVE.]  The round-7 text read "every clause of
%       A5 prices a RATE IN n; these are qualitative finiteness statements and an
%       O_P statement about them prices the wrong kind of object."  A5(vi) IS
%       ||h||_{H^s} = O_P(1) -- qualitative, no rate in n -- and A5(vi') is a
%       w.p.->1 envelope at a fixed radius.  So "A5 prices rates only" is not a
%       ground this block can use, and it was used twice more downstream
%       (results.tex rem:dirmember(a), rem:obligations) before dispatch 11 caught
%       it.  THE CORRECT DISCRIMINANT: every clause of A5 constrains an object
%       THE MAINTAINED BLOCK ITSELF INTRODUCES, at a single estimator -- h,
%       ell-hat, delta-hat_n, varrho-hat.  ||D||_{0,h} and ||D||_{1,h} are
%       seminorms of def:dirclass evaluated at a DIRECTION PAIR no maintained
%       clause fixes, so a clause about them would be a clause about a
%       definition's argument.  See rem:rootnhyp(5) for the same correction
%       applied to (R-a)-(R-e) of thm:rootn;
%   (b) both seminorms are indexed by (h, tau-bar) with h RANDOM, so a clean
%       clause would have to quantify over a deterministic class containing h
%       w.p.a.1 -- the ball V_n of L-19, which Step 6 constructs and which does
%       not exist in this block yet;
%   (c) it is IMPLIED, not assumed: rem:dirmember's eq:dirmember shows both
%       seminorms finite for EVERY h in V and every tau-bar as soon as L-hat is
%       an admissible weight (def:weight(b)) generated by a bounded, u-continuous
%       ell-hat.  Assuming an implied condition is the addition this round
%       removed everywhere else.
% SO WHAT A5(iv) PRICES IS ||ell-hat - ell_star||_{L^2(P)} AND NOTHING ELSE.
% SCOPE OF THE CLAIM -- ROUND-8 NARROWING (auditor Finding 7).  The round-7 text
% here read: "... so that NO BOUND IN THE BLOCK is denominated in a norm this
% assumption does not price."  READ LITERALLY THAT IS FALSE, and the file
% refutes it two hundred lines later: eq:crossshift is denominated in
% ||r||_infty and ||D||_{0,h}; lem:secondord(iv)'s SECOND bound was denominated
% in ||r||_infty (round-8 restates it in ||r||_{L^2(P)}, which A4(iii) does
% price -- Finding 5); and proofs.tex says in italics "Nothing in the assumption
% block prices ||D||_{0,h} or ||D||_{1,h}".  THE CORRECT, SCOPED CLAIM IS THIS:
%
%   In the D-DIRECTION of clause (iv) -- and there only -- lem:secondord(iv) is
%   printed in the norm A5(iv) prices, via eq:prodbound's FIRST inequality, with
%   ||D||_{1,h} shown as a stated weakening.
%
% The block makes NO block-wide claim that every bound in it is priced.  Three
% seminorms in Part F are deliberately unpriced -- ||D||_{0,h}, ||D||_{1,h} and,
% before round 8, ||r-hat - r||_infty -- and each is either a MEMBERSHIP
% condition on the direction class (the first two) or was a MIS-DENOMINATION
% since repaired (the third).
%
% The forward obligation -- a VERIFICATION, not a hypothesis -- sits in
% rem:dirmember, alongside A2'(i.b)/L-21 and N1' as the third of its kind; all
% three are inventoried in ONE PLACE, rem:obligations in results.tex, which is
% where a reader should be sent and where A5(iv)'s own presupposition (Finding 2)
% is now visible from the same page that says the obligation is open.

% A5 (iii) IS A FRONTIER, NOT A THRESHOLD.  The condition r > d_Z/2, which the
% memo headlines in four places, is the mu -> infinity LIMIT of
% \eqref{eq:frontier} and is NOT the operative condition; the correct
% derivation is the memo's own Sec.3.3 Step 3.  DO NOT WRITE r > d_Z/2.
%
% Smoothness in varphi_dagger and smoothness in varrho SUBSTITUTE along
% \eqref{eq:frontier}; the substitution is real and is worth displaying,
% because hiding it makes the block look weaker than it is:
%
%      mu       d_Z = 3           d_Z = 4           d_Z = 5
%      -----    ---------------   ---------------   ---------------
%      0.75     r > 7.5           r > 10            r > 12.5
%      1        r > 4.5           r > 6             r > 7.5
%      1.5      r > 3             r > 4             r > 5
%      2        r > 5/2           r > 10/3          r > 25/6
%      3        r > 2.1           r > 2.8           r > 3.5
%      infinity r > 1.5           r > 2             r > 2.5   (unattainable)
%
% THE mu = 2 ROW IS PRINTED AS EXACT FRACTIONS, DELIBERATELY.  Round 2 wrote
% "r > 3.33" against a true threshold of 10/3 = 3.333..., ROUNDING A STRICT
% LOWER BOUND DOWN: every r in (3.33, 10/3] passes the printed test and fails
% the frontier.  ("4.17" rounded 25/6 = 4.1666... UP and was safe, but the two
% must not be printed to different conventions.)  All other entries in the
% table are exact.
% Read the table by ROWS, not by its last line.  At mu = 1, d_Z = 3 the
% requirement is r > 4.5; at d_Z = 5 it is r > 7.5.  A binary propensity that
% is 7.5 times differentiable in five dimensions is not a condition a referee
% will grant, and the honest response is that the frontier is where the
% paper's feasibility lives: BUY mu, NOT r.  Under H1, mu = s/(2a), so buying
% mu means assuming varphi_dagger smoother or the problem less ill-posed.
%
% mu > 1/2 is a NEW NECESSARY CONDITION, absent from the memo's list.  Under
% H1 it is implied by plan Sec.2.6's own scope s > a + 1/2, with room:
% mu > 1/2 + 1/(4a).
%
% A WARNING FOR THE SIMULATION GRID (memo Sec.7, SG-20).  "A propensity
% estimator of order 2 clears n^{-1/4}" is a statement about delta-HAT_n, not
% about A5(iii).  A5(ii) constrains
% delta_n = (delta-hat_n ln n)^{2 mu / (2 mu + 1)}   [eq:deltalog, dispatch 10],
% which is LARGER -- larger by the power 2 mu/(2 mu + 1) < 1 of a quantity below
% one AND by the logarithmic factor.  SWEPT IN DISPATCH 11: this comment read the
% log-free form, which is the same staleness rem:phirscope(2) carried.
% At r = 2 the frontier reads
%     (2 mu)/(2 mu + 1) * 2/(4 + d_Z) > 1/4,
% i.e. mu > 5/6 at d_Z = 1 but mu > 7/2 at d_Z = 3.  This is consistent with
% the table above, whose mu = 3 row already demands r > 2.1 at d_Z = 3.
% QUOTING THE UNDEFLATED VERSION REPEATS, IN A NEW PLACE, EXACTLY THE MISTAKE
% THE r > d_Z/2 PROHIBITION EXISTS TO PREVENT.  SG-20's framing needs the
% deflation; flagged for whoever owns the grid.

\begin{assumption}[A8: Sieve window]\label{ass:sieve}
The balancing sieve $b_1,\dots,b_{L_x}$ is a \emph{B-spline} basis of \emph{degree exceeding
$s-\tfrac12$} on a quasi-uniform mesh over $\mathcal X$ \emph{with simple interior knots}, with
\[
n^{1/(2s)}\;\ll\;L_x\;\ll\;n^{1/2},
\]
and $\beta_1,\dots,\beta_{L_u}$ is a basis on the residual range with $L_u\le L_x$. \emph{Consequently
the sieve is contained in the Sobolev space the source conditions are written in}:
\begin{equation}\label{eq:sieveinhs}
\mathcal S_n\;=\;\operatorname{span}(b_1,\dots,b_{L_x})\cap\mathcal V\;\subset\;H^{s}(\mathcal X),
\end{equation}
so $\|\zeta\|_{H^{s}}<\infty$ at every $\zeta\in\mathcal S_n$, with $\mathcal S_n$ the sieve of
\cref{def:lfstep} and $H^{s}(\mathcal X)$ the \emph{classical} Sobolev space of H1 ---
\emph{the one that carries no boundary condition}, which is what makes the containment available to a
basis that satisfies none.

\emph{The word ``consequently'' is a derivation, and it is printed here rather than asserted}: ten live
citations rest on it. Write $p_b$ for the degree, so $p_b>s-\tfrac12$; and write $s'$ for a
\emph{generic} order of the Sobolev scale, bound in this derivation and in H1's
boundary-condition paragraph and used nowhere else.
% ARCHAEOLOGY.  Until dispatch 22 the containment carried neither a derivation nor a
% pin while the weaker statement beside it carried two pins and was called
% load-bearing.  Both letters are dispatch-22 replacements and neither changes
% anything derived: the degree was bare italic d, one glyph from the dimensions
% d_Z, d_K and from the differentials dx, du, in the FIRST file of the compiled
% reading order; the generic order was written tau at two sites and sigma at three,
% against the reserved discrepancy constant and against the four-way sigma rule.
% Glossary C87 records both.
A polynomial spline of
degree $p_b$ on a mesh with simple interior knots lies in $C^{p_b-1}(\mathcal X)$
\citep{schumaker2007}; its $p_b$th derivative is constant on each mesh cell, hence a finite step function;
and a step function on a bounded interval lies in $H^{s'}$ for every $s'<\tfrac12$, and --- unless it
is constant, i.e.\ unless the spline is a single polynomial --- in no $H^{s'}$ with $s'\ge\tfrac12$.
Hence $\mathcal S_n\subset H^{s'}(\mathcal X)$ for every $s'<p_b+\tfrac12$, and the exponent
$p_b+\tfrac12$ cannot be raised uniformly over the span. Since $p_b>s-\tfrac12$ gives $s<p_b+\tfrac12$,
this is
\eqref{eq:sieveinhs}. \emph{Three ingredients of that derivation are load-bearing and each is now stated
rather than read in.} The interior knots must be \emph{simple}: at an interior knot of multiplicity two
the spline is only $C^{p_b-2}$ there, its $(p_b-1)$st derivative already jumps, and membership drops by a
full order to $s'<p_b-\tfrac12$ --- so the containment fails at the stated degree.
% ARCHAEOLOGY.  The clause above did not say "simple interior knots" through
% dispatch 20, where the derivation was in a comment and the hypothesis it needs was
% in neither the clause nor the comment.
The space must be the \emph{classical} one, for the
reason H1 now records. And $p_b>s-\tfrac12$ is \emph{strict}, $s'=p_b+\tfrac12$ being
exactly where membership fails.
\end{assumption}

% ---------------------------------------------------------------------------
% THE BASIS CLAUSE IS STATED IN *DEGREE*, AND THE HALF-ORDER IS NOT A ROUNDING.
% Through dispatch 19 this clause read "order exceeding s".  That is half an
% order short of eq:sieveinhs AND it is ambiguous, and the two defects compound.
%
% WHY IT IS SHORT.  eq:jackson is pinned to Schumaker (2007, Thm. 6.27) -- twice,
% and proofs.tex calls that pin "load-bearing" -- so Schumaker's convention
% governs: ORDER m = DEGREE m-1, and a spline of order m on a simple-knot mesh
% lies in C^{m-2}.  Its (m-1)-st derivative is a step function, which lies in
% H^sigma exactly for sigma < 1/2, so
%       spline of order m  \in  H^sigma   <==>   sigma < m - 1/2.
% Hence S_n \subset H^s <==> m > s + 1/2, while "order exceeding s" gives only
% m > s.  The gap {s < m <= s + 1/2} is EMPTY for integer s and for {s} < 1/2,
% and NON-EMPTY otherwise -- and proofs.tex's scope note on eq:jackson states
% outright that "s need not be an integer".
%
% THE WITNESS, AND IT IS ADMISSIBLE.  Take s = 1.7, which clears eq:lxwindow's
% s > 27/16 = 1.6875 at alpha = 3/8, the block's own boundary configuration.
% The old clause admitted m = 2 -- LINEAR splines, C^0 -- and m - 1/2 = 1.5 < 1.7,
% so S_n \not\subset H^{1.7}.  (A second witness, found independently: s = 2.6,
% m = 3.)  The consequence is not a weaker bound but an EMPTY EVENT: every iterate
% lies in S_n by def:lfstep(b),(d) and (P-a.1), and varphi_dagger \in H^s by H1, so
% ||varphi-hat^{(k)} - varphi_dagger||_{H^s} = +infinity at EVERY realisation and
% four printed conditions -- A5(vi), A5(vi'), (P-a.1) and eq:foldrates' second
% display -- held on the empty set rather than with probability tending to one.
% (P-a.1) was the sharpest case: "varphi-hat^{(0)} \in S_n AND
% P(||varphi_dagger - varphi-hat^{(0)}||_{H^s} <= R) -> 1" is two halves of ONE
% clause, jointly unsatisfiable at the witness.
%
% WHY THE REPAIR IS STATED IN DEGREE, AND WHY THAT IS CONVENTION-FREE.  The fix
% deliberately does NOT adjudicate whether the retired "order" meant Schumaker
% order or degree.  DEGREE is the quantity on which the two readings agree:
%       S_n \subset H^s   <==>   degree > s - 1/2,
% which is m > s + 1/2 in Schumaker's order and the SAME inequality on the SAME
% quantity under an order-means-degree reading.  So the amended clause is correct
% under either convention and the ambiguity cannot be re-opened.
%
% WHAT THIS IS *NOT*.  It is NOT a counterexample to prp:phirate or thm:rootn.
% At the witness configuration those results' hypotheses are unsatisfiable, so
% both hold VACUOUSLY; nothing proved in this corpus is false at s = 1.7.  What
% the witness exhibits is an admissible configuration at which a MAINTAINED
% assumption's event is empty -- a defect of the assumption block, not of a
% theorem.  Recorded here so that no later reader converts it into the stronger
% claim.
%
% eq:jackson IS NOW RESTATED IN DEGREE TOO, AND THAT REVERSES THE DISPATCH-20
% DECISION RECORDED IN THIS PARAGRAPH.  The paragraph read: "eq:jackson IS NOT
% TOUCHED ... eq:jackson's own lead-in in proofs.tex still reads 'order exceeding
% s'; that is eq:jackson's HYPOTHESIS and not a quotation of this clause, and it
% is left standing deliberately."  THE PREMISE IS RIGHT AND THE DECISION WAS
% WRONG IN CONSEQUENCE.  Before the amendment the two strings were IDENTICAL, so
% this clause implied eq:jackson's hypothesis trivially under either convention.
% Amending one and not the other put them in DIFFERENT UNITS, and the implication
% then held only under the Schumaker reading -- at s = 2.4 this clause admits
% degree 2 while "order exceeding s" read as a degree needs degree 3.  Two proof
% steps consume it and NEITHER converts (lem:balance Step 4's term (III);
% prp:phirate Step 1, which does not even cite eq:jackson by label), and the
% conversion was printed exactly once, at def:lfstep bullet 4, which does not
% consume it -- and in the looser non-strict form s <= d+1 against eq:jackson's
% strict d+1 > s, so the two printed statements of one hypothesis disagreed at
% s = d+1.  eq:jackson's lead-in now reads "degree at least s-1", the non-strict
% form, which is what the arithmetic below actually gives and what def:lfstep
% bullet 4 already printed; the two statements now agree at s = d+1 as well.
% THE MATHEMATICS WAS FINE THROUGHOUT; THE UNITS WERE NOT.  The arithmetic:
%   APPROXIMATION needs the spline space to have approximation order at least s.
%     A degree-d space reproduces polynomials of degree d, so its approximation
%     order is d + 1 (= Schumaker order m), and Thm. 6.27's rate L_x^{-sigma}
%     is available for sigma <= m.  So eq:jackson's L_x^{-s} needs
%         m >= s,  i.e.  degree >= s - 1.
%   MEMBERSHIP needs  degree > s - 1/2.
% MEMBERSHIP IS THE STRICTLY STRONGER OF THE TWO, by half an order.  So this
% clause is now driven by eq:sieveinhs and eq:jackson rides free with room:
% degree > s - 1/2 > s - 1.  The amendment can only give eq:jackson slack, and
% eq:jackson's exponent L_x^{-s} is consistent with the reading adopted here --
% AND with the other reading, since order-means-degree makes the clause stronger
% still.  NOTHING TO FIX; nothing was changed.
% THE ONE RESIDUAL CHECK, NOW PERFORMED AND RESOLVED.  The above turns on
% Thm. 6.27's smoothness cap being the ORDER m and not the degree m-1.  It is the
% order: the direct theorem runs through the m-th-order modulus of smoothness, so
% the cap is m = d + 1.  Hence APPROXIMATION needs degree >= s - 1 and MEMBERSHIP
% needs degree > s - 1/2, and MEMBERSHIP IS STRICTLY THE STRONGER, by half an
% order.  Had the cap been the degree, eq:jackson would have needed degree >= s
% and would have been the binding clause instead of membership -- it is not.  The
% resolution is printed at eq:jackson in proofs.tex, where the hypothesis is now
% stated in degree, so that no site has to convert between conventions.
%
% WHERE eq:sieveinhs IS CITED.  A5(vi) and A5(vi') above; (P-a.1) and
% eq:foldrates' second display in results.tex; and def:lfstep's fourth bullet,
% whose flag ("the proviso is A8's business") carried the whole burden alone and
% pointed at eq:jackson -- WRONG IN DIRECTION, since eq:jackson says how well
% splines APPROXIMATE H^s functions and what is needed is that splines LIE IN
% H^s.  Approximation power does not give membership.  The pointer is redirected
% to this clause, which is what def:lfstep's own first sentence already said.
% ---------------------------------------------------------------------------

% THE WINDOW IS NON-EMPTY IF AND ONLY IF s > 1, AND THE BASIS IS PART OF THE
% CONDITION.  The lower end comes from the approximation term: the first-order
% bias is O_P(n^{-1/2} delta_n) + O_P(||nu_g-hat|| L_x^{-s}), and
% L_x^{-s} = o(n^{-1/2}) requires L_x >> n^{1/(2s)}.  The upper end is the
% CONDITIONING constraint on the sieve solve and is BASIS-DEPENDENT:
%     splines       L_x << n^{1/2}   ==>  window non-empty iff s > 1;
%     polynomials   L_x << n^{1/3}   ==>  window non-empty iff s > 3/2.
% The memo prices the window without naming the basis; the basis is what
% decides whether the smoothness cost is s > 1 or s > 3/2, and it competes
% with A2, H1 and A5(iii) for the same s budget.  Splines are chosen.
%
% ROUND-7: THE TWO ENDS ARE ROUTE-DEPENDENT, AND THE BLOCK NOW SAYS SO.  N-18,
% delivered at Step 5 and confirmed by proofs.tex:lem:balance, is that A5(iv) and
% A8 are ALTERNATIVE sufficient conditions for the first-order bias, NOT
% cumulative ones:
%     RATE ROUTE      lem:balance(ii): A5(ii) + A5(iv) alone.  NO L_x condition,
%                     NO balancing constraint, NO eq:nuhat, NO A5(vi).  In
%                     particular A8's LOWER end -- and with it the "s > 1 /
%                     s > 3/2" line of the smoothness-budget table -- IS NOT USED.
%     BALANCING ROUTE lem:balance(iii): A8 in full + eq:nuhat in full + A5(vi),
%                     and NO rate for ell-hat.
% A8's UPPER end L_x << n^{1/2} survives on both routes: it is the conditioning
% constraint for the sieve solve and L-13 consumes it.
%
% DISPATCH-8: prp:phirate NOW CITES A8 BY CLAUSE, AND IT HAD TO.  Through
% round 10 prp:phirate listed "A8" in full among its hypotheses while its own
% Step 1 wrote "by A8's lower end and the spline direct theorem" -- two errors
% in one line.  The direct theorem needs the BASIS clause (B-splines of DEGREE
% exceeding s - 1/2 on a quasi-uniform mesh with simple interior knots), not the
% lower end.  DISPATCH-21 DESTALING: this parenthesis read "of order exceeding s"
% -- the pre-amendment wording of the clause, in the PRESENT TENSE, 120 lines
% below the amended clause and inside the same file.  The dispatch-20 sweep
% corrected four such quotations elsewhere and missed this one at home, which is
% the register defect the B1(vi) header already names: a clause amended in one
% place and quoted unamended in another is worse than one nobody amended.  The
% direct theorem's own hypothesis is DEGREE AT LEAST s - 1 and is implied by the
% amended clause with half an order to spare, so nothing in the sentence's
% CONCLUSION changes; only its quotation of the hypothesis does.  And A8's lower end
% together with (P-d) requires s > 1 / (2(1 - 2 alpha)), which at the boundary
% configuration alpha = 3/8 is s > 2 against the s > 1.75 the block recorded --
% AN EMPTY WINDOW.  prp:phirate now takes A8's BASIS and UPPER END and NOT its
% lower end, exactly as this comment already said the rate route should; (P-c)
% is the rate route's own lower end for L_x, calibrated to h_u delta-hat_n
% rather than to n^{-1/2}; and the joint line (P-c) ^ (P-d), eq:lxwindow, is now
% a row of the smoothness-budget table, as is the stronger both-routes line.
% THE FIX IS C2'(i)'s FROM DISPATCH 6, APPLIED TO A8: cite the clause, not the
% assumption.
% CONSEQUENCE FOR THE BUDGET TABLE: the "Sieve window non-empty" row is priced
% ROUTE-CONDITIONALLY there, since s is already contested against A2 (via
% mu = s/2a), H1 and A5(iii), and charging s > 1 unconditionally over-charges the
% one budget the paper's feasibility depends on.
%
% Conservative but honest, and still true on the balancing route: the third term
% is really bounded by ||nu_g-hat|| min(delta_n, L_x^{-s}), so even there the
% s > 1 requirement may relax further.  That refinement is not claimed.

% ---------------------------------------------------------------------------
% PART V.  REGULARISATION OF THE NONLINEAR OPERATOR EQUATION
%
% DISPATCH-7 RESTRUCTURE.  Through round 9 this part fixed the reading of
% A-hat (def:hatA) and then stated two maintained assumptions, A6 and A7,
% ABOUT AN OBJECT -- T-hat -- THAT NO FILE DEFINED.  Part H exhibited it;
% def:hatT below now DEFINES it here, where A7 needs it, so that the maintained
% block no longer names an undefined map.  The integrated kernel def:intker,
% which eq:hatT is written in, moves here with it; results.tex keeps
% lem:intker, which is where its four properties are derived from C3.
% A6 IS RETIRED -- see the retired-assumption block below and prp:tccfails.
% ---------------------------------------------------------------------------

\begin{definition}[The integrated kernel]\label{def:intker}
Let $C_{h_u}$ be the kernel of C3 and, for $s_0\in[0,1]$, define
\begin{equation}\label{eq:intker}
\mathsf C_{h_u}(v;s_0)\;:=\;\int_{-\infty}^{v/h_u}\Theta_{s_0}(\varsigma)\,d\varsigma
\;=\;\int_{-\infty}^{v/h_u}C_{h_u}\big(h_u\varsigma,s_0\big)\,d\varsigma
\;=\;\frac1{h_u}\int_{-\infty}^{v}C_{h_u}(\varsigma,s_0)\,d\varsigma ,
\qquad v\in\IR ,
\end{equation}
with $\Theta_{s_0}$ of \eqref{eq:rescaled}, and write $\mathsf C_{h_u}(v):=\mathsf C_{h_u}(v;1)$. The same definition with $h_x$ in place of $h_u$
defines $\mathsf C_{h_x}$. \emph{The boundary index is carried only in the $x$-direction}: in the
residual direction the argument $u$ ranges over all of $\IR$ and there is no boundary, so
\eqref{eq:hatT} uses $\mathsf C_{h_u}(\cdot)=\mathsf C_{h_u}(\cdot;1)$ throughout; in the
$x$-direction $\mathcal X$ is a compact interval by A0 and the index is used.
Its properties are \cref{lem:intker}.
\end{definition}

\begin{definition}[The sample operator $\hat T$]\label{def:hatT}
Adopt the sample-splitting convention of \cref{lem:balance}: $\{1,\dots,n\}=\mathcal I_a\cup\mathcal I_m$,
$n_{\mathrm a}:=\#\mathcal I_a$, $n_{\mathrm m}:=\#\mathcal I_m$, the two folds independent by
A0, with $\hat\varrho$ computed on $\mathcal I_a$ as A4(iii) requires.
\emph{The split is \textbf{balanced}}: $n_{\mathrm a}$ and $n_{\mathrm m}$ are non-random with
\begin{equation}\label{eq:foldbal}
n_{\mathrm m}\;\asymp\;n
\qquad\text{(equivalently, }\liminf_n n_{\mathrm m}/n>0\text{)} .
\end{equation}
% EXTRACTED FROM THE STATEMENT OF def:hatT ON 2026-08-03, under the house rule
% (.claude/references/theorem-environment-style.md) that a definition carries
% CONSTRUCTION ONLY.  WHAT eq:foldbal COSTS, and the three sites that consume it:
% rem:hatTread(1).  eq:foldbal ITSELF STAYS, with its \label.
% ARCHAEOLOGY.  Through round 11 eq:foldbal was assumed nowhere.

\emph{Construction convention, and it costs nothing}: $\hat\varrho$ is taken to have values in
$[\varepsilon_\varrho,1-\varepsilon_\varrho]$, with $\varepsilon_\varrho$ the overlap constant of
A4(i); if a first-stage estimator does not, replace it by its truncation to that
interval. Writing $\widehat{\widetilde W}_i:=W_i-\hat\varrho(\Delta Z_i)$, the \emph{sample operator} is
% WHAT THE TRUNCATION BUYS -- (a) it can only decrease the L^2 distance, so A4(iii)
% is preserved; (b) the resulting two-sided bound on |hatWtilde_i| -- and the scope
% note for a smaller known lower bound epsilon: rem:hatTread(2) and (3).
% THE DEFINITION OF hatWtilde_i DID NOT MOVE, AND MUST NOT.  Clause (b) opened by
% DEFINING the symbol and then derived a bound on it; eq:hatT immediately below
% USES the symbol, so cutting "(a) and (b)" wholesale would have left it
% used-but-undefined.  The definitional half stays here; only the derived bound
% epsilon_varrho <= |hatWtilde_i| <= 1 moved.
\begin{equation}\label{eq:hatT}
\hat T(\varphi,\theta)(u,t)\;:=\;\frac1{n_{\mathrm m}}\sum_{i\in\mathcal I_m}
\widehat{\widetilde W}_i\;\mathsf C_{h_u}\big(u-\Delta U_i(\theta,\varphi)\big)\;w\big(t'\Phi(\Delta Z_i)\big),
\qquad(u,t)\in\IR\times\mathcal T ,
\end{equation}
with $\Delta U_i(\theta,\varphi)$ the trial residual of \eqref{eq:trialres}. Its $x$-smoothed variant
$\hat T_{h_x}$, which is what \cref{def:lfstep} iterates on and what
$\hat{\mathcal A}_\varphi$ differentiates, is \eqref{eq:hatTx}.

% MIRROR DRIFT, PRE-EXISTING AND DELIBERATELY NOT CONVERGED (2026-08-03).  The
% paragraph below is commentary of exactly the class this pass extracts, and the
% SPLIT MIRROR ALREADY CARRIES IT OUTSIDE THE STATEMENT, at assumptions_supp.tex
% under the paragraph heading "the sample operator hatT".  It is left here rather
% than folded into rem:hatTread because assumptions_supp.tex already holds this
% text, so adding it to the parent's remark and not to the split's would give
% rem:hatTread a different clause list on each side -- and every cross-document
% citation of the form \cref{rem:hatTread}(N) would then resolve to two different
% clauses.  Recorded for whoever converges the two files.
\emph{Why the truncation is in the maintained block and not in a proof.} Without it
$\hat\varrho\in\{0,1\}$ has positive probability, hence so does $\widehat{\widetilde W}_i=0$; on the
event that this holds at every $i\in\mathcal I_m$ the sample operator is identically zero and every
statement about it is vacuous. A4(i) constrains $\varrho$ and says nothing about
$\hat\varrho$, so the degeneracy cannot be excluded by assumption --- it is excluded by construction,
which is where it belongs.
\end{definition}

\noindent\emph{Notes on \cref{def:hatT}.} What the balanced split \eqref{eq:foldbal} costs and where it is
consumed, what the truncation buys, the scope note for a smaller known lower bound, and which assumption
names \eqref{eq:hatT} are \cref{rem:hatTread}.

\begin{definition}[The smoothed sample linearisation $\hat{\mathcal A}$]\label{def:hatA}
For $\varphi\in\mathcal V$ let $\mathcal A_\varphi:\mathcal V\to L^{2}(\Pi)$ be the Fr\'echet derivative of
$T(\cdot,\theta)$ at $\varphi$, so that $\mathcal A_{\mathrm{cond}}=\mathcal A_{\vdagger}$, with kernel
$\mathsf k(u,t;x)=\nu_{g_{u,t}}(x)/\pi(x)$ as in \cref{prp:compact}. Write $\hat{\mathcal A}_\varphi$ for
the \emph{kernel-smoothed sample analogue} of $\mathcal A_\varphi$: the integral operator whose kernel is
obtained from $\mathsf k$ by replacing the indicator $\un\{\Delta U(\theta,\varphi)\le u\}$ inside
$g_{u,t}$ by its $C_{h_u}$-smoothed version, and each of the two products
$E[\,\cdot\mid X_j=x]f_{X_j}(x)$ constituting $\nu_{g_{u,t}}$ by its $C_{h_x}$ kernel estimate. Write
$\hat{\mathcal A}^\ast_\varphi$ for its adjoint in the $L^{2}(\pi)$--$L^{2}(\Pi)$ pairing, and
\[
\hat{\mathcal A}_{\mathrm{cond}}:=\hat{\mathcal A}_{\vdagger},
\qquad
\hat{\mathcal A}^\ast_{\mathrm{cond}}:=\hat{\mathcal A}^\ast_{\vdagger},
\]
the running Landweber--Fridman iterate being written $\hat{\mathcal A}_{\hat\varphi_m}$ wherever the base
point matters.

\emph{This fixes a reading, and the reading is substantive rather than notational.}
$\hat{\mathcal A}_\varphi$ is a \textbf{separately smoothed operator} and \emph{not} the Fr\'echet
derivative of $\hat T$ in $\varphi$; the hat marks a sample analogue of $\mathcal A_\varphi$, not a
derivative of $\hat T$. \emph{The bandwidth is what forces the choice}: C3 attaches
$C_{h_x}$ to $\hat{\mathcal A}_{\mathrm{cond}}$ and $\hat{\mathcal A}^\ast_{\mathrm{cond}}$, and a
derivative of $\hat T$ would carry no bandwidth of its own. \emph{And the rejected reading is now known
to be unavailable rather than merely different}: by \cref{prp:hatAprop}(v), $\hat T(\cdot,\theta)$ is
nowhere Fr\'echet differentiable on $\mathcal V$, so ``the derivative of $\hat T$'' names no operator.

\emph{Forward obligation, DISCHARGED at \cref{def:hatAcon} and \cref{prp:hatAprop}, with one negative
answer.} This definition names a construction; it does not perform one. Step~6 was to exhibit
$\hat{\mathcal A}_\varphi$ explicitly --- in particular the smoothed indicator and the two kernel
estimates --- and to verify the two conditions then maintained for the object exhibited.
\eqref{eq:hatAcon} exhibits it; A7 is verified in the amended form
\eqref{eq:a7amend} by \cref{prp:stopfeas}; and \emph{the tangential cone condition A6 is not verified
but refuted} (\cref{prp:tccfails}), which is why the block below carries it as retired.
Recorded in \cref{rem:obligations} as obligation (O4).
\end{definition}

\begin{assumption}[A6: RETIRED --- tangential cone condition for the \emph{random} operator]\label{ass:tcc}
\textbf{\sout{RETIRED.}} \emph{This assumption is not maintained, it is not available for use, and no
result in this programme takes it as a hypothesis.} It read: \emph{there exist
$\eta_{\mathrm{TCC}}<\tfrac12$, not depending on $n$, and a ball
% ARCHAEOLOGY.  It was maintained through round 9.
$\mathcal V_n\subset\mathcal V$ around $\vdagger$ containing the Landweber--Fridman path, such that, with
$\hat{\mathcal A}_{\tilde\varphi}$ the smoothed sample linearisation of \cref{def:hatA},}
\[
P\Big(\big\|\hat T(\varphi,\theta)-\hat T(\tilde\varphi,\theta)-\hat{\mathcal A}_{\tilde\varphi}(\varphi-\tilde\varphi)\big\|_{L^2(\Pi)}
\;\le\;\eta_{\mathrm{TCC}}\,\big\|\hat T(\varphi,\theta)-\hat T(\tilde\varphi,\theta)\big\|_{L^2(\Pi)}
\ \ \forall\,\varphi,\tilde\varphi\in\mathcal V_n\Big)\;\longrightarrow\;1 .
\qquad\text{\sout{[retired]}}
\]
\emph{Why it is retired, and it is not a matter of strength.} \Cref{prp:tccfails} proves that the
displayed event has probability zero --- at every $n$, at every $\eta_{\mathrm{TCC}}<1$, at every ball
of positive $L^{2}(\pi)$-radius, and \textbf{for every bounded operator $\mathcal V\to L^{2}(\Pi)$ put
in place of $\hat{\mathcal A}_{\tilde\varphi}$}. An assumption whose event is null cannot be maintained
and cannot be repaired by lowering a constant or shrinking a ball.
\emph{What it should have said is \cref{def:a6b} (A6b)}, which pairs $\hat{\mathcal A}_\varphi$ with the
operator it actually differentiates, $\hat T_{h_x}$; A6b is a \emph{property}, on the A1d/A2b/A2c model,
never asserted and consumed by no theorem. \emph{Scope of the refutation, stated because it is
narrower than the retirement}: \cref{prp:tccfails} is a statement about $L^{2}(\pi)$-balls in
$\mathcal V$. On the spline sieve $\mathcal S_n$ the spike it uses does not exist, so a cone condition
\emph{restricted to $\mathcal V_n\cap\mathcal S_n$} is untouched by it --- and would still not deliver
the inherited route, for the independent reason recorded at \cref{rem:tccwhy}\textup{(3)}.
\emph{Nothing downstream is affected}: \cref{prp:phirate} uses neither A6 nor A6b.
\end{assumption}

\begin{assumption}[A7: Stopping feasibility, at the amended threshold]\label{ass:stopfeas}
The level used by the discrepancy principle dominates the true operator noise level with probability
approaching one: for some $\tau_{\mathrm{DP}}\ge1$,
\begin{equation}\label{eq:a7amend}
P\Big(\big\|\hat T(\vdagger,\theta)\big\|_{L^2(\Pi)}\le\tau_{\mathrm{DP}}\,\hat\delta_n\,\ln n\Big)
\;\longrightarrow\;1 ,
\end{equation}
with $\hat T$ as in \eqref{eq:hatT}. \emph{Any divergent factor serves in place of $\ln n$}; $\ln n$ is
fixed here so that one symbol, not a family, propagates downstream.
\end{assumption}

% A7'S THRESHOLD IS AMENDED IN THIS ROUND, AND THE AMENDMENT IS MANDATORY RATHER
% THAN OPTIONAL.  Rounds 1-9 wrote the threshold tau_DP * delta-hat_n.  Two
% facts force the change.
%  (a) IT IS NOT DERIVABLE AT A FIXED CONSTANT.  prp:stopfeas(i) closes the
%      T-component at a fixed tau_DP; prp:stopfeas(ii) shows the varrho
%      component does NOT, because A4(iii) is an O_P statement and an O_P
%      hypothesis carries no fixed-constant exceedance bound.  The alternative
%      repair -- strengthen A4(iii) to P(||varrho-hat - varrho|| > c r_n) -> 0 --
%      is also sufficient (prp:stopfeas(iii)) and is a strictly stronger
%      hypothesis on the propensity estimator, so the threshold moves instead.
%  (b) def:lfstep ALREADY IMPLEMENTS THE AMENDED LEVEL.  Leaving A7 as written
%      would put the maintained block in contradiction with the estimator.
%
% WHAT THE AMENDMENT COSTS, AND IT IS NOT NOTHING.  THE SUBSTITUTION IS AT THE
% LEVEL OF delta-hat_n, NOT OF delta_n -- SWEPT IN DISPATCH 11, because the old
% wording ("delta_n -> delta_n (ln n)^{2 mu/(2 mu+1)}") reads as a SECOND
% inflation once A5(ii) carries the logarithm in the symbol, which is exactly the
% one-symbol-one-object collision eq:deltalog was written to end.
% prp:stopfeas(iv) shows the substitution delta-hat_n -> delta-hat_n ln n, whose
% effect on the structural rate IS eq:deltalog's
% delta_n = (delta-hat_n ln n)^{2 mu / (2 mu + 1)}, is absorbed by the
% STRICTNESS of eq:frontier, so A5(ii)'s o_P(n^{-1/4}) survives.  IT IS NOT FREE
% AGAINST (P-f).  (P-f) is an EXACT domination, delta_n^2 = O(h_u^2 delta-hat_n
% ln n), and both sides now carry logarithms: the left carries (ln n)^{4 mu /
% (2 mu + 1)} and the right (ln n)^1, with 4 mu / (2 mu + 1) > 1 exactly when
% mu > 1/2 -- i.e. always, in this block.  The consequence is that (P-f)'s
% thresholds are STRICT: alpha (2 mu - 1)/(2 mu + 1) > 1/8 at C5'(i)'s bandwidth,
% and mu > 5/6 at the re-tuned bandwidth of C5'(i) as amended.  At the reference
% configuration (mu, d_Z, r) = (1, 3, 4.5), which sits exactly at the boundary in
% powers of n, (P-f) FAILS by (ln n)^{1/3}.  See rem:h6bandwidth.

% THE DISCREPANCY CONSTANT IS tau_DP, NOT tau -- ROUND-7 RENAME, AND IT IS THE
% CHEAP SIDE OF A COLLISION.  Part F uses tau = (tau_1, tau_2, tau_3) for the
% perturbation parameter of lem:secondord and prp:orthog, and tau-bar for its
% radius; A7's constant is a different object.  The perturbation parameter is
% used at roughly thirty sites and the discrepancy constant at two, so the
% constant moves.  The word subscript follows the eta_TCC precedent, which the
% glossary adopted for exactly this reason: tau_DP never appears without it.

% A7 IS STATED FOR THE RANDOM OPERATOR, and that is the point.
% hanke1995 proves the discrepancy principle for an EXACTLY KNOWN operator with
% NOISY DATA.  Here the noise is IN THE OPERATOR: T-hat is built from smoothed
% indicators and from varrho-hat, so T-hat(varphi_dagger) != 0.  This is the
% Dunker, Florens, Hohage, Johannes & Mammen (2014) noise model, which the
% project ADOPTS AS ITS SETTING even though it rejects IRGN as the regulariser.
%
% Without A7 THE DISCREPANCY PRINCIPLE NEVER STOPS: if the threshold sits below
% the noise floor, the residual never falls under it and the estimator is
% whatever the iteration cap bar-m_n returns, with no rate.  THAT IS EXACTLY THE
% FAILURE THE ROUND-9 THRESHOLD INVITED, and it is why def:lfstep(c) is
% STAGE-ADAPTIVE: the data error of the stage-k LINEARISED equation is not
% delta-hat_n but delta-hat_n ln n + h_u^{-2} eps_{k-1}^2, and a fixed threshold
% sits below it in exactly the regime the recursion must pass through.  A7
% supplies the FIRST of those two terms; the second is supplied by the
% deterministic schedule eq:lfsched, which the estimator computes.
% A7 is testable in simulation (probe CP-13): report the realised stopping index
% BY STAGE and how often the rule fails to stop before the cap, running the
% fixed and the stage-adaptive thresholds side by side.
%
% NOTE what prp:frechet(ii) supplies and what it does not.  It gives
% ||A_{varphi+h} - A_varphi|| <= C ||h||, i.e. the DETERMINISTIC operator's
% derivative map is Lipschitz.  A cone condition on the RANDOM operator T-hat
% does not follow from it -- which is now moot for A6 (retired) and is retained
% because it is the same distinction def:a6b turns on.

% ---------------------------------------------------------------------------
% PART VI.  REGULARITY CONDITIONS CARRIED OVER FROM THE DRAFT APPENDIX
%           (assapp1--assapp5, restated for the panel and for the new operator)
% ---------------------------------------------------------------------------

\begin{assumption}[C1: Densities]\label{ass:app1}~
Write $x=(x_1,x_2)\in\mathcal X^{2}$.
\begin{enumerate}
\item[(i)] The conditional density $f_{\Delta Y,X_1,X_2\mid\Delta Z,W}$ and the joint density
$f_{\Delta Y,X_1,X_2\mid\Delta Z}$ are $p_K$ times continuously differentiable and uniformly bounded
above.
\item[(ii)] $f_{X_1,X_2}$, $f_{X_1}$, $f_{X_2}$ and $\pi$ are uniformly bounded above and below away
from $0$ on $\mathcal X^{2}$ and $\mathcal X$ respectively; consequently
$\bar c_\pi:=\max_{j\in\{1,2\}}\|f_{X_j}/\pi\|_\infty<\infty$.
\end{enumerate}
\end{assumption}

% WHICH CLAUSE IS CONSUMED WHERE -- ROUND-7 NARROWING.  Every result proved
% through Step 5 that cites C1 consumes ONLY C1(ii): prp:frechet, prp:compact,
% lem:nusigma, lem:adjfubini, lem:adjoint, prp:kerchar, cor:a1cprim,
% cor:rangegap, prp:nophifree, lem:solset, lem:secondord, prp:orthog and
% lem:balance all use it for bar-c_pi, for ||1/pi||_inf, for ||pi||_inf, or for
% strict positivity of pi.  Those citations are now written C1(ii).
% C1(i): WHERE IT IS ACTUALLY CONSUMED.  THE ROUND-7 FORM OF THIS NOTE WAS
% FALSE AND IS STRUCK.  It read: "C1(i) HAS NO CONSUMER BEFORE STEP 6 ... it
% is consumed by prp:phirate (Step 6) together with C3 and C5'.  It is retained,
% and this note is the forward reference that keeps the usage matrix honest."
% THE FORWARD REFERENCE NEVER LANDED.  prp:phirate's proof cites C1 at no line
% at all, neither clause, and neither lem:unifop nor lem:opnoise carries C1 on
% its hypothesis list.  A note written to keep the usage matrix honest was the
% one place the matrix was wrong, which is the defect species this block keeps
% producing and the reason the correction is printed rather than deleted.
%
% THE TRUE STATEMENT.  C1(i) is cited at exactly two lines in the corpus,
% proofs.tex Step 2 of lem:aerr, where the p_K derivatives of the products
% E[breve-g | X_j = .] f_{X_j} are what make the x-smoothing bias O(h_x^{p_K})
% at order p_K rather than at floor(s - 1/2).  That is its only consumer.
% AND THAT CONSUMER IS ITSELF CONSUMED BY NOTHING: lem:aerr opens "Consumers:
% none".  So C1(i) is a hypothesis of one lemma and, transitively, of no
% theorem.  IT IS NOT RETIRED, because lem:aerr's proof needs it and lem:aerr is
% retained for the two reasons its own reading (iv) gives -- it is the
% re-derivation C5'(ii) asks for, and it is the price-tag on the route Step 0 of
% prp:phirate DECLINES.  A clause whose only consumer is live is not in the
% ass:rank / ass:tcc / A2'(ii) / B1(vi) class and must not be filed there.
% The printed statement of all this is in the Discussion below, under the
% Part VI clause-status list, on the A5(v) model.
%
% CONSEQUENCE FOR prp:phirate'S OWN LIST, closed in the same pass: that list
% read "C1" in full and now reads C1(ii), which it consumes transitively
% through prp:hatAprop and lem:nusigma.  It does not consume C1(i), and the
% struck note above was the reason a reader would have thought it did.
%
% CHANGES FROM THE DRAFT'S assapp1:
%  (a) restated for the PANEL, with x = (x_1,x_2) -- as plan Sec.2.6b directs;
%  (b) conditioned on Delta Z as well as W, which A1 requires;
%  (c) the draft's smoothness order "d" is renamed p_K, the kernel order, both
%      to remove the collision with d_Z and because the two must match for the
%      bias order to be realised;
%  (d) the ratio bound bar-c_pi is made explicit because every boundedness
%      statement about A_cond and about the second-order form uses it, and the
%      memo's L-20 states it on the wrong object.

\begin{assumption}[C2$'$: Smoothness of the residual law]\label{ass:app2}~
\begin{enumerate}
\item[(i)] The conditional density $f_{\Delta Y\mid X_1,X_2,\Delta Z,W}(\cdot\mid\cdot)$ exists, is
$p_K-1$ times continuously differentiable in its first argument, and admits envelopes
$\bar f_k(x_1,x_2,z,\mathsf w):=\sup_{u}\big|\partial_u^{k}f_{\Delta Y\mid X_1,X_2,\Delta Z,W}(u\mid x_1,x_2,z,\mathsf w)\big|$
with $E[\bar f_k]<\infty$ for $k=0,1,\dots,p_K-1$ and $\|\bar f_0\|_\infty+\|\bar f_1\|_\infty<\infty$.

\emph{Two readings of that display are fixed here --- \textbf{no letters are attached to them}, so that
nothing below can be cited as a sub-clause of C2$'$(i) --- because without them neither $E[\bar f_k]$ nor
$\|\bar f_k\|_\infty$ is a well-defined quantity.}

\emph{First, the version, and why each $\bar f_k$ is measurable.} The density is read at a version that
is \textbf{jointly Borel} in $(u,x_1,x_2,z,\mathsf w)$ and, at \emph{every} conditioning point rather than
almost every one, $p_K-1$ times continuously differentiable in $u$ --- which is what ``continuously
differentiable in its first argument'' is read to mean above. \emph{The supplier is one Part away and is now
cited rather than left implicit}: \textup{(P6a)} of \cref{def:pregver} fixes exactly such a version for the
conditional law of $\Delta Y_\theta-\Delta\varphi$ given $\mathcal F_{XZW}=\sigma(X_1,X_2,\Delta Z,W)$, for
every $\omega$ --- at $(\theta,\varphi)=(0,0)$ that law \emph{is} the one this clause states, and every other
is a translate of it by the shift identity of \textup{(P6)} --- and the Doob--Dynkin lemma transports an
$\mathcal F_{XZW}$-measurable function of $\omega$ to
a Borel function of the conditioning \emph{value}, which is the argument $\bar f_k$ carries.
\emph{Those two supply the version \textbf{one variable at a time}, and the declaration above is a
\textbf{joint} one, so a third step is named rather than elided.} \textup{(P6a)} gives joint measurability
in $(\omega,u)$; Doob--Dynkin runs at each \emph{fixed} $u$ and returns, for that $u$, a Borel function of
the conditioning value; neither, nor the two together, is joint Borel measurability in
$(u,x_1,x_2,z,\mathsf w)$. \emph{What delivers it is the rationals-plus-continuity step this paragraph
already performs below for the envelope}: at every conditioning point the map $u\mapsto\partial^{k}_uf$ is
continuous, and a
function that is Borel in one argument at each fixed value of the other and continuous in the other at each
fixed value of the first is jointly Borel --- the Carath\'eodory criterion, reached here by taking $u$ along
the rationals and passing to the limit, which is the same device \cref{def:crossfit}\textup{(c$'$.1)}
runs for \eqref{eq:prodmeas} and the proof of \cref{prp:frechet} runs for $G_\varphi$.
\emph{Nothing that consumes this clause is affected}: every site reads $\bar f_k$ or
$\|\bar f_k\|_\infty$, and the sentences below build both from the version fixed above at each separate
conditioning point. \emph{The shortfall was in the \textbf{declaration} and not in the consumption}, and it is
repaired here because a declaration a reader cannot derive from its own stated suppliers is a gap
wherever it sits. \emph{It is also
the version the corpus already runs}: the Carath\'eodory argument opening the proof of \cref{prp:frechet}
establishes joint measurability of $G_\varphi$ and of its first two $u$-derivatives $G'_\varphi$ and
$G''_\varphi$ --- which \emph{are} $\partial_u^{k}f$ at $k=0$ and $k=1$ --- from measurability in $\omega$ and
continuity in $u$, and says there that it is used silently below. \emph{Extending it to $k\le p_K-1$ is the same argument and no more}: each $\partial_u^{k}f$ is the
pointwise limit of the difference quotients of $\partial_u^{k-1}f$ in $u$ taken along the reciprocals of the
positive integers, and a pointwise limit of countably many jointly Borel functions is jointly Borel --- the
device \cref{def:crossfit}\textup{(c$'$.1)} uses for $\partial_u\hat\ell^{(-k)}$. \emph{Then the envelope is a
countable supremum}: continuity in $u$ makes $\sup_{u\in\IR}=\sup_{u\in\IQ}$ at every conditioning point, so
each $\bar f_k$ is Borel as a map into $[0,+\infty]$, $E[\bar f_k]$ is well defined in $[0,+\infty]$, and the
display above is the requirement that it be finite. \emph{Without this the display is an expectation of a
function not known to be measurable} --- the ground on which \cref{lem:reprsmooth} Step~3 refuses the same
elision, and the same species as the defect repaired at A4(iii).

\emph{Second, the outer norm is an \textbf{essential} supremum, and the inner one is genuine.}
$\|\bar f_k\|_\infty:=\operatorname{ess\,sup}\bar f_k$ under the law of $(X_1,X_2,\Delta Z,W)$, the supremum
in $u$ inside $\bar f_k$ remaining a genuine supremum --- \emph{the same split
\cref{def:crossfit}\textup{(c$'$.2)} settles for \eqref{eq:supnormconv}, and settled the same way and for the
same reason}. \emph{The essential reading is what every consumer needs and is the weaker of the two}: in the
proof of \cref{prp:frechet} at Steps~3--5, in that of \cref{prp:compact}, in that of \cref{lem:adjfubini} and
at \cref{lem:aerr} Step~4,
$\|\bar f_0\|_\infty$ and $\|\bar f_1\|_\infty$ appear only on the right of an inequality \emph{inside} an
expectation or a conditional expectation, so an almost-sure bound is exact there and a sure bound would be
surplus strength. \emph{One thing does \textbf{not} carry over from \eqref{eq:supnormconv}}: there the outer
essential supremum needed a measurability argument of its own because the object inside it was random; here
$\bar f_k$ is a \emph{deterministic} function and $\|\bar f_k\|_\infty$ is a number, so no measurability
question arises for the outer norm and none is settled. What had to be settled was only which of the two
readings it is.
\item[(ii)] For every $t\in\mathcal T$ the map $u\mapsto\mathcal M(u,t)$ is $p_K$ times continuously
differentiable, with $p_K=8$ and
\[
\max_{0\le k\le p_K}\ \sup_{t\in\mathcal T}\ \bar w(t)^{-1}\sup_{u\in\IR}\big|\partial_u^{k}\mathcal M(u,t)\big|\;<\;\infty .
\]
\end{enumerate}
\end{assumption}

% THE DRAFT'S assapp2 IS INSUFFICIENT AND IS REPLACED.  It names only f_U.  Two
% distinct jobs need two distinct objects.
%
%  (a) L-4, the step that produces the Frechet derivative A_cond, needs the
%      density of Delta Y CONDITIONAL ON (X_1, X_2, Delta Z, W) -- i.e.
%      conditional on the ENDOGENOUS regressors -- and needs it C^1 in u with a
%      bounded derivative, one order MORE than the memo's L-4 requests.  The
%      extra order is what turns Gateaux differentiability into Frechet
%      differentiability with a QUADRATIC remainder (prp:frechet(i)) and what
%      makes the derivative map Lipschitz (prp:frechet(ii)).  Clause (i).
%
%  (b) The order-p_K kernel smoothing of the indicator has bias O(h_u^{p_K}),
%      and realising that order needs the SMOOTHED FUNCTION to have p_K
%      derivatives.  ROUND-2 CORRECTION: the smoothed function is NOT the
%      marginal f_{Delta U}.  What T-hat smooths is the CDF-level object
%      u |-> M(u,t), whose derivative is
%        partial_u M(u,t)
%          = E[Wtilde w(t'Phi(Delta Z)) f_{Delta U | X_1,X_2,Delta Z,W}(u)],
%      i.e. the Wtilde-w-WEIGHTED CONDITIONAL density; and the budget must be
%      uniform in t, since the operator norm is taken over the whole index set.
%      Round 1 put the budget on the marginal f_{Delta U}, which is neither
%      necessary nor sufficient.  Clause (ii) states it directly on M.
%      Clause (i), with envelopes to order p_K - 1, is a sufficient primitive
%      for (ii): differentiate under the expectation with dominating functions
%      bar-w(t) bar-f_k.
%
%  This also disposes of the round-1 "one derivative of slack" note, which was
%  computed on the wrong density and is withdrawn.

\begin{assumption}[C3: Kernel, and the defining property of the boundary adjustment]\label{ass:app3}
Fix $p_K=8$ and a finite $c_C>0$. \emph{Write, for a bandwidth $h>0$ and a boundary index
$s_0\in[0,1]$,}
\begin{equation}\label{eq:rescaled}
\Theta_{s_0}(\varsigma)\;:=\;C_{h}\big(h\varsigma,\,s_0\big),
\qquad\varsigma\in\IR,
\end{equation}
\emph{for the kernel read as a function of its scaled argument}. The smoothing kernel used in the residual
direction is a univariate generalized kernel $C_{h_u}$ with:
% EXTRACTED FROM THE STATEMENT OF C3 ON 2026-08-03, under the house rule
% (.claude/references/theorem-environment-style.md) that an assumption carries
% CONDITIONS ONLY.  C3 is very largely math displays, which ARE conditions, so
% only two things left it and the block is otherwise unchanged.
%   -- WHAT eq:rescaled IS FOR -- the object every bound below is stated for and
%      the object eq:intker integrates:                     rem:app3read(1).
%   -- CLAUSE (ii)'s over-long heading, which asserted two readings rather than
%      naming the clause:                                   rem:app3read(2).
% NO \label MOVED: eq:rescaled and eq:truncmom both stayed in the statement.
\begin{enumerate}
\item[(i)] \emph{Order $p_K$ at the interior index.} $\Theta_1$ is supported in $[-c_C,c_C]$ and
\[
\int_{-c_C}^{c_C}\Theta_1(\varsigma)\,d\varsigma\;=\;1,
\qquad
\int_{-c_C}^{c_C}\varsigma^{\,j}\,\Theta_1(\varsigma)\,d\varsigma\;=\;0
\quad (1\le j\le p_K-1),
\qquad
\int_{-c_C}^{c_C}|\varsigma|^{p_K}\big|\Theta_1(\varsigma)\big|\,d\varsigma\;<\;\infty .
\]
\emph{The normalisation $\int\Theta_1=1$ is printed rather than left inside the words ``of order
$p_K$''}: it is consumed at \cref{lem:intker}(i), (iii) and (iv).
% ARCHAEOLOGY.  Rounds 1--11 left it implicit.
\item[(ii)] \emph{Boundary adjustment.} For each $s_0\in[0,1]$, $\Theta_{s_0}$ is supported
in $[-c_Cs_0,\;c_C]$ and
\begin{equation}\label{eq:truncmom}
\int_{-c_Cs_0}^{c_C}\varsigma^{\,j}\,\Theta_{s_0}(\varsigma)\,d\varsigma\;=\;\delta_{j0}
\qquad (0\le j\le p_K-1),
\end{equation}
with $\delta_{j0}$ the Kronecker symbol; and, uniformly in $h>0$ and $s_0\in[0,1]$,
\[
\sup_{|\varsigma|\le c_C}\big|\Theta_{s_0}(\varsigma)\big|\;<\;\infty
\qquad\text{and}\qquad
\int_{-c_Cs_0}^{c_C}|\varsigma|^{p_K}\big|\Theta_{s_0}(\varsigma)\big|\,d\varsigma\;<\;\infty .
\]
\emph{At $s_0=1$ clause (ii) \textbf{is} clause (i)}, so the interior kernel is $\Theta_1$ and there is
no second convention. \emph{At a point within $c_Ch$ of the \textbf{right} endpoint of a compact interval
the kernel used is the reflection $\varsigma\mapsto\Theta_{s_0}(-\varsigma)$}, supported in
$[-c_C,c_Cs_0]$, whose truncated moments are $(-1)^{j}\delta_{j0}=\delta_{j0}$; no separate clause is
needed and the reflection is recorded here rather than left implicit.
\item[(iii)] \emph{Regularity at the interior index.} $\Theta_1$ is symmetric and Lipschitz, with
Lipschitz constant $\mathsf c_4$.
\end{enumerate}
The same three conditions with the same $p_K$ apply to the kernel $C_{h_x}$ used in the $x$-direction
inside $\hat{\mathcal A}_{\mathrm{cond}}$ and $\hat{\mathcal A}^\ast_{\mathrm{cond}}$, both as defined in
\cref{def:hatA}.
\end{assumption}

\noindent\emph{Notes on C3.} What \eqref{eq:rescaled} is for, and how clause (ii)'s boundary index and
defining property are read, are \cref{rem:app3read}.

% ROUND-12 RESTATEMENT OF C3.  THREE THINGS CHANGE AND EACH CLOSES A NAMED
% FINDING; NOTHING PROVED UNDER THE OLD C3 IS LOST, BECAUSE EVERY BOUND
% CONSUMED WAS |Theta_{s_0}| <= c_2 ON A SET OF LENGTH <= 2 c_C PLUS THE ORDER
% CONDITIONS, AND BOTH SURVIVE VERBATIM.
%
%  (a) Theta_{s_0} IS NOW DEFINED IN THE ASSUMPTION, eq:rescaled.  It was
%      previously introduced as a "standing convention (H2c)" in
%      archive/proofs_partH.tex -- stamped DO NOT COMPILE -- and used at six live
%      sites in proofs.tex with NO DECLARATION ANYWHERE IN A LIVE FILE.  That is
%      the assumption-auditor's U-1 and the regularity-checker's A8, and the
%      repair is to state the object where it belongs.  proofs.tex now carries
%      (P9) as a pointer to eq:rescaled rather than as an independent
%      declaration.
%
%  (b) int Theta_{s_0} = 1 IS PRINTED.  The old C3 said only "of order p_K" and
%      the normalisation is what "of order" means for a generalized kernel; it is
%      consumed at lem:intker(i), (iii) and (iv), so leaving it implicit charged
%      three clauses to a phrase.  The old (H2c) note (a) already said it should
%      be printed.
%
%  (c) THE BOUNDARY INDEX NOW HAS A DEFINING PROPERTY, AND THE C3(ii)/C3(iii)
%      CLASH IS DISSOLVED RATHER THAN RE-FILED.  Three incompatible meanings of
%      s_0 were in play: the OLD C3(ii)'s abstract index, whose support clause
%      [(s_0-1)/h, s_0/h] cap [-c_C, c_C] is VACUOUS for fixed s_0 in (0,1) and
%      small h and ONE-SIDED at s_0 = 1, contradicting (iii)'s SYMMETRY; the
%      literature's boundary distance IN BANDWIDTHS; and lem:intker(iv)'s proof,
%      which chose s_0 = (x* - underline-x)/(bar-x - underline-x), a FRACTION OF
%      THE INTERVAL and dimensionally wrong.  The new (ii) is the second of the
%      three, NORMALISED so that s_0 = 1 remains the interior index: support
%      [-c_C s_0, c_C], i.e. left truncation at the boundary distance measured in
%      bandwidths and divided by c_C.  At s_0 = 1 the support is exactly
%      [-c_C, c_C], TWO-SIDED, so (iii)'s symmetry is consistent and the clash
%      DISAPPEARS -- it was an artefact of the old interval, not a real ambiguity.
%      THE OLD VERDICT "nothing depends on which is meant" WAS FALSE, and it was
%      false at exactly one place: lem:intker(iv)'s proof uses the two-sided
%      reading at interior points and the one-sided reading at boundary points.
%      Under the new (ii) both are the SAME clause at different s_0.
%
%  (d) WHAT (ii) NOW SUPPLIES AND THE OLD C3 DID NOT: eq:truncmom, the TRUNCATED
%      moment conditions.  This is the defining property of a boundary-adjusted
%      generalized kernel and it is what lem:intker(iv) needs; the old C3 pinned
%      nothing of the kind, and (iv)'s proof appealed to "the rescaling implicit
%      in C3(ii)", which was not a statement.  LOAD-BEARING:
%      prp:hatAprop(iii) -- that A-hat* maps into V -- rests ENTIRELY on (iv),
%      and its own proof says that without it the identity fails by an O(1)
%      amount and A-hat* alpha leaves the parameter space, which (L2) inherits.
%
% The draft's assapp3a is otherwise retained, with three renamings: its order
% symbol "l" becomes p_K (l is one glyph from ell, the debiasing weight); its
% boundary-adjustment parameter "t" becomes s_0 (t is the Bierens index); and its
% support interval, written C_0 in round 2, is WRITTEN OUT as [-c_C, c_C].
% ROUND-7: THE SUPPORT INTERVAL LOSES ITS SYMBOL.  Calligraphic C_0 collided with
% cor:udep(c)'s admissible weight class -- a set of functions on R^p -- and the
% glossary carried both rows without noticing.  The kernel's support interval is
% a purely local object mentioned twice in one assumption, so it is written out,
% exactly as C12 wrote out the compact cube [-c_Phi, c_Phi]^p rather than naming
% it.  c_C joins c_Phi and c_Pi in the constant family.  The symbol C_0 is
% therefore free and belongs to cor:udep alone.
%
% THE DRAFT'S assapp3b IS RETIRED.  It introduces a discrete product kernel for
% smoothing over W, with bandwidth h_w, and a Dirac-delta representation of
% f_W.  Under the architecture of plan Sec.17, W enters ONLY through Wtilde
% (a multiplication) and, inside A-hat, through exact cell-splitting on the
% binary W, which A4(i) makes feasible.  So h_w = 0 identically, assapp3b
% reduces to the indicator, and its symbol L_{h_w} -- which collides with the
% debiasing weight L -- disappears with it.  The Dirac device at draft `:257'
% goes too, removing a third meaning of Delta.

\begin{assumption}[C4$'$: Residual truncation]\label{ass:app4}
Let $\ell_n$ be a real sequence that is either bounded or diverges slowly to $\infty$ with $n$. The
marginal density $f_{\Delta U}$ exists, is continuous, and is strictly positive on $[-\ell_n,\ell_n]$
for every $n$; set $\kappa_n:=\inf_{|u|\le\ell_n}f_{\Delta U}(u)$, so that $\kappa_n>0$. If $\ell_n$ is
bounded then $\inf_n\kappa_n>0$.
\end{assumption}

% STATUS OF C4' -- PRINTED, NOT LEFT HERE.  ell_n and kappa_n occur nowhere in
% this project outside the clause above, and nothing is truncated to
% [-ell_n, ell_n].  The clause is INERT and is NOT RETIRED; the printed statement
% of that standing is in the Discussion, under "Five clauses of Part VI that no
% theorem rests on", which is the A5(v) model applied uniformly to C1(i), C4',
% C5'(ii), C5'(iii) and C5'(iv).
%
% THREE CORRECTIONS TO THE DRAFT'S assapp4.
%  (a) The draft requires kappa_n -> 0 AS PART OF THE ASSUMPTION.  That is not
%      an assumption anyone would want: it is a statement that the truncation
%      degenerates.  It also contradicts the same assumption's own "either
%      bounded or diverges slowly", since with ell_n bounded kappa_n is a
%      positive constant.  The requirement is struck.
%  (b) ROUND 2: kappa_n > 0 is a CONCLUSION and needs a hypothesis.  Continuity
%      alone gives an infimum over a compact set that is attained; strict
%      POSITIVITY of f_{Delta U} on [-ell_n, ell_n] is what makes it nonzero,
%      and it is now stated.  C2'(ii) gives smoothness of M, not positivity of
%      the marginal, so it cannot supply this.
%  (c) The tail caveat.  With ell_n bounded there is NO rate penalty.  If ell_n
%      diverges the penalty is SUB-POLYNOMIAL, not logarithmic: for Gaussian
%      f_{Delta U} and ell_n ~ sqrt(2 c log n), kappa_n^{-1} ~ n^{c}.  The plan
%      Sec.2.6 echo asserting a kappa_n^{-1} rate penalty is struck (memo Sec.9
%      item 11).

\begin{assumption}[C5$'$: Smoothing parameters]\label{ass:app5}~
\begin{enumerate}
\item[(i)] \emph{Residual direction, operator component.} $p_K=8$ and
\[
h_u\;\asymp\;\mathrm{sd}(\Delta U)\,\hat\delta_n^{1/p_K}\;=\;\mathrm{sd}(\Delta U)\,n^{-\alpha/8},
\qquad \alpha:=\frac{r}{2r+d_Z}\in\big(0,\tfrac12\big),
\]
% WHY h_u -> 0 AND h_u^{p_K} asymp hatdelta_n FOLLOW from the display above, and
% why no lower bound on h_u is required for this component: rem:app5read(1).
\item[(ii)] \emph{$x$-direction, derivative and adjoint components.} $h_x\to0$ and
\[
\big(n\,h_x\,h_u\big)^{-1}\ln n\;\longrightarrow\;0 .
\]
% EXTRACTED FROM THE STATEMENT OF C5' ON 2026-08-03.  Fourteen lines of WORKED
% ARITHMETIC stood here -- which lemma consumes the displayed form, the stronger
% inherited form and where it may be imposed, the (P-g) upper bound it must meet,
% and the resulting window on s.  A derivation is not a condition, so it is now
% rem:app5read(2), at the foot of this file.
% THE DISPLAY eq:hxwindow MOVED WITH IT, INTACT AND WITH ITS OWN \label; it stays
% declared exactly once in this file, and all 12 citing sites are \eqref.
\item[(iii)] \emph{Propensity.} $\hat\varrho$ satisfies A4(iii) with $r$ obeying the
frontier \eqref{eq:frontier}.
\item[(iv)] \emph{No further smoothing.} No bandwidth is used in $t$, in $\Delta Z$, or in $W$.
\end{enumerate}
\end{assumption}

\noindent\emph{Notes on C5$'$.} What clause (i)'s display buys, and clause (ii)'s stronger inherited form
together with the window on $s$ it leaves against (P-g), are \cref{rem:app5read}.

% STATUS OF CLAUSES (ii), (iii), (iv) -- PRINTED, NOT LEFT HERE.
%
% THE HEADER OF THIS NOTE WAS FALSE THROUGH DISPATCH 20 AND IS CORRECTED HERE.
% It read "No hypothesis list in this corpus carries any of the three."  THREE
% LISTS CARRIED ALL THREE, by block citation of the parent tag C5': lem:aerr,
% prp:phirate and rem:rootnhyp(3) -- the last being thm:rootn's (T-5) inheritance.
% Dispatch 19's own register rule states that a block citation of the parent tag
% IS a citation of every clause under it; the same dispatch narrowed B1 and C1 on
% those lines by clause and left C5' un-narrowed.  All three now read C5'(i), and
% the header is true as of this dispatch and not before.
%
% (ii) is consumed ARITHMETICALLY, as the lower end of the admissible h_x window,
% at rem:tccwhy and lem:scaleinv, and lem:aerr reading (i) shows the operator-error
% bound needs only (n h_x h_u)^{-1} ln n -> 0, which is now what the clause states.
% (iii) is a POINTER to A4(iii) and eq:frontier, both cited directly wherever the
% propensity rate is needed.  (iv) is a SCOPE DECLARATION discharged by
% def:hatTexp and def:hatAcon.  NONE OF THE THREE IS RETIRED.  The printed
% statement is in the Discussion, under "Five clauses of Part VI that no theorem
% rests on" -- the A5(v) model, applied uniformly to C1(i), C4', C5'(ii),
% C5'(iii) and C5'(iv).
%
% assapp5 WAS WRITTEN FOR A DIFFERENT OPERATOR AND HAD TO BE REBUILT.  The
% draft's condition is  h_u, h_w, h_x -> 0  and  (n h_x^p h_u^3)^{-1} ln n -> 0.
% Three things are wrong with it here.
%
%  (a) THE EXPONENT p IS NOT dim(Z).  It was dim(x) in the cross-sectional
%      companion.  In the panel X is SCALAR and x = (x_1,x_2), so the exponent
%      is 2.  Plan Sec.1's reserved-symbol table lists draft `:302' under
%      "p = dim(Z)", which is a mis-attribution.
%
%  (b) THE OPERATOR NEEDS NO LOWER BOUND ON h_u.  T(varphi,theta)(u,t) is an
%      UNCONDITIONAL expectation of a bounded summand: there is no
%      local-averaging denominator, so the stochastic error is O_P(n^{-1/2})
%      pointwise and O_P(n^{-1/2} sqrt(ln n)) uniformly over (u,t), IRRESPECTIVE
%      of h_u.  The smoothing bias is O(h_u^{p_K}) under C2'(ii).  At p_K = 8
%      and h_u ~ n^{-1/16} the two balance exactly at n^{-1/2}.  This is where
%      clause (i) comes from, and it RECONCILES plan Sec.16.1's inherited
%      bandwidth choice with the theory rather than treating it as a MATLAB
%      artefact.
%
%      GUARD, REWRITTEN IN DISPATCH 7 BECAUSE THE OLD GUARD PROTECTED A MATCHING
%      TO THE WRONG TARGET.  ROUND-9 TEXT, STRUCK: "h_u ~ n^{-1/16} is simply the
%      largest bandwidth whose bias h_u^8 sits at the FLOOR n^{-1/2}."  Two parts
%      of that survive and one does not.
%       SURVIVES.  The exponent comes from matching bias to a level, NOT from a
%       bias-variance trade-off.  Because there is no local-averaging
%       denominator, the stochastic term does not inflate as h_u decreases, so
%       the standard balance h_u^8 against (n h_u)^{-1/2} -- which would give
%       h_u ~ n^{-1/17} -- DOES NOT APPLY HERE.  A later editor must not
%       "correct" the exponent by that route.
%       DOES NOT SURVIVE.  THE LEVEL TO MATCH IS delta-hat_n, NOT n^{-1/2}.  The
%       stochastic floor n^{-1/2} is what eq:opnoise's FIRST term sits at; the
%       operator noise the whole of Step 6 is denominated in is
%       delta-hat_n = max{n^{-1/2} sqrt(ln n), r_n}, and A5(i) records that in
%       every configuration of interest delta-hat_n = r_n >> n^{-1/2}.  Matching
%       the bias to n^{-1/2} therefore chooses a bandwidth SMALLER than necessary
%       by a factor n^{(1 - 2 alpha)/16}, and pays for it: (P-f) of prp:phirate
%       is delta_n^2 = O(h_u^2 delta-hat_n ln n), which is HARDER the smaller
%       h_u is.  At h_u ~ n^{-1/16} it reads alpha (2 mu - 1)/(2 mu + 1) > 1/8 --
%       a JOINT restriction on (mu, alpha) that is IMPOSSIBLE at mu = 3/4 (it
%       would need alpha > 5/8 > 1/2) and FAILS at the BOUNDARY configuration
%       (mu, d_Z, r) = (1, 3, 4.5) once logarithms are counted.  At the amended
%       h_u ~ delta-hat_n^{1/8} it reads mu > 5/6, free of alpha.
%      WHAT THE AMENDMENT COSTS.  (1) The bandwidth now depends on alpha, hence
%      on r and d_Z, which the analyst must supply -- the same input A4(ii)-(iii)
%      already require, so this is not a new infeasibility, but it IS a new
%      dependence and is stated rather than hidden.  (2) The smoothing bias is no
%      longer negligible against the stochastic floor: it is asymp delta-hat_n,
%      which is all eq:opnoise needs (lem:opnoise(ii)) but which means
%      prp:stopfeas(i) closes at delta-hat_n and not at the sharper
%      delta-hat_n^{(T)}.  See the A5(i) comment, clause (c).
%      (3) mu > 5/6 STRICT joins the smoothness budget table.
%
%  (c) h_w DISAPPEARS (see the note under C3), and no bandwidth in t or in
%      Delta Z ever existed: the Bierens device replaces conditioning on
%      Delta Z by the scalar index t'Phi(Delta Z), and the propensity is a
%      SCALAR regression, priced separately in (iii).
%
% Clause (ii) is CONSERVATIVE and is flagged as such.  The exponent 3 on h_u is
% inherited from the companion paper's density-derivative estimator; how many
% u-derivatives the Landweber step actually needs is prp:phirate's business
% (Step 6), and the exponent should be re-derived there.  A stronger condition
% is safe in the meantime.
%
% STATUS OF CLAUSE (ii): THE RE-DERIVATION WAS DONE IN DISPATCH 7, THE RELAXATION
% WAS DECLINED THERE, AND IT IS TAKEN NOW.  lem:aerr consumes
% (n h_x h_u)^{-1} ln n -> 0 -- ONE power of h_x and ONE of h_u -- and eq:gamman
% shows why: the exponents 2 and 3 pair with an estimator of a DENSITY DERIVATIVE
% in x and a SECOND derivative in u, which is the cross-sectional companion's
% object; A-hat_varphi estimates a density-weighted CONDITIONAL MEAN.
%
% THE DISPATCH-7 GROUND FOR KEEPING THE STRONGER FORM WAS AN INVALID INFERENCE
% AND IS STRUCK.  It read: "Clause (ii) as written IMPLIES the relaxed condition
% whenever h_x, h_u <= 1, so nothing proved is lost by keeping it."  The premise
% is true and the conclusion does not follow.  A LOWER bound that is stronger
% than needed is free only in isolation; it is NOT free once it must intersect an
% UPPER bound imported from elsewhere, because the strengthening moves the floor
% up towards a ceiling it does not control.  The ceiling here is prp:phirate's
% (P-g), R h_x^{s-1/2} = O(delta-hat_n), and the intersection is EMPTY at
% admissible parameters.
%
% THE WITNESS, AND IT IS ADMISSIBLE.  (mu, s, a, r, d_Z) = (2, 1.2, 0.3, 8/3, 3),
% so alpha = r/(2r+d_Z) = 8/25 = 0.32 and kappa = 2 mu/(2 mu + 1) = 4/5.
%   FRONTIER  eq:frontier: kappa alpha = 0.256 > 1/4.               CLEARS.
%   eq:lxwindow: s > 9 alpha / (8 (1 - 2 alpha)) = 1.  s = 1.2.      CLEARS.
%   STRONG FORM (n h_x^2 h_u^3)^{-1} ln n -> 0 at h_u ~ n^{-alpha/8} = n^{-0.04}:
%       h_x >> n^{-0.44}.
%   (P-g) at s - 1/2 = 0.7:  h_x = O(n^{-0.32/0.7}) = O(n^{-0.45714}).
%   WINDOW EMPTY, by the exponent gap 0.01714.
% The threshold is eq:hxwindow, s > 1/2 + 16 alpha/(8 - 3 alpha), which at
% alpha = 0.32 is s > 1.22727 and is FAILED by s = 1.2.  Emptiness holds on a band
% of s that is NON-EMPTY at every alpha in (1/4, ~0.348) -- and eq:frontier forces
% alpha > 1/4 always, so the band is never vacuous at the low end.  The band
% closes at alpha ~ 0.348 and is empty at alpha = 3/8, which is why the block's
% two spot-checks below missed it: both are run at s = 3 and s just above 2.
%
% WHAT THIS IS *NOT*.  It is NOT a counterexample to lem:aerr, prp:phirate or
% thm:rootn.  At the witness the block-cited hypothesis sets of those results were
% UNSATISFIABLE, so each held VACUOUSLY; nothing proved in this corpus is false at
% (s, alpha) = (1.2, 0.32).  What the witness exhibits is an admissible
% configuration at which a MAINTAINED clause and an imported proposition
% hypothesis have no common solution -- a defect of the conditions, not of a
% theorem.  Recorded so that no later reader converts it into the stronger claim.
%
% THE TWO REPAIRS, BOTH TAKEN.  (a) The three hypothesis lists that block-cited
% C5' now cite C5'(i), so no hypothesis set in the corpus contains the conflict.
% (b) The clause itself now states the weaker form, with the stronger one retained
% as an option, gated on eq:hxwindow -- so the MAINTAINED BLOCK read as a whole is
% satisfiable at every admissible (s, alpha) as well, which (a) alone does not
% deliver.  Both are needed and neither substitutes for the other.
%
% WHAT THE RETENTION OF THE STRONGER FORM IS STILL FOR.  It is the only place in
% the block where h_x carries a LOWER bound, and two printed computations read the
% h_x window at that end: rem:tccwhy(3) and lem:scaleinv's cap inflation.  Both are
% read at s = 3, where eq:hxwindow holds with room, so both stand unchanged.
% RECORDED, NOT DONE, AND STILL NOT DONE: promoting (P-g) into C5' as a clause.
% Whoever performs it must carry eq:hxwindow with it, since promotion puts the
% floor and the ceiling in the same maintained block.
%
% DISPATCH-8: THE WINDOW BELOW IS RECOMPUTED AT THE AMENDED BANDWIDTH.  The
% round-10 version of this comment read the window at h_u ~ n^{-1/16}, which
% C5'(i) RETIRED in the same pass, and at s > 1.75, which the joint (P-c) ^
% (P-d) ^ A8 line in the budget table now rules out as a DESIGN (it clears
% eq:lxwindow but not A8 in full).  Both numbers are replaced.
%
%   AT THE FRONTIER-CLEARING CONFIGURATION (mu, s, a, alpha) = (2, 3, 3/4, 0.32),
%   which is the admissible one:
%     h_u ~ n^{-alpha/8} = n^{-0.04}.
%     Clause (ii), (n h_x^2 h_u^3)^{-1} ln n -> 0, is
%         h_x >> n^{-1/2} h_u^{-3/2} (ln n)^{1/2} = n^{-0.44} (ln n)^{1/2}.
%     (P-g), R h_x^{s - 1/2} = O(delta-hat_n), is h_x = O(n^{-alpha/(s - 1/2)})
%         = O(n^{-0.128}).
%     WINDOW: (n^{-0.44}, n^{-0.128}) -- non-empty with room.
%     The RELAXED lower bound lem:aerr consumes, (n h_x h_u)^{-1} ln n -> 0, is
%         h_x >> n^{-0.96} ln n, which is comfortable, so the relaxation buys a
%         much wider window and is available whenever it is wanted.
%
%   AT THE BOUNDARY CONFIGURATION (mu, d_Z, r) = (1, 3, 4.5), alpha = 3/8, kept
%   only as the clause-(0) illustration and NOT a design:
%     h_u ~ n^{-3/64}.  Clause (ii) gives h_x >> n^{-55/128} = n^{-0.4297};
%     (P-g) at s just above 2 gives h_x = O(n^{-alpha/(s - 1/2)}) = O(n^{-0.25}).
%     WINDOW: (n^{-0.4297}, n^{-0.25}) -- non-empty.  Under the RETIRED
%     bandwidth h_u ~ n^{-1/16} the lower end was n^{-13/32} = n^{-0.40625} and
%     the upper end, at the then-quoted s > 1.75, was n^{-0.3}; the amended
%     bandwidth WIDENS the window at the lower end and the corrected s
%     requirement NARROWS it at the upper end.
%
% Whoever promotes (P-g) into C5'(ii) must re-run these two lines, because both
% ends move with h_u and with s.

\paragraph{Discussion.}
A0 and A1 are the sampling and exogeneity primitives; A1 is the draft's own \texttt{ass:condind} and is
the weaker of the two candidate exogeneity conditions, since it permits the instrument to be assigned on
characteristics that also enter $\Delta Z$. A1b(iii) rules out a second structural function generating
the same conditional law of the residual given $(W,\Delta Z)$; A1b(iv) is what identifies
$\theta_\dagger$, and it involves no source condition and no debiasing weight; A1c is the linearisation
of A1b(iii) and governs the local behaviour of the regularised inverse. B1 is a set of conditions on
objects the analyst \emph{chooses}, not on the data-generating process, with the exception of
B1(iii)--(iv), which restrict how coarse an index family may be used. \emph{B1(vi) is \textbf{retired} and is
not part of that set}: it was an identification restriction, it is not a condition of this estimator, and it
is not available as a hypothesis --- see the struck clause and the reason recorded with it. A2 is the
Severini--Tripathi condition and is the binding
assumption for $\sqrt n$ \emph{through this moment function}: in the linear model its necessity for
$\sqrt n$-estimability is known \citep{severini2006}, but in the present nonlinear problem necessity
remains conjectural, and two claims must be kept apart --- necessity for this moment, and necessity for
$\sqrt n$-estimability by \emph{some} estimator. A2$'$ adds what is needed for the second-order term to
be a finite number. \emph{A3 is retired, and it never ruled out ``a design in which the
debiased moment has no local information about $\theta_\dagger$''.} By \cref{lem:momzero}\textup{(iv)} the debiasing
% ARCHAEOLOGY.  A3 was described here through dispatch 12 in those terms, which was
% an overstatement of what it could rule out.
term contributes nothing to the $\theta$-Jacobian at all --- $E[\psi(\theta,\zeta_\star)]$ is affine with
gradient $-E[\Delta Z\Delta Z']$ --- so \emph{all} of the moment's information about $\theta_\dagger$
comes from restriction (B), which is what \cref{lem:shiftinv} and \cref{rem:R1} said in words from Step~1
onward and what A1b(iv) prices. \emph{The debiasing term buys orthogonality in $\varphi$, not information
about $\theta$}, and \eqref{eq:sigsplit} shows it is paid for in variance. A4 rules out a design in which
the instrument is degenerate on part of the support of the controls. The application's $689/300$ split is not itself a difficulty: at $\varrho\approx0.303$ the
factor $\varrho(1-\varrho)\approx0.211$ costs $0.25/0.211\approx1.18$ relative to a balanced design ---
about $18\%$ in variance and about $9\%$ in standard error. A4(i) binds only if $\varrho$ approaches $0$
or $1$ somewhere on the support of $\Delta Z$, which is a question about the conditional propensity and
not about the marginal split. N1 is a \emph{property} of a debiasing weight, not a normalisation built
into it: it can always be arranged by re-anchoring, at no cost to $\partial_uL$ or to $E[\psi]$, and
\cref{prp:orthog}(iii) shows that it is \emph{necessary and sufficient} for exact orthogonality in the
propensity. Stating it as a property rather than as a construction is what leaves that equivalence with
content. A5 collects what the regularised first step must deliver; clauses (i)--(iii) are
discharged, not assumed, by \texttt{prp:phirate}, and clause (iii) is the paper's real feasibility
frontier. \emph{Clauses (vi) and (vi$'$) are not discharged}: they are unproven inputs, and from Step~6
onward they are consumed on \emph{both} routes --- (vi) by \cref{lem:balance}(iii) on the balancing
route, (vi$'$) by \cref{lem:pathcont} on the rate route, which needs the Sobolev bound at every stage
and not only at the final iterate.
\emph{A7 is the condition under which a discrepancy-principle stopping rule is a theorem rather than a
heuristic when the noise sits in the operator, and it is now the only one}: A6 is retired, having been
refuted for the operator it named (\cref{prp:tccfails}), and the stopping theorem
\cref{prp:phirate} runs on A7 in the amended form \eqref{eq:a7amend} together with the
stage-adaptive schedule the estimator computes --- with \emph{no} cone condition anywhere. A8 fixes the
sieve window and, with it, the basis; from Step~6 onward its \emph{basis clause} and its \emph{upper
end} are consumed by the $\varphi$-step as well as by the balancing solve, through
\cref{lem:srctransfer}(iii) and \cref{lem:unifop}(3d) --- \emph{but its lower end is not}, and
\cref{prp:phirate} cites A8 by clause for that reason; (P-c) is the rate route's own lower end. C1--C5 are
regularity conditions for uniform consistency of the kernel components, carried over from the draft's
appendix and restated for the panel operator.

\paragraph{Five clauses of Part VI that no theorem rests on, and the status of each.}
\emph{These five are collected here because a maintained clause with no theorem behind it is exactly what a
reader cannot check by grepping, and because the block already does this for A5(v)
(\cref{rem:rootnscope}\textup{(4)}) --- that is the model, and the treatment below is uniform with
it.} \emph{None of the five is retired, and the distinction matters}: this block has four retirements resting
on three grounds --- A3 and A$2'$(ii) because they are \emph{proved}, A6 because it is
\emph{false}, B1(vi) because it is \emph{about a different estimator} --- and none of the five below meets any
of the three. Unconsumed is not retired, and only a retired clause is unavailable as a hypothesis. What each
of the five is instead:
\begin{itemize}
\item[--] \textbf{C1(i)} has exactly one consumer in the corpus, Step~2 of \cref{lem:aerr}, where it supplies
the $p_K$ derivatives of the two products $E[\breve g_{u,t}\mid X_j=\cdot]f_{X_j}$ --- \emph{the breve
marks the \textbf{infeasible} version of a hatted object, the one built at the true $\varrho$ rather
than at $\hat\varrho$; it is declared at \cref{lem:opnoise} Step~1, which is where the two stems that
carry it are defined, and this is its first occurrence in compiled order} --- behind that step's
$O(h_x^{p_K})$ $x$-smoothing bias. \emph{That consumer is itself consumed by nothing} --- \cref{lem:aerr}
opens ``Consumers: none'' and is retained for the two reasons its reading~\textup{(iv)} gives --- so C1(i) is
a hypothesis of one lemma and, transitively, of no theorem. It is kept because that lemma's proof needs it.
\emph{\Cref{prp:phirate} does \textbf{not} consume C1(i)}.
% ARCHAEOLOGY.  A comment in this file asserted that it did, through dispatch 17.
\item[--] \textbf{C4$'$.} $\ell_n$ and $\kappa_n$ occur nowhere in this project outside C4$'$ itself, and no
object anywhere in these files is truncated to $[-\ell_n,\ell_n]$. The clause prices a residual truncation
the estimator does not perform. \emph{It is inert}: no bound stated anywhere degrades if it is deleted. It is
nevertheless not retired, because \emph{inert} and \emph{unavailable} are different registers --- a truncated
variant of the residual direction would need exactly this clause, and \cref{def:a6b} is carried in this block
on the same footing, being ``never asserted and consumed by no theorem''.
\item[--] \textbf{C5$'$(ii).} No hypothesis list consumes it --- \emph{and that statement is true only from
this dispatch}: \cref{lem:aerr}, \cref{prp:phirate} and \cref{rem:rootnhyp}\textup{(3)} block-cited C5$'$ and
therefore carried this clause, on the register rule stated with the retirements above, and all three now cite
C5$'$(i). \Cref{lem:aerr} reading~\textup{(i)} shows that the one operator-error bound stated at this pair of
bandwidths needs only $(nh_xh_u)^{-1}\ln n\to0$, which is what the clause now states.
\emph{What it does do is fix the lower end of the admissible $h_x$ window}, and two printed computations turn
on that end: \cref{rem:tccwhy}'s third reading,
which collects C5$'$(ii) with A8, \eqref{eq:entropysieve}, (P-c) and \cref{lem:unifop} Step~4 and concludes
that $h_xL_x$ is free in both directions, and \cref{lem:scaleinv}'s cap inflation, which reads $\bar m_n$
across the $h_x$ window that C5$'$(ii) and (P-g) leave. Both are computed at the strengthened form of the
clause and at $s=3$, where \eqref{eq:hxwindow} holds with room, so both stand. So C5$'$(ii) is consumed
\emph{arithmetically} and not as a hypothesis, and the two readings must be kept apart: an auditor grepping
hypothesis lists will find nothing, and deleting the clause would nonetheless leave two displays without a
lower limit. \emph{And a lower bound stated more strongly than any proof needs is \textbf{not} free here},
which is the correction \eqref{eq:hxwindow} records: it must intersect the upper bound (P-g) imports, and the
strengthened form and (P-g) have no common solution at any $s\le\tfrac12+16\alpha/(8-3\alpha)$ ---
a band that is non-empty at every $\alpha\in(\tfrac14,\approx0.348)$, and \eqref{eq:frontier} forces
$\alpha>\tfrac14$.
\item[--] \textbf{C5$'$(iii)} is a pointer and not a restriction. It repeats A4(iii) and
\eqref{eq:frontier}, and every result that needs the propensity rate cites those two directly. Nothing cites
C5$'$(iii) and nothing should; it is here so that the smoothing-parameter block can be read on its own.
\item[--] \textbf{C5$'$(iv)} is a scope declaration about the estimator, not a restriction on $P$. It is
discharged by construction --- \cref{def:hatTexp} and \cref{def:hatAcon} introduce $h_u$ and $h_x$ and no
bandwidth in $t$, in $\Delta Z$ or in $W$ --- so there is nothing for a proof to cite it for.
\end{itemize}
\emph{Why a status sentence and not a retirement, applied to all five alike.} A retirement is a claim that the
clause is unavailable as a hypothesis, and for C1(i) and C5$'$(ii) that claim would be false --- the first is a
live hypothesis of \cref{lem:aerr} and the second is what puts the lower end on the $h_x$ window the two
printed computations above read. Retiring the remaining three alone would file clauses of one species under two
registers and leave a reader inferring the difference. \emph{The uniform remedy is the one A5(v) already
carries: state the standing, in one place, in the printed text.}

\paragraph{The smoothness budget, in one place.}
Four smoothness parameters appear above and they do not compete for the same budget once H1 is imposed:
\begin{center}
\begin{tabular}{llll}
\toprule
Condition & In terms of $(s,a,r)$ & Route & Source \\
\midrule
Mild ill-posedness\textsuperscript{\dag} & $s>a+\tfrac12$ & both & plan \S2.6 scope \\
$\mu_{\text{\textup{H1}}}>\tfrac12$\textsuperscript{\dag} & $s>a$ & both & implied by the line above, with room \\
Sieve window non-empty & $s>1$ (splines); $s>\tfrac32$ (polynomials) & \emph{balancing only} & A8 \\
Sieve conditioning ($L_x\ll n^{1/2}$) & --- & both & A8 \\
Propensity frontier & $\dfrac{2\mu}{2\mu+1}\cdot\dfrac{r}{2r+d_Z}>\dfrac14$ & both & A5(iii) \\
Kernel bias & $u\mapsto\mathcal M(u,t)$ in $C^{p_K}$ uniformly in $t$, $p_K=8$ & both & C2$'$(ii) \\
Representer & source index $\mu''>\tfrac32$ for $\ell^\Pi_\star$, independent of $\mu$ & both & A2c, \cref{lem:reprsmooth} \\
\emph{Linearisation remainder} & $\mu>\tfrac56$ \emph{strict}, at C5$'$(i) as amended & \emph{rate only} & (P-f), \cref{rem:h6bandwidth} \\
\emph{Sieve entropy} & $L_x\ln n=o\big(n\hat\delta_n^{2}\big)$ & \emph{rate only} & (P-d), \eqref{eq:entropysieve} \\
\emph{$x$-smoothing bias} & $\mathsf R\,h_x^{s-1/2}=O(\hat\delta_n)$ (an \emph{upper} end for $h_x$) & \emph{rate only} & (P-g), \cref{lem:unifop} Step 4 \\
\emph{Sobolev stability, stagewise} & $P\big(\max_{k\le\bar k}\|\hat\varphi^{(k)}-\vdagger\|_{H^{s}}\le\mathsf R\big)\to1$ & \emph{both} & A5(vi$'$) \\
\midrule
\emph{$L_x$ window, rate route} & $s>\dfrac{9\alpha}{8(1-2\alpha)}$ & \emph{rate only} & (P-c)$\,\wedge\,$(P-d), \eqref{eq:lxwindow} \\
\emph{$L_x$ window, both routes at once} & $s>\dfrac{1}{2(1-2\alpha)}$ & \emph{both, jointly} & A8 lower end $\wedge$ (P-d) \\
\midrule
\emph{Kernel order vs.\ Sobolev order} & $p_K\ge s-\tfrac12$, i.e.\ $s\le\tfrac{17}{2}$ at $p_K=8$ & \emph{rate only} & \cref{lem:unifop} Step 4, gates (P-g), (N1) \\
\emph{Source index, now primitive} & $\mu>\tfrac56$ at (P-a.3), \textbf{not} derived from $s/(2a)$ & \emph{rate only} & (P-a.3), \cref{rem:pasrc} \\
\bottomrule
\end{tabular}
\end{center}
\noindent\textsuperscript{\dag}\emph{These two lines carry $\mu_{\text{\textup{H1}}}=s/(2a)$, H1's index at
the \textbf{population} operator $\mathcal A_{\mathrm{cond}}$ on all of $\mathcal V$ at $\theta_\dagger$ ---
\textbf{not} the primitive $\mu$ of (P-a.3), which is what every other $\mu$ in this table and in
\eqref{eq:frontier} means.} They are \emph{interpretations} and transfer to the primitive index only under
the commutativity proviso \cref{rem:pasrc}\textup{(4)} states and declines. \cref{rem:pasrc}\textup{(5)}
carries the binding reading.

\noindent Two readings of the table are needed rather than one. \emph{The ``Route'' column is not decoration:} by
\cref{lem:balance} the first-order bias is controlled by \emph{either} A5(ii)+A5(iv) \emph{or}
A8+\eqref{eq:nuhat}+A5(vi), not by both, so the $s>1$ line is charged on the balancing route only. On the
rate route the $s$ budget faces the first two lines, the frontier, the representer condition and
\emph{the last two rows of the table}, and \emph{nothing from A8's lower end} --- which is why
\cref{prp:phirate} cites A8 by clause, taking its basis and its upper end and \emph{not} its lower end.
The upper end of A8, which is the conditioning constraint on the
sieve solve, is charged on both routes and carries no $s$ requirement. The trade-off between
$\mu$ and $r$ in the frontier should be displayed rather than hidden: at
$(\mu,d_Z)=(1,3)$ the frontier reads $r>4.5$, at $(1,5)$ it reads $r>7.5$, and at $(2,3)$ it reads
$r>2.5$.

\emph{The representer line is the one that changed at Step 4, and the change is a discharge rather than a
budget.} Through round 8 it read ``source index $\mu''$ for $\ell^\Pi_\star$'' and pointed at
A2$'$(i) --- a \emph{primitive}, with no route to it. It now points at A2c and at
\cref{lem:reprsmooth}: the index is pinned at $\mu''>\tfrac32$, which is the one-dimensional
$C^{1}$-embedding threshold and not a free choice, and the line records a \emph{sufficient condition}
rather than an assumption. It still competes for no part of the $s$ budget, because $\mu''$ indexes
smoothness of $\ell^\Pi_\star$ in the residual and $s$ indexes smoothness of $\vdagger$ in $x$.

\emph{The four lines added at Step 6 should be read as one addition and three bookkeeping entries.}
\emph{The addition is $\mu>\tfrac56$}, and it \emph{subsumes} the $\mu_{\text{\textup{H1}}}>\tfrac12$ line,
which is therefore now slack rather than binding: under H1 and the scope condition $s>a+\tfrac12$ one has
$\mu_{\text{\textup{H1}}}>\tfrac12+\tfrac1{4a}$, so $\mu>\tfrac56$ is \emph{implied} whenever
$a\le\tfrac34$ --- \emph{under the interpolation proviso of \cref{rem:pasrc}\textup{(4)}, without which the
computed $\mu_{\text{\textup{H1}}}$ is not the $\mu$ of \eqref{eq:muhalf}}; see the dagger note above --- and is a
genuine extra restriction otherwise. \emph{It is strict}, and the strictness is not decoration: at the
\emph{boundary} configuration $(\mu,d_Z,r)=(1,3,4.5)$ the underlying condition (P-f) holds in powers of
$n$ with no margin and \emph{fails} once the logarithms A7's amended threshold introduces are counted
--- see \cref{rem:h6bandwidth}, clauses (1) and (5). \emph{The three bookkeeping entries} are the
sieve-entropy line (which replaces the stronger Sobolev-entropy condition \eqref{eq:entropy}, not
consumed by any result), the $x$-smoothing line (which is the block's \emph{only} upper bound on $h_x$
and is currently carried as a proposition hypothesis rather than as a clause of C5$'$), and the
stagewise Sobolev line (which is A5(vi$'$), an input and not a discharge). \emph{None of those four
competes with $s$ against the frontier}; they compete with $\mu$, with $L_x$ and with $h_x$
respectively.

\emph{Two further lines \textbf{do} compete with $s$.} (P-c) and (P-d) are the two ends of one window for
$L_x$, and reading
% ARCHAEOLOGY.  Through round 10 the text asserted that "the s budget is unchanged";
% it is not.  The two lines were added in the dispatch-8 pass.
them together at $h_u\asymp\hat\delta_n^{1/p_K}$ gives $n^{9\alpha/(8s)}\ll L_x=o(n^{1-2\alpha}/\ln n)$,
which is non-empty if and only if \eqref{eq:lxwindow}, $s>9\alpha/\big(8(1-2\alpha)\big)$. That is
$s>\tfrac{27}{16}=1.6875$ at $\alpha=\tfrac38$ and $s>1$ at $\alpha=0.32$: satisfiable, but at
$\alpha=\tfrac38$ it consumes most of the available $L_x$ margin, the admissible window there being
only $\big(n^{0.2411},\,n^{0.25}/\ln n\big)$. \emph{And if A8's \textbf{lower} end is carried as well
--- which it must be by any configuration on which both the rate route and the balancing route are to
be available --- the requirement is instead $s>\big(2(1-2\alpha)\big)^{-1}$, i.e.\ $s>2$ at
$\alpha=\tfrac38$}: at $\alpha<\tfrac49$ the two thresholds differ by
$(4-9\alpha)/\big(8(1-2\alpha)\big)>0$ and A8's is strictly the larger.
\emph{The practical consequence, and it is the reason the line is recorded rather than absorbed}: the
boundary configuration $(1,3,4.5)$ was priced at $s>1.75$ in earlier rounds, which clears
\eqref{eq:lxwindow} and does \emph{not} clear A8 in full. It is therefore not an admissible design for
a paper that keeps both routes open --- independently of the fact that it already fails (P-f) by
$(\ln n)^{1/3}$. The frontier-clearing configuration $(\mu,s,a,\alpha)=(2,3,\tfrac34,0.32)$ clears
every line in this table. \emph{Read that configuration as \textbf{four} numbers and not as three plus a
derivation}: $\mu=2$ is (P-a.3)'s primitive index, assumed at the compression of the realised operator at
$\theta_\circ$, while $(s,a)=(3,\tfrac34)$ fix the scale and the population decay at $\theta_\dagger$; that
$s/(2a)=2$ \emph{coincides} with the assumed $\mu$ is a property of the numbers chosen and is not an
identity this block can invoke, the bridge being the proviso \cref{rem:pasrc}\textup{(4)} declines. \emph{The
coincidence is deliberate and is worth keeping} --- a configuration on which the two indices agree is the
one a reader will find least surprising --- but it is not what makes the configuration admissible.
\emph{Mechanised at \texttt{scripts/symbolic/specs/phirate-sieve-joint-window.yaml}.}

\emph{Two lines were added in round 12, and the second changes how the whole table is read.}
\emph{The first is a genuine but slack budget line}: $p_K\ge s-\tfrac12$, which \cref{lem:unifop}
Step~4's Taylor expansion consumes and which no earlier round stated. At $p_K=8$ it permits every
$s\le\tfrac{17}{2}$, so it binds at no configuration this block contemplates --- but it is a condition,
it gates (P-g) and (N1), and it was previously supplied by an appeal to
$H^{s}\hookrightarrow C^{0,s-1/2}$ that is \textbf{vacuous} for every $s>\tfrac32$ and hence at every
admissible configuration. \emph{The rate $\mathsf Rh_x^{s-1/2}$ survives; its justification did not.}

\emph{The second is that $\mu$ is no longer a function of $(s,a)$ on the rate route.} (P-a.3) states the
source condition \textbf{spectrally}, at the compression \eqref{eq:compress}, with $\mu$ and
$\mathsf R_\ast$ primitive; the identity $\mu=s/(2a)$ is \emph{struck} from \cref{prp:phirate}'s
hypothesis line and re-appears in \cref{rem:pasrc}(4) as an interpretation under an explicitly stated
interpolation proviso. \emph{Three consequences for reading the table.} (i) The frontier line and the
$\mu>\tfrac56$ line are now conditions on a \emph{primitive} $\mu$, so they no longer translate
automatically into conditions on $s$ --- the entry ``$\mu=s/(2a)$'' has been removed from the frontier
row. (ii) \emph{H1 is not retired and the rest of the table is unchanged}: H1 still fixes the scale
$H^{s}$ that \cref{def:srcball}, A5(vi$'$), (P-c) and (P-g) are written in, still supplies A2's Picard
form, and still delivers the mild-ill-posedness and $\mu>\tfrac12$ lines \emph{as an interpretation}.
(iii) \emph{The honest reading of the collapse}: H1's matching is a statement about
$\mathcal A_{\mathrm{cond}}$ on all of $\mathcal V$, while (P-a.3) is about a finite-rank, sieve-compressed
operator at a random base point, and \cref{lem:srctransfer}(iii.a) records that a link condition
\emph{cannot} hold on $\mathcal V$. So H1 never delivered $\mu=s/(2a)$ for the operator (P-a) is about,
and round 12 makes visible an assumption the block was already carrying rather than adding one. See
\cref{rem:pasrc}, all four clauses.

% ---------------------------------------------------------------------------
%  THE REMARKS BELOW WERE EXTRACTED FROM ASSUMPTION STATEMENTS ON 2026-08-03,
%  under the house style (.claude/references/theorem-environment-style.md): an
%  ASSUMPTION CARRIES CONDITIONS ONLY, and motivation, worked arithmetic, scope
%  notes and "what this rules out" go after it.  NOTHING WAS DELETED: every
%  sentence removed from a statement is below, verbatim except for the connectives
%  and the splice-point punctuation needed to make it stand outside the
%  enumeration it used to sit in.
%
%  THEY ARE AT THE FOOT OF THE FILE ON PURPOSE.  Nothing printed follows them
%  here, so they renumber nothing in this file.  They are the first numbered
%  environments in this file, which is why the placement is stated rather than
%  assumed.
%
%  ASSUMPTIONS ARE STILL CITED BY TAG.  \begin{assumption} is \newtheorem*, so a
%  bare \ref{ass:...} prints nothing; the remarks below name C5', A2c, C4' and C3
%  in prose and never \ref them.
%
%  MIRROR NOTE.  This file is the undivided parent; the split mirror carries the
%  same remarks in assumptions_supp.tex.  Where the two sides of a moved sentence
%  already differed on disk, EACH SIDE KEEPS ITS OWN WORDING -- the divergences
%  are pre-existing and converging them is a content decision, not a formatting
%  one.  For C5'(ii) the divergence is "the one in the corpus is" here against
%  "the upper bound is" there.
% ---------------------------------------------------------------------------

\begin{remark}[How the smoothing-parameter block C5$'$ is read]\label{rem:app5read}
\emph{Two clauses, both of which stood inside the statement of C5$'$. Neither is a condition, and the
block restricts exactly the same bandwidth sequences without them.}
\begin{enumerate}
\item[(1)] \emph{What clause \textup{(i)}'s display buys.} It is chosen so that $h_u\to0$ and the smoothing
bias $h_u^{p_K}$ is $\asymp\hat\delta_n$. \emph{No lower bound on $h_u$ is required for this component.}
\item[(2)] \emph{Clause \textup{(ii)}'s displayed form, the stronger inherited form, and the window they
leave.} The displayed form is the one \cref{lem:aerr} consumes. \emph{The stronger inherited form
$\big(nh_x^{2}h_u^{3}\big)^{-1}\ln n\to0$ remains available, and is what the two arithmetic readings
recorded elsewhere in this Discussion are computed at, but it is imposed \textbf{only} where it is
compatible with an upper bound on $h_x$}, because clause (ii) is a \emph{lower} bound and the block
carries no upper bound at
all: the one in the corpus is \cref{prp:phirate}'s hypothesis (P-g), $\mathsf R h_x^{s-1/2}=O(\hat\delta_n)$.
At the tuning of clause (i) --- $h_u\asymp\hat\delta_n^{1/p_K}=n^{-\alpha/8}$, $\hat\delta_n=n^{-\alpha}$,
$p_K=8$ --- the stronger form reads $h_x\gg n^{-(1-3\alpha/8)/2}$ and (P-g) reads
$h_x=O\big(n^{-\alpha/(s-1/2)}\big)$, so the two leave a non-empty window \emph{if and only if}
\begin{equation}\label{eq:hxwindow}
s\;>\;\tfrac12\;+\;\frac{16\,\alpha}{8-3\alpha} ,
\end{equation}
whereas the form displayed in the statement reads $h_x\gg n^{-(1-\alpha/8)}$ and is compatible with (P-g)
as soon as
$s>\tfrac12+8\alpha/(8-\alpha)$ --- hence, since $\alpha<\tfrac12$ by A4(ii), at
\emph{every} $s>\tfrac{31}{30}$, uniformly in $\alpha$.
\end{enumerate}
\end{remark}

\begin{remark}[What A2c is, and what it is for]\label{rem:reprsourceread}
\emph{One clause, which stood inside the statement of A2c. It is a status note and not a condition, and
A2c restricts exactly the same representers without it.}
\begin{enumerate}
\item[(1)] \emph{A2c is NOT MAINTAINED.} It is a hypothesis \cref{lem:reprsmooth} imposes on a representer it
quantifies over, on the model of A1d, whose sole consumer is \cref{cor:a1cprim}. Its only role is to
\emph{imply} A2$'$(i) at $\ell^\Pi=\ell^\Pi_\star$ --- where the budget is written
$\Lambda^{(\mu'')}_\star$ --- and it is never asserted. \emph{A2c is stated in the same
$\bar w$-weighted $L^{1}(\Pi_t)$ geometry as A2$'$(i.a)--(i.b) themselves}, which is what
makes \cref{lem:reprsmooth}(iii) a one-line integration of a slicewise inequality rather than an
interchange of a supremum with an integral.
\end{enumerate}
\end{remark}

\begin{remark}[How \cref{def:hatT}'s two construction conventions are read]\label{rem:hatTread}
\emph{Four clauses, each of which stood inside the definition. None of it is construction:
\eqref{eq:foldbal}, the truncation, the definition of $\widehat{\widetilde W}_i$ and \eqref{eq:hatT} are
unchanged by its removal.}
\begin{enumerate}
\item[(1)] \emph{What the balanced split costs, and where it is consumed.}
\emph{This is a construction convention and it costs nothing --- the analyst chooses the split, and
$n_{\mathrm m}=\lceil n/2\rceil$ serves --- but it is \textbf{used}}: every variance bound in Part~H is
denominated in $n_{\mathrm m}$ and every
rate in $n$, so the two must be tied. It is consumed at \cref{lem:opnoise} Step~5,
\cref{prp:stopfeas}(i) and \cref{prp:phirate} Step~1 (N2.3), each of which converts an
$n_{\mathrm m}^{-1/2}$ or $n_{\mathrm m}/L_x$ into the corresponding statement in $n$. \emph{Without
\eqref{eq:foldbal} nothing in Part~H is false, but nothing is stated in $n$ either.}
\item[(2)] \emph{Two things follow from the truncation and both are used.} (a) The truncation can only decrease
$\|\hat\varrho-\varrho\|_{L^{2}}$, because $\varrho\in[\varepsilon_\varrho,1-\varepsilon_\varrho]$
almost surely by A4(i) and the metric projection onto an interval containing the target
is a pointwise contraction; so A4(iii) is preserved. (b) With
$\widehat{\widetilde W}_i:=W_i-\hat\varrho(\Delta Z_i)$ as \cref{def:hatT} defines it, we have
$\varepsilon_\varrho\le|\widehat{\widetilde W}_i|\le1$ \emph{surely}, for every $i$, because
$W_i\in\{0,1\}$.
\item[(3)] \emph{If only a smaller lower bound $\varepsilon\in(0,\varepsilon_\varrho]$ is known
to the analyst, truncating at $\varepsilon$ serves equally and every constant below carries
$\varepsilon$ in place of $\varepsilon_\varrho$.}
\item[(4)] \emph{\eqref{eq:hatT} is the object A7 names, and through round 9 no file defined it.}
\end{enumerate}
\end{remark}

\begin{remark}[How the kernel block C3 is read]\label{rem:app3read}
\emph{Two clauses, both of which stood inside the statement of C3. Neither is a condition, and the block
restricts exactly the same kernels without them.}
\begin{enumerate}
\item[(1)] \emph{What \eqref{eq:rescaled} is for.} $\Theta_{s_0}$ is the object every bound below C3 is
stated for, and the object \eqref{eq:intker} integrates.
\item[(2)] \emph{How clause \textup{(ii)} is read.} The boundary index is a distance in
\textbf{bandwidths}, and the defining property \eqref{eq:truncmom} is a \textbf{truncated} moment
condition.
\end{enumerate}
\end{remark}

% ============================================================================
%  ROUND-2 CHANGE LOG (against proofs_review.md, 56/100)
%
%  C-1  B1(ii) restricted to w = exp; the cos/sin claim deleted; the general-w
%       route recorded as rem:swintercept in results.tex with the correct
%       anchor (Stinchcombe & White 1998 Thm 3.1, intercept-indexed).
%  C-2  A4 note (c) rewritten: the identity holds at sigma(Delta Z) ONLY; the
%       "same range, same kernel" gloss WITHDRAWN, with the exact extra
%       condition (p_X = varrho) named and rejected.
%  M-1  A2, A2' and A3's referent moved from nu to nu/pi, with the Picard
%       orthogonality test recorded as the check that settles it.
%  M-3  A1b(iv) added: nonsingularity of E[Delta Z Delta Z'] and the second
%       moments.  Identification of theta_dagger no longer routed through A3.
%  M-4  Discussion: A2's necessity described as conjectural in this nonlinear
%       problem, per the contract; the "no estimator is root-n" claim removed.
%  M-5  Domain of A_cond fixed to V (see def:spaces); H1's note added.
%  M-6  B1(iv) weakened to full topological support (Pi_u) and non-mutual
%       singularity (Pi_t) -- exactly what Step 5 consumes.
%  M-7  C2'(ii) restated on u |-> M(u,t) uniformly in t, not on f_{Delta U};
%       C2'(i) reduced to p_K - 1 derivatives, which is what it is used for.
%  M-9  C6 retagged A1c and added to the stable-tag list.
%  M-10 A4(i) removed from prp:wellposed / prp:frechet / prp:compact; A0's
%       non-use in this dispatch stated explicitly for the assumption-auditor.
%  m-4  assapp3's boundary parameter renamed s_0; varsigma and s_0 declared
%       reserved local bound variables in the glossary.
%  m-5  Tag-vs-number bridge stated in the header.
%  m-6  B1(v) now carries the explicit dominating function for differentiating
%       L under the t-integral.
%  m-7  C4' now assumes positivity of f_{Delta U}, which kappa_n > 0 requires.
%  F-1  comment corrected: |Wtilde| <= 1, A4(i) idle; cross-reference F-3 -> F-1.
% ============================================================================
%
% ============================================================================
%  DISPATCH-2 CHANGE LOG.  FOUR CHANGES; ONE IS A STRENGTHENING, ONE IS A
%  DELETION, ONE IS AN ADDITION, ONE IS A TYPE CORRECTION.  TAGS ARE UNCHANGED.
%
%  1. B1(iv.a) STRENGTHENED.  Pi_u must now be absolutely continuous with a
%     bounded, strictly positive density varpi_u.  This IMPLIES the round-2 full
%     topological support clause, so lem:bierens Step 5 -- (iv.a)'s only
%     consumer before this dispatch -- is unaffected and no existing proof
%     changes.  It is a condition on a CHOSEN reference measure, not on P.  It
%     is what closes escalation N-10 by construction: see lem:adjoint(ii)-(iii)
%     and convention (P7) in proofs.tex.
%
%  2. A2'(0) DELETED.  It is now a theorem, lem:adjoint(ii).  Clauses (i) and
%     (ii) KEEP THEIR NUMBERS, so every downstream citation is unaffected.  A2'
%     now names the relation ell_star = W ell^Pi_star in its own statement
%     rather than in a comment.
%
%  3. A1d ADDED, NOT MAINTAINED.  Instrument completeness given the residual; a
%     primitive sufficient condition for A1c, cited by cor:a1cprim alone.  The
%     analogue of Feve & Florens (2014, Prop. 2.2)'s "strongly identified by W".
%     It is added rather than left inline so that cor:a1cprim can cite by tag
%     and so that the assumption-auditor sees a single declared consumer.
%
%  4. TYPE CORRECTION (theory_memo.md finding F-15), affecting no assumption
%     statement.  L, ell, ell_star and cor:udep's c are R^p-VALUED, one
%     coordinate per coordinate of theta.  A2's own wording -- "each coordinate
%     of nu/pi" -- already implies it; the glossary typed them as scalars, which
%     made psi and A3's J dimensionally ill-formed.  Found by
%     scripts/symbolic/specs/jacobian-dims-numeric.
%
%     WHICH TAGS ARE COORDINATEWISE, CORRECTED IN ROUND 7.  A2'(i), A2'(iii) and
%     N1 are read COORDINATEWISE -- each is p copies of a scalar statement.
%     A2'(ii) is written with the Euclidean norm, which is the coordinatewise
%     statement summed.  A3 IS NOT COORDINATEWISE AND MUST NOT BE READ THAT WAY:
%     J = E[Delta Z Delta Z'] + E[Wtilde ell_star Delta Z'] is p x p precisely
%     because ell_star is a genuine p x 1 VECTOR, so that
%     E[Wtilde ell_star Delta Z'] is (p x 1)(1 x p).  A coordinatewise reading
%     returns p scalars and J is ill-formed.  The round-2 change log listed A3
%     with the coordinatewise tags and the glossary Sec.8.4 item 3 said the
%     opposite; the glossary was right.


%
%  ALSO RECORDED, in the A2 comment: A2 acquires a primitive reading through
%  lem:adjoint(iii), and MAY BE STRICTLY STRONGER than what the estimator needs
%  (F-20, recorded not claimed); and, in rem:Tcond, the first half of B1(v) is
%  IMPLIED by B1(iv.c) AT gamma = Pi_t (F-19).
%
%  ROUND-7 SCOPE CORRECTION TO F-19, AND IT RUNS THE OTHER WAY FROM THE ROUND-6
%  SUMMARY.  It is NOT the case that "the second half of B1(v) is now the only
%  live half".  F-19 retires the first half ONLY AT gamma = Pi_t, which is the
%  POPULATION-side identification of the correction measure made by
%  lem:adjoint(ii).  lem:density(ii) quantifies over ALL finite signed gamma
%  satisfying B1(v) and its proof CONSTRUCTS AN ATOMIC gamma != Pi_t, verifying
%  the first half against that gamma directly -- so the FIRST half is live on the
%  estimation side.  The SECOND half is live in rem:Tcond only, since cor:elldiff
%  declines it twice and says why (it dominates partial_u lambda, one
%  u-derivative short of what eq:Wderiv consumes).  BOTH HALVES ARE RETAINED, and
%  each for its own reason.
% ============================================================================
%
% ============================================================================
%  DISPATCH-3 ROUND-7 CHANGE LOG (fix-only; against regularity-audit
%  dispatch2-partF, notation-table.md, usage-matrix.md).  TAGS ARE STABLE:
%  N1' and A5(vi) are ADDITIONS, A2'(iii) is an APPENDED CLAUSE, and no existing
%  tag or clause number moves.
%
%  A0     "compact" -> "COMPACT INTERVAL" (F-9).  Free; four results consume
%         connectedness.  The lem:adjfubini appeal to A0 is STRUCK as
%         over-attribution: P is sigma-finite by def:model.
%  A1b(i) reduced to its content, varphi_dagger in V; the generic-varphi half is
%         def:spaces and not an assumption.
%  A2'    "admits a VERSION" added (m-3); (ii) written with the Euclidean norm;
%         NEW CLAUSE (iii), the anchor of L_star, which closes CRITICAL F-2 --
%         L_star was consumed by eq:Psitau, lem:secondord, prp:orthog(ii)-(iii),
%         A5(iv) and the glossary and was DEFINED NOWHERE.  def:weight in
%         results.tex supplies the definition; (iii) prices its anchor and
%         nothing else, and costs nothing at either anchor of interest.
%  A4(d)  variance referent corrected from E[Wtilde^2 ell_star^2] to
%         E[Wtilde^2 ||L_star||^2].
%  N1     restated as a PROPERTY of a weight, not a construction (auditor
%         Finding 1: three incompatible statuses across five sites).  Struck from
%         lem:secondord, lem:solset and lem:balance, where it was unused.
%  N1'    NEW TAG for the direction condition E[D | Delta Z] = 0, which round 6
%         cited as "N1 on the direction" -- a third object under one tag.
%  A5(vi) NEW CLAUSE tagging ||h||_{H^s} = O_P(1), the last untagged substantive
%         condition in Part F; single consumer lem:balance(iii).
%  A7     tau -> tau_DP, to free tau for Part F's perturbation parameter.
%  A8     window priced ROUTE-CONDITIONALLY, per N-18; budget table gains a
%         "Route" column.
%  C1     five citations narrowed to C1(ii); C1(i)'s forward reference to Step 6
%         recorded.  THAT FORWARD REFERENCE WAS FALSE AND IS STRUCK IN DISPATCH
%         18: Step 6 (prp:phirate) cites C1 at no line and its list is narrowed
%         to C1(ii); C1(i)'s one consumer is lem:aerr Step 2.  See the corrected
%         note at C1 and the Part VI clause-status list in the Discussion.
%  C3     support interval C_0 written out as [-c_C, c_C], freeing C_0 for
%         cor:udep(c) (notation collision T2-7).
%  A1c    head comment DESTALED: prp:kerchar proves the passage the head called
%         unavailable (m-11).
%  A3     removed from the coordinatewise list; it is a genuine vector statement.
% ============================================================================

% ============================================================================
%  ROUND-8 CHANGE LOG (fix-only; against notation-table_round8.md 0/100 and
%  usage-matrix_round8.md 7/100).  TAGS REMAIN STABLE: A2b and def:hatA are
%  ADDITIONS; no existing tag or clause number moves.
%
%  A1b(i)  MODE SENTENCE ADDED, and it is the fix for auditor Finding 1 (-30).
%          A1b(i) is consumed for WELL-POSEDNESS, never for a bound; a result
%          about a generic h in V consumes def:spaces.  Four idle citations
%          struck or redirected: lem:nusigma(ii) (proofs.tex + the statement
%          line), lem:levels (hypothesis line -- it appeared NOWHERE in the
%          proof), prp:nophifree (hypothesis line + roadmap sentence), and
%          cor:a1cprim (which inherited it).  The four genuine consumers are
%          NAMED in the clause.
%  A2b     NEW TAG, NOT MAINTAINED: "admissible representer", the PROPERTY half
%          of A2'(i).  Auditor Finding 4, which GATED STEP 4.  A2'(i) gains an
%          explicit status sentence, on the N1 model.  cor:elldiff is restated at
%          a generic admissible representer with ell^Pi_star as the
%          SPECIALISATION, so that L-21 knows which object it discharges.
%  A2'     CITED BY CLAUSE at the four Part F results (Finding 8): A2, A2'(i),
%          A2'(iii).  A2'(ii) is consumed by none of them and is a Step-8 clause;
%          this matches the C1 -> C1(ii) and B1(iv) -> B1(iv.a) narrowings the
%          block performed in round 7 for exactly this reason.
%  A5(iv)  RESTATED (Finding 2, -15).  The clause now prices
%          ||ell-hat - ell_star||_{L^2(P)} ONLY.  The first equality of the
%          round-7 display, ||partial_u(L-hat - L_star)|| = ||ell-hat -
%          ell_star||, IS AN ASSERTION AND NOT NOTATION: it presupposes
%          L-hat = L^{b-hat}_{ell-hat} per def:weight(b), which is precisely the
%          obligation rem:dirmember records as OWED.  A forward obligation had
%          migrated into the block unstated.  It is now stated as a CONDITIONAL
%          consequence, with the two consumers separated by which form each uses.
%  A5(iv)  SCOPE NOTE NARROWED (Finding 7).  The round-7 claim "no bound in the
%   note   block is denominated in a norm this assumption does not price" is
%          block-wide and FALSE as printed -- proofs.tex refutes it in italics
%          two hundred lines later.  Scoped to the D-direction of clause (iv).
%  A4(iii) is now the tag cited for the r-direction rate at lem:secondord(iv)'s
%          instantiation (Finding 5), replacing a mis-citation to A5(iv); and
%          the bound itself is re-denominated in ||r||_{L^2(P)}, the norm A4(iii)
%          prices.  NOTHING prices ||varrho-hat - varrho||_infty, and nothing
%          now needs it.
%  A6/C3   def:hatA ADDED (warden V-8).  The three symbols A-hat_{phi-tilde},
%          A-hat_cond, A-hat*_cond appeared in a MAINTAINED assumption and in C3
%          with no defining display in any file.  The hat convention alone is
%          insufficient because A_phi is a Frechet derivative, so "sample
%          analogue" is ambiguous between "the derivative of T-hat" and "a
%          separately smoothed operator" -- and C3, by attaching a bandwidth,
%          implies the second.  def:hatA FIXES THE SECOND READING and records the
%          construction as a Step-6 obligation.
%  B1(ii)  comment: Phi(+-1) = +-a written out as +-Phi(1) (warden V-6); a is the
%          degree of ill-posedness and Sec.1.6 of the glossary built the sans-serif
%          constant slot expressly to keep constants off that letter.
% ============================================================================

% ============================================================================
%  DISPATCH-14 CHANGE LOG (fix-only, against four narrow audits arriving
%  together).  TWO CLAUSES ARE RETIRED IN PLACE AND NO TAG OR CLAUSE NUMBER
%  MOVES.  No new maintained assumption.  No conclusion of any result changes.
%
%  A2'(ii)  RETIRED IN PLACE, on the ass:rank precedent -- retired because
%           PROVED, not because false and not because too strong.  Under A2'(i)
%           IN FULL with B1(i), (ii), (iv.a), cor:elldiff gives
%           ||ell_star||_inf < infinity and (P1) gives |Wtilde| <= 1, so
%           E[Wtilde^2 ||ell_star||^2] <= ||ell_star||_inf^2 < infinity.  The
%           clause had NO CONSUMER -- results.tex says so three times and this
%           file's own ROUND-8 CHANGE LOG says so -- and A4's comment (d)
%           justified retaining it on a ground that contradicts that log.  The
%           forward Step-8 need under tier T3, where A2'(i) is relaxed and the
%           clause becomes a genuine Khan-Tamer irregularity boundary, is a
%           FORWARD OBLIGATION and is recorded as one at the retirement note.
%           CLAUSE NUMBER KEPT as a struck placeholder, so a stale citation lands
%           somewhere.  Two DESCRIPTIVE mentions amended: rem:adjointreconcile(i)
%           in results.tex and its rendering in paper/sections/theorems/.
%
%  B1(vi)   RETIRED IN PLACE, and for the THIRD of the three retirement reasons:
%           neither "it is a theorem" (A3) nor "it is false" (A6) but "it is
%           about a different estimator".  T_n occurred NOWHERE outside its own
%           sentence; def:hatT integrates against Pi_t over ALL of T.  Its
%           heading claimed identification of varphi_dagger, which in this block
%           is A1b(iii), assumed outright.  The counterexample-hunter's
%           finite-grid witness is recorded WITH the retirement, so a later
%           dispatch reinstating the growing-grid design finds both together.
%
%  A1b(i)   CONSUMER LIST REPAIRED.  lem:biasexp was missing from a list the
%           clause declares BINDING ON EVERY CITATION, while lem:biasexp lists
%           and cites A1b(i) CORRECTLY -- it is what puts h_k in V and makes the
%           balancing pairing the object lem:nusigma(iii) computes.  A reader
%           following the list would have struck a correct citation: the
%           deletion-inviting direction.  THE LIST WAS WRONG AND THE CITATION WAS
%           RIGHT.  Found independently by the assumption-auditor and the
%           regularity-checker, the latter as "a citation gap and not a
%           mathematical one".
%           Three further sites (lem:prelimtheta Step 4, lem:epterm Step 3,
%           lem:sandwich Steps 3 and 6') were NOT added, because lem:nusigma(i)
%           is weakened this round from "for every h in V" to "for every h in
%           L^2(X,pi)" -- all its proof establishes -- after which none of the
%           three needs the tag at all.
%
%  A4(d)    TWO FALSE STATEMENTS STRUCK, one legitimate ground kept.  (1) "A2'(ii)
%           is the clause the second-order and bias terms consume" -- false, and
%           contradicted by this file's own round-8 log.  (2) The bare "(i) +
%           (iii)", read inside a comment attached to A4, attributes finiteness
%           of E[Wtilde^2 ||L_star||^2] to OVERLAP; A4(i) plays NO role in that
%           bound, whose tags are A2'(i), A2'(iii), A1b(iv), B1(i), (ii), (iv.a)
%           and (P1).  What survives, and is now the only ground stated, is that
%           overlap governs the SIZE of the variance: with the conditional laws
%           given W held fixed and varrho == eps, ranges are invariant (so A2 is
%           untouched) while E[Wtilde^2 ell_star^2] asymp eps^{-1}.
%
%  A4(ii)   d_Z >= 1 STATED, and it was assumed silently at four sites.  At
%           d_Z = 0 the exponent alpha = r/(2r+d_Z) is exactly 1/2, and (P-d)'s
%           sieve-entropy condition L_x ln n = o(n delta-hat_n^2) then reads
%           L_x = o(1) -- no integer L_x >= 1 satisfies it, eq:lxwindow is empty
%           at every s, and (T-5) is unsatisfiable.  Three further sites print
%           "alpha < 1/2" as a REASON and are false as printed although their
%           conclusions survive at the boundary.  d_Z = 0 is an EASIER problem
%           (a constant propensity) excluded for bookkeeping, not a pathology,
%           and the note says so.
%
%  H1       THE IDENTITY IS SUBSCRIPTED AND SCOPED: mu_{H1} := s/(2a) is the
%           source index of A2, at the POPULATION operator on ALL of V at
%           theta_dagger.  BARE mu is (P-a.3)'s PRIMITIVE index at the
%           finite-rank sieve compression of the REALISED operator at theta_o.
%           They are equal only under the interpolation proviso rem:pasrc(4)
%           DECLINES.  The two budget-table lines that compute mu_{H1} carry a
%           dagger note; the sentence deducing mu > 5/6 from s > a + 1/2 carries
%           the proviso; the flagship (mu,s,a,alpha) = (2,3,3/4,0.32) is read as
%           FOUR numbers of which two happen to satisfy s/(2a) = mu.
%           rem:pasrc(5) carries the binding reading; glossary Sec.8.8 documents
%           the hazard.
% ============================================================================

% ============================================================================
%  ROUND-9 CHANGE LOG (memo Sec.10 STEP 4).  ONE TAG IS ADDED AND NO EXISTING
%  TAG, CLAUSE NUMBER OR STATEMENT MOVES.
%
%  A2c     NEW TAG, NOT MAINTAINED: the source condition of index mu'' > 3/2 on a
%          representer, in the u-Sobolev scale on R with respect to LEBESGUE
%          measure, integrated against bar-w d Pi_t.  Modelled on A1d: a primitive
%          sufficient condition with exactly one consumer (lem:reprsmooth = L-21),
%          never asserted.  It DISCHARGES obligation (O1) of rem:obligations by
%          supplying a route to A2'(i.a)-(i.b) IN SUP-INSIDE FORM ON ell^Pi_star.
%          A2'(i) REMAINS MAINTAINED; what it loses is the status of an
%          unverifiable primitive.
%  A2'     comment blocks DESTALED at two places, both saying L-21 was unproved.
%          The (alpha)-versus-(beta) note and item 2 of "FOUR THINGS" now record
%          that L-21 is proved, that the threshold is mu'' > 3/2, and that the
%          sup-inside target is MET rather than approximated -- because A2c is
%          stated in the same bar-w L^1(Pi_t) geometry as the clauses it implies,
%          so no supremum is ever moved across the t-integral.
%  A2b     comment block CORRECTED at three lines that said "R1" where they meant
%          A2b -- the tag this block itself records as REJECTED because rem:R1
%          occupies it.  Found by sweeping the retired string across the file
%          rather than by any finding.  Glossary Sec.1.5b's rule applied to a
%          comment block: the round-8 rename reached every statement and citation
%          and not the prose that explains the rename.
%  BUDGET  the "Representer" line of the smoothness table now points at A2c and
%   TABLE  lem:reprsmooth instead of A2'(i), pins mu'' > 3/2 rather than leaving
%          the index free, and says in the following paragraph that the line
%          records a SUFFICIENT CONDITION and not an assumption.
%
%  WHAT IS DELIBERATELY NOT DONE.
%   - NO CLAUSE OF A2' IS WEAKENED OR REMOVED.  An assumption with a proved
%     sufficient condition is still an assumption; deleting (i) and keeping A2c
%     would make A2c maintained, which is exactly the A1c/A1d relation the block
%     already declines to collapse.
%   - NO NEW CITATION.  The three classical facts consumed at Step 4 -- the
%     one-dimensional Sobolev embedding, the numerical-radius identity for a
%     self-adjoint operator, and the essential-spectrum invariance named as Weyl's
%     theorem at lem:solset(iii) -- are each PROVED INLINE in proofs.tex, at the
%     level of detail the step consumes.  The reason is recorded rather than
%     silent: no spectral-theory or Sobolev-space entry resolves in this project's
%     bibliography (conway2000 was DEPINNED in round 6), and adding an
%     unresolvable \cite breaks the build.  The intended pins are named in a
%     comment at each site so that the bibliography pass has a target.
%     **THIS CLAUSE IS A RECORD OF THIS DISPATCH'S SCOPE AND IS NOW STALE AS A
%     DESCRIPTION OF proofs.tex.  AMENDED 2026-08-02 AND AGAIN 2026-08-03; the
%     amendment is here rather than by deletion, because the clause is what
%     licensed the choice at the time.**  What is now true, fact by fact:
%       * ~~no spectral-theory or Sobolev-space entry resolves~~ -- brezis2011,
%         halmos1982, conway1990 and kato1995 all resolve, every field and locator
%         read off an object on disk.
%       * ~~The intended pins are named in a comment~~ -- they are no longer
%         "intended".  Three are attached: conway1990 Thm II.1.5 at
%         lem:formstruct(i), which is why a FOURTH fact is no longer proved inline
%         at all; three integer-order sub-facts of lem:reprsmooth Step 1' to
%         brezis2011; and kato1995 Thm IV.5.35 at BOTH of lem:solset(iii)'s
%         previously bare "Weyl's theorem" invocations.
%       * The three facts listed above ARE still proved inline, and now for stated
%         reasons rather than for want of a book: the Sobolev embedding because
%         Brezis has no fractional space and no usable constant (M-3); the
%         numerical-radius identity because Halmos is complex-only and the
%         statement is false over R for a normal operator (M-5); and the
%         essential-spectrum CRITERION because Kato is complex-only by a global
%         convention and his nearest statement prices approximate nullity with a
%         normalized rather than an orthonormal sequence (M-4d).  The INVARIANCE,
%         by contrast, is pinned.
%     The register at the head of proofs.tex carries all of it, item by item.
% ============================================================================
